\documentclass[UTF8,a4paper,11pt]{ctexart}

\usepackage[margin=1.8cm]{geometry}

\usepackage{amsmath,amssymb,mathtools,bm}

\usepackage{enumitem,longtable,booktabs,xcolor,hyperref,tikz,pgfplots}

\hypersetup{hidelinks}

\usetikzlibrary{arrows.meta,calc,intersections,patterns,positioning,quotes,angles}

\pgfplotsset{compat=1.18}

\setlist{leftmargin=2em}

\title{圆锥曲线射影几何延伸阅读（非高中常规教案）}

\author{教师参考稿}

\date{\today}

\begin{document}

\maketitle

\begin{quote}

本稿保留为教师延伸阅读，不作为常规高中圆锥曲线专题教案直接使用。射影几何可用于教师理解结构与预判方向；学生课堂与高考书写必须回到课程标准范围内的定义、坐标法、联立消元、韦达定理、点差法、向量和必要的平面几何。

\end{quote}

\tableofcontents

\clearpage

\section{第1课：从透视投影到圆锥曲线}

\noindent 建议课时：1课时。

\subsection*{教学目标}

\begin{itemize}

\item 理解射影平面中“补上无穷远点后，平行与相交可统一表述”的基本思想，知道它只用于观察结构，不直接作为高考书面解答工具。

\item 能从统一视角认识圆、椭圆、双曲线、抛物线都属于圆锥曲线，知道“透视投影一般不保长度、角度，但有些结构关系可保留”。

\item 会把射影视角下得到的结构预判，翻译回高中课标允许的方法：坐标法、定义法、韦达定理、点差法、向量或平面几何。

\item 能用圆锥曲线第一定义、第二定义和轨迹法，处理给定高考题中的基础结论，体会“统一视角 $\to$ 可计算表达”的转化路径。

\end{itemize}

\subsection*{重点与难点}

\begin{itemize}

\item “无穷远点”只作观念统一：欧氏平面中的平行线在射影平面中看作交于无穷远点。

\item 圆锥曲线的统一来源：平面截圆锥可得到椭圆、抛物线、双曲线；圆可视为二次曲线中的特殊情形。

\item 射影视角的教学边界：可以帮助猜结构、找不变量、理解定义来源，但正式解答必须回到高中方法。

\item 从定义出发的高中落地：如抛物线用 $|MF|=d(M,l)$，椭圆用 $|MF_1|+|MF_2|=2a$，双曲线用 $\bigl||MF_1|-|MF_2|\bigr|=2a$，隐圆用轨迹方程配方识别。

\end{itemize}

\subsection*{高观点}

本节只把射影几何作为“统一视角”。核心信息有三层：其一，透视投影告诉我们，图形外观会变，但某些结构关系值得追踪；其二，把平行看成在无穷远处相交后，很多分类会更统一，例如“相交弦”“平行弦方向”“抛物线的轴向”都可放在同一框架下理解；其三，圆、椭圆、双曲线、抛物线都可纳入圆锥曲线家族。课堂上只做概念溯源与结构预判，不把交比、极点极线作为本节正式求解工具。

\subsection*{高中落地}

统一视角最终要翻译为高中可写的方法。若结构指向抛物线，就回到第二定义 $|MF|=d(M,l)$；若指向椭圆或双曲线，就回到第一定义；若轨迹经变形出现 $x^2+y^2+Dx+Ey+F=0$，就用配方法识别圆；若题目给出坐标与向量关系，就直接设点列式、化简、配方、求最值。课堂强调：结论的发现可以借助“统一视角”，结论的证明必须用高中工具完成。

\subsection*{教学流程}

\paragraph{5分钟：导入}

教师展示生活中的透视现象：铁轨、走廊、道路边线“看起来相交”。提出问题：欧氏几何里平行线不相交，为什么视觉上却会“相交”？顺势引入“把平行线看成在无穷远处相交”的统一想法，并说明本节目的不是学竞赛式射影证明，而是为一轮复习建立圆锥曲线的整体框架。

\noindent\textbf{教师追问：}两条平行铁轨在照片中为什么看起来会交到一点？\\

\textbf{预期回答：}因为透视投影会改变图形外观，远处看起来汇聚。\par

\noindent\textbf{教师追问：}如果把这个现象数学化，你觉得“平行”能不能被当成“特殊的相交”？\\

\textbf{预期回答：}可以把平行看成在很远的地方相交，这样表述更统一。\par

\paragraph{8分钟：概念建立}

教师用简图说明射影平面的最少语言：补入无穷远点、无穷远直线；强调这不是高考考点表述，而是帮助理解。接着说明圆锥曲线的统一来源：平面截圆锥可得到椭圆、抛物线、双曲线。板书统一关系：$\text{圆}\subset\text{二次曲线}$，而椭圆、抛物线、双曲线同属圆锥曲线。

\noindent\textbf{教师追问：}为什么补上无穷远点后，‘平行’与‘相交’的处理会更统一？\\

\textbf{预期回答：}因为原来不相交的平行线也被视为有交点，只是交点在无穷远处。\par

\noindent\textbf{教师追问：}椭圆、抛物线、双曲线之间，今天最重要的共同点是什么？\\

\textbf{预期回答：}它们都属于圆锥曲线，可以放在同一类中研究。\par

\paragraph{10分钟：从统一视角回到高中定义}

教师把“统一家族”翻译成高考常用定义。椭圆：$|MF_1|+|MF_2|=2a$；双曲线：$\bigl||MF_1|-|MF_2|\bigr|=2a$；抛物线：$|MF|=d(M,l)$。指出：射影视角告诉我们它们是一家人，但真正做题时，最稳定的方法仍是定义、坐标与方程。结合抛物线的准线思想，说明“无穷远方向”在直观上有助于理解抛物线只开向一个方向。

\noindent\textbf{教师追问：}如果题目出现‘到定点与到定直线距离相等’，你首先想到哪一类曲线？\\

\textbf{预期回答：}想到抛物线，用第二定义。\par

\noindent\textbf{教师追问：}如果题目出现‘到两个定点距离和为定值’或‘差的绝对值为定值’，分别想到什么？\\

\textbf{预期回答：}前者是椭圆，后者是双曲线，用第一定义。\par

\paragraph{8分钟：例题观察一}

引入题目 ID：`2023-北京卷::6`。不改写原题，只做结构观察。教师引导学生先不算，先判断：题中给出抛物线 $C:y^2=8x$、焦点 $F$、点 $M$ 在曲线上、点到直线 $x=-3$ 的距离。这里最关键的不是二次项，而是“点到准线距离”。因为 $y^2=8x$ 可写成 $y^2=2px$，得 $p=4$，准线为 $x=-2$。随后强调正式求解应回到第二定义：$|MF|=d(M,\text{准线})$，再结合点到 $x=-3$ 的距离进行比较。

\noindent\textbf{教师追问：}看到 `2023-北京卷::6`，第一反应是用哪条定义？\\

\textbf{预期回答：}用抛物线第二定义。\par

\noindent\textbf{教师追问：}为什么这道题不必急着设 $M(x,y)$ 大量运算？\\

\textbf{预期回答：}因为题目已经给了点到直线的距离信息，和准线距离天然相关。\par

\paragraph{6分钟：例题观察二}

引入题目 ID：`2022-全国乙卷-理科::5`。教师提示：焦点、点在抛物线上、又出现 $|AF|=|BF|$，仍应优先想到抛物线定义，把 $|AF|$ 转成点 $A$ 到准线的距离，再与点 $B(3,0)$ 的焦点距离联系。这里借机强调“统一视角”的价值不在于替代计算，而在于更快识别入口。

\noindent\textbf{教师追问：}这道题里，哪一个等量关系最值得先翻译？\\

\textbf{预期回答：}先翻译 $|AF|=|BF|$。\par

\noindent\textbf{教师追问：}如果把 $|AF|$ 换成到准线的距离，题目会变得更像哪类基础题？\\

\textbf{预期回答：}会变成点到直线距离与已知线段长度比较的题。\par

\paragraph{8分钟：例题观察三}

引入题目 ID：`2017-全国II卷-理科::20`。教师带学生只看第（1）问的结构：点 $M$ 在椭圆 $C:\dfrac{x^2}{2}+y^2=1$ 上，过 $M$ 作 $x$ 轴垂线得 $N$，再由 $\overrightarrow{NP}=\sqrt{2}\overrightarrow{NM}$ 定义点 $P$。这里的关键是向量伸缩导致坐标线性变化，轨迹可能从椭圆“变形”为更熟悉的二次曲线。计算后若出现 $x^2+y^2=2$，就应立即识别为圆。这正体现“隐圆”的高中方法，而不是射影方法。

\noindent\textbf{教师追问：}这道题从哪一步能看出适合用坐标法？\\

\textbf{预期回答：}因为点在椭圆上且有垂足、向量倍数关系，坐标表达最直接。\par

\noindent\textbf{教师追问：}当轨迹方程化成 $x^2+y^2=2$ 时，你应立刻想到什么？\\

\textbf{预期回答：}想到轨迹是圆，圆心在原点，半径为 $\sqrt{2}$。\par

\paragraph{7分钟：课堂小结与方法归纳}

教师归纳本节方法链：先用射影视角统一认知，再用高中定义落地求解。给出三条识别线索：若见“焦点+准线/到直线距离”，先想抛物线第二定义；若见“两焦点距离和差”，先想第一定义；若见轨迹方程中 $x^2,y^2$ 系数相同，优先检查是否为隐圆。再明确边界：本节不要求学生掌握交比和极点极线计算，也不能在高考书面解答中直接写射影结论代替证明。

\noindent\textbf{教师追问：}今天这节课最重要的不是哪个公式，而是哪种思考顺序？\\

\textbf{预期回答：}先统一识别结构，再回到定义和坐标求解。\par

\noindent\textbf{教师追问：}正式高考解答时，能不能写‘由射影几何知结论成立’？为什么？\\

\textbf{预期回答：}不能，因为高考解答必须用课标允许的方法写出可评分过程。\par

\paragraph{3分钟：当堂检测}

学生口答分类，不展开完整计算。教师依次报出题目 ID：`2023-北京卷::6`、`2022-全国乙卷-理科::5`、`2017-全国II卷-理科::20`，要求学生说出首选方法分别是什么：第二定义、第二定义、坐标轨迹法识别隐圆。

\noindent\textbf{教师追问：}对 `2023-北京卷::6` 的首选工具是什么？\\

\textbf{预期回答：}抛物线第二定义。\par

\noindent\textbf{教师追问：}对 `2022-全国乙卷-理科::5` 的首选工具是什么？\\

\textbf{预期回答：}抛物线第二定义。\par

\noindent\textbf{教师追问：}对 `2017-全国II卷-理科::20` 第（1）问的首选工具是什么？\\

\textbf{预期回答：}坐标法列轨迹，再识别圆。\par

\subsection*{课堂真题}

\paragraph{2023 年高考数学 北京卷第6题}

已知抛物线 $C: y^2 = 8x$ 的焦点为 $F$ ，点 $M$ 在 $C$ 上．若 $M$ 到直线 $x = -3$ 的距离为 $5$，则 $|MF| =$ （ ）

\textbf{答案：}D\par

\textbf{解答：}因为抛物线 $C: y^2 = 8x$ 的焦点 $F(2,0)$ ，准线方程为 $x = -2$ ，点 $M$ 在 $C$ 上，所以 $M$ 到准线 $x = -2$ 的距离为 $|MF|$，又 $M$ 到直线 $x = -3$ 的距离为 $5$，所以 $|MF| + 1 = 5$，故 $|MF| = 4$.

\paragraph{2022 年高考数学 全国乙卷-理科第5题}

设$F$为抛物线$C: y^2 = 4x$的焦点，点$A$在$C$上，点$B(3, 0)$，若$|AF| = |BF|$，则$|AB| =$（\qquad）

\textbf{答案：}B\par

\textbf{解答：}由题意得，$F(1,0)$，则$|AF| = |BF| = 2$，\\即点$A$到准线$x=-1$的距离为2，所以点$A$的横坐标为$-1+2=1$，\\不妨设点$A$在$x$轴上方，代入得，$A(1,2)$，\\所以$|AB| = \sqrt{(3-1)^2 + (0-2)^2} = 2\sqrt{2}$.

\paragraph{2017 年高考数学 全国II卷-理科第20题}

设 $O$ 为坐标原点，动点 $M$ 在椭圆 $C:\frac{x^2}{2}+y^2=1$ 上，过 $M$ 作 $x$ 轴的垂线，垂足为 $N$，点 $P$ 满足 $\overrightarrow{NP}=\sqrt{2}\overrightarrow{NM}$．
（1）求点 $P$ 的轨迹方程；
（2）设点 $Q$ 在直线 $x=-3$ 上，且 $\overrightarrow{OP}\cdot\overrightarrow{PQ}=1$．证明：过点 $P$ 且垂直于 $OQ$ 的直线 $l$ 过 $C$ 的左焦点 $F$．

\textbf{答案：}（1）$x^2+y^2=2$；（2）略\par

\textbf{解答：}（1）设 $M(x_0,y_0)$，由题意可得 $N(x_0,0)$，设 $P(x,y)$，由点 $P$ 满足 $\overrightarrow{NP}=\sqrt{2}\overrightarrow{NM}$．可得 $(x-x_0,y)=\sqrt{2}(0,y_0)$，可得 $x-x_0=0$，$y=\sqrt{2}y_0$，即有 $x_0=x$，$y_0=\frac{y}{\sqrt{2}}$，代入椭圆方程 $\frac{x^2}{2}+y^2=1$，可得 $\frac{x^2}{2}+\frac{y^2}{2}=1$，即有点 $P$ 的轨迹方程为圆 $x^2+y^2=2$；
（2）证明：设 $Q(-3,m)$，$P(\sqrt{2}\cos\alpha,\sin\alpha)$（$0\le\alpha<2\pi$），$\overrightarrow{OP}\cdot\overrightarrow{PQ}=1$，可得 $(\sqrt{2}\cos\alpha,\sin\alpha)\cdot(-3-\sqrt{2}\cos\alpha,m-\sin\alpha)=1$，即为 $-3\sqrt{2}\cos\alpha-2\cos^2\alpha+\sqrt{2}m\sin\alpha-2\sin^2\alpha=1$，当 $\alpha=0$ 时，上式不成立，则 $0<\alpha<2\pi$，解得 $m=\frac{3(1+\sqrt{2}\cos\alpha)}{\sqrt{2}\sin\alpha}$，即有 $Q(-3,\frac{3(1+\sqrt{2}\cos\alpha)}{\sqrt{2}\sin\alpha})$，椭圆 $\frac{x^2}{2}+y^2=1$ 的左焦点 $F(-1,0)$，由 $\overrightarrow{PF}\cdot\overrightarrow{OQ}=(-1-\sqrt{2}\cos\alpha,-\sin\alpha)\cdot(-3,\frac{3(1+\sqrt{2}\cos\alpha)}{\sqrt{2}\sin\alpha})=3+3\sqrt{2}\cos\alpha-3(1+\sqrt{2}\cos\alpha)=0$．可得过点 $P$ 且垂直于 $OQ$ 的直线 $l$ 过 $C$ 的左焦点 $F$．

\subsection*{课后作业}

\begin{enumerate}

\item 2023 年高考数学 北京卷第6题

\item 2022 年高考数学 全国乙卷-理科第5题

\item 2017 年高考数学 全国II卷-理科第20题

\item 2021 年高考数学 上海春季卷第19题

\end{enumerate}

\noindent\textbf{作业答案索引：}

\begin{itemize}

\item 2023 年高考数学 北京卷第6题：D

\item 2022 年高考数学 全国乙卷-理科第5题：B

\item 2017 年高考数学 全国II卷-理科第20题：（1）$x^2+y^2=2$；（2）略

\item 2021 年高考数学 上海春季卷第19题：$(1)\ \frac{x^2}{100}-\frac{y^2}{300}=1,\ P(\frac{15\sqrt{2}}{2},\frac{5\sqrt{6}}{2})$；$(2)\ |OQ|\approx 19$ 米，$Q$ 点位置北偏东 $66^\circ$

\end{itemize}

\subsection*{板书设计}

一、透视现象与统一视角
\[
\text{平行线在透视图中“汇聚”} \Rightarrow \text{补上无穷远点后更统一}
\]
\[
\text{射影视角：用于结构理解\text{，}不直接作高考书面证明}
\]

二、圆锥曲线统一家族
\[
\text{椭圆、抛物线、双曲线}\in\text{圆锥曲线},\qquad \text{圆}\subset\text{二次曲线}
\]

三、回到高中定义
\[
\text{椭圆： } |MF_1|+|MF_2|=2a
\]
\[
\text{双曲线： } \bigl||MF_1|-|MF_2|\bigr|=2a
\]
\[
\text{抛物线： } |MF|=d(M,l)
\]

四、题型入口
\[
\text{焦点+准线} \Rightarrow \text{第二定义}
\]
\[
\text{两焦点和/差} \Rightarrow \text{第一定义}
\]
\[
\text{轨迹化简得 } x^2+y^2+Dx+Ey+F=0 \Rightarrow \text{圆}
\]

五、课堂例题 ID
\[
\texttt{2023-北京卷::6}\quad \texttt{2022-全国乙卷-理科::5}\quad \texttt{2017-全国II卷-理科::20}
\]

\subsection*{易错点}

\begin{itemize}

\item 把“射影几何统一视角”误当成“高考可直接使用的证明工具”。纠正：它只能帮助观察与预判，正式解答必须回到高中方法。

\item 见到抛物线方程就立刻设点硬算，忽视题目已给出的“到直线距离”信息。纠正：优先检查是否可直接调用第二定义。

\item 把“圆也是特殊曲线”与“圆属于圆锥曲线”混为同一层面表述。课堂可说明圆常视为二次曲线中的特殊情形，本节重在统一认识，不深究严格分类细节。

\item 轨迹题化简后没有及时识别 $x^2$ 与 $y^2$ 系数相等这一信号。纠正：一旦出现此结构，先尝试配方判断是否为圆。

\item 只会背定义，不会选入口。纠正：先看题中给的是焦点、准线、距离和差，还是向量伸缩与轨迹。

\end{itemize}

\subsection*{教师备注}

1. 本节定位是“一轮复习中的总领课”，目标不是推导交比、极点极线，也不是做高难射影证明，而是建立学生的结构意识。
2. 课堂语言要反复强调边界：射影视角只用于说明“为什么这些曲线值得放在一起看”，不可替代高中课标方法。
3. 若学生追问“抛物线为什么也能与透视、无穷远方向联系”，教师可用直观语句说明：抛物线只有一个开口方向，可以把它理解为某种“朝向无穷远”的极限情形，但不在课堂上展开严格射影论证。
4. 处理 `2017-全国II卷-理科::20` 时，建议只完成第（1）问的方法观察与结果识别，把第（2）问留作课后提升，避免本节被大量计算拖慢。
5. 作业布置建议分层：基础层完成 `2023-北京卷::6`、`2022-全国乙卷-理科::5`；提高层加做 `2017-全国II卷-理科::20` 与 `2021-上海春季卷::19`。批改时重点看“方法入口是否正确”，不只看最终答案。
6. 若需一句话收束本课，可板书：
\[
\boxed{\text{先用统一视角看懂一家人\text{，}再用高中方法把题做出来}}
\]

\clearpage

\section{第2课：交比、调和与极点极线}

\noindent 建议课时：2课时。

\subsection*{教学目标}

\begin{itemize}

\item \(\text{理解并会表述交比 }(A,B;C,D)\text{ 的定义、调和点列 }(A,B;C,D)=-1\text{ 及其基本含义\text{。}}\)

\item \(\text{会用“外点两切线 }\Rightarrow\text{ 切点弦为极线”“割线 }AB\text{ 与两切线交点 }P\text{ 互为极点极线”识别图形结构\text{。}}\)

\item \(\text{能用 }\mathrm{La\ Hire}\text{ 互易性 }X\in \ell_P\iff P\in \ell_X\text{ 进行结构预判\text{，}但书面解答回到坐标法、定义法、韦达定理、点差法或平面几何\text{。}}\)

\item \(\text{能把射影视角翻译为高中可写结论：如椭圆极线 }\frac{x_0x}{a^2}+\frac{y_0y}{b^2}=1\text{、抛物线切点弦 }y_0y=p(x+x_0)\text{\text{。}}\)

\item \(\text{通过真题建立“定点、定直线、切点弦、斜率和/积不变量”之间的联系\text{。}}\)

\end{itemize}

\subsection*{重点与难点}

\begin{itemize}

\item \(\text{交比的定义与调和点列： }(A,B;C,D)=-1\text{。}\)

\item \(\text{极点极线的三种识别：切点弦定义、调和定义、解析方程定义\text{。}}\)

\item \(\mathrm{La\ Hire}\text{ 互易性： }X\in \ell_P\iff P\in \ell_X\text{。}\)

\item \(\text{圆锥曲线中“动点在定直线上 }\Longleftrightarrow\text{ 动极线过定点”的结构翻译\text{。}}\)

\item \(\text{正式解题回归高中方法：联立方程、韦达定理、斜率关系、切线方程、点差法\text{。}}\)

\end{itemize}

\subsection*{高观点}

\(\text{本课中的射影几何只承担三项功能：}\;\boxed{\text{统一视角}}\;\boxed{\text{结构预判}}\;\boxed{\text{概念溯源}}\text{。}\)\\
\(\text{统一视角：把“外点两切线、切点弦、割线、共点共线、定点定直线”放入同一语言：极点极线\text{。}}\)\\
\(\text{结构预判：若某点 }Q\text{ 在定直线 }l\text{ 上运动\text{，}则往往预判其极线 }\ell_Q\text{ 过 }l\text{ 的极点\text{。}}\)\\
\(\text{概念溯源：极线公式并非死记硬背\text{，}而可追溯到“割线上的调和共轭点轨迹”\text{。}}\)\\
\(\text{但任何正式高考解答均不直接写“由射影几何可得”\text{，}必须落实为高中课标允许的方法\text{。}}\)

\subsection*{高中落地}

\(\text{常用翻译规则如下：}\)\\
\(1.\;\text{椭圆 }\frac{x^2}{a^2}+\frac{y^2}{b^2}=1\text{ 中\text{，}点 }P(x_0,y_0)\text{ 的极线写成 }\frac{x_0x}{a^2}+\frac{y_0y}{b^2}=1\text{。}\)\\
\(2.\;\text{双曲线 }\frac{x^2}{a^2}-\frac{y^2}{b^2}=1\text{ 中\text{，}点 }P(x_0,y_0)\text{ 的极线写成 }\frac{x_0x}{a^2}-\frac{y_0y}{b^2}=1\text{。}\)\\
\(3.\;\text{抛物线 }x^2=2py\text{ 中\text{，}点 }P(x_0,y_0)\text{ 的切点弦/极线写成 }xx_0=p(y+y_0)\text{；}\text{若写成 }y^2=2px\text{\text{，}则为 }yy_0=p(x+x_0)\text{。}\)\\
\(4.\;\text{若题目出现“过曲线外一点作两条切线\text{，}切点为 }A,B\text{”\text{，}则 }AB\text{ 往往是切点弦\text{，}可直接转入方程法\text{。}}\)\\
\(5.\;\text{若题目出现“动点在定直线上、动弦过定点”\text{，}可先用射影视角预判\text{，}再用联立方程与韦达定理证明\text{。}}\)

\subsection*{教学流程}

\paragraph{6分钟：导入与目标}

\(\text{教师出示本课主线：}\quad \text{交比}\rightarrow\text{调和}\rightarrow\text{极点极线}\rightarrow\text{定点定直线题型\text{。}}\)\\
\(\text{强调边界：射影几何用于看结构\text{，}书面解答必须回到高中方法\text{。}}\)

\noindent\textbf{教师追问：}\(\text{看到“外点向圆锥曲线作两条切线\text{，}连接切点”\text{，}你会联想到什么对象？}\)\\

\textbf{预期回答：}\(\text{切点弦\text{，}或该外点的极线\text{。}}\)\par

\noindent\textbf{教师追问：}\(\text{如果一个结论总表现为“动点在定直线上”\text{，}它的对偶结论可能是什么？}\)\\

\textbf{预期回答：}\(\text{可能对应“动极线过定点”\text{。}}\)\par

\paragraph{12分钟：概念建构一：交比与调和}

\(\text{在一条直线上定义 }(A,B;C,D)=\dfrac{\overrightarrow{CA}/\overrightarrow{CB}}{\overrightarrow{DA}/\overrightarrow{DB}}\text{。}\)\\
\(\text{若 }(A,B;C,D)=-1\text{\text{，}则称 }C,D\text{ 关于 }A,B\text{ 调和共轭\text{。}}\)\\
\(\text{只要求学生会识别“内分点+外分点且比值相同”这一调和原型\text{，}不在课堂上做抽象证明\text{。}}\)

\noindent\textbf{教师追问：}\(\text{为什么调和点列里常常同时出现一个内分点和一个外分点？}\)\\

\textbf{预期回答：}\(\text{因为交比为 }-1\text{ 反映了比值大小相同、方向相反\text{。}}\)\par

\noindent\textbf{教师追问：}\(\text{若 }(A,B;C,D)=-1\text{\text{，}它最核心的信息是什么？}\)\\

\textbf{预期回答：}\(\text{说明 }C,D\text{ 关于 }A,B\text{ 构成调和共轭\text{，}可在投影下保持\text{。}}\)\par

\paragraph{10分钟：概念建构二：极点极线与 La Hire 互易性}

\(\text{从几何定义进入：外点 }P\text{ 向圆锥曲线作两切线\text{，}切点为 }T_1,T_2\text{\text{，}则 }T_1T_2\text{ 为 }P\text{ 的极线 }\ell_P\text{。}\)\\
\(\text{给出互易性： }X\in\ell_P\iff P\in\ell_X\text{。}\)\\
\(\text{说明：这是结构工具\text{，}不是考场直接作答语言\text{。}}\)

\noindent\textbf{教师追问：}\(\text{若 }AB\text{ 是圆锥曲线的一条割线\text{，}而 }A,B\text{ 处切线相交于 }P\text{\text{，}那么 }P\text{ 与 }AB\text{ 是什么关系？}\)\\

\textbf{预期回答：}\(\text{ }P\text{ 是 }AB\text{ 的极点\text{，}}AB\text{ 是 }P\text{ 的极线\text{。}}\)\par

\noindent\textbf{教师追问：}\(\text{La Hire 互易性在“证明共线/共点”中能做什么？}\)\\

\textbf{预期回答：}\(\text{可把“某点在某线上”转换成“另一点在对应极线上”\text{，}便于预判结构\text{。}}\)\par

\paragraph{12分钟：公式落地：高中可写形式}

\(\text{板书并整理三类常用方程：}\)\\
\(\text{椭圆 }\frac{x^2}{a^2}+\frac{y^2}{b^2}=1\Rightarrow \frac{x_0x}{a^2}+\frac{y_0y}{b^2}=1;\)\\
\(\text{双曲线 }\frac{x^2}{a^2}-\frac{y^2}{b^2}=1\Rightarrow \frac{x_0x}{a^2}-\frac{y_0y}{b^2}=1;\)\\
\(\text{抛物线 }y^2=2px\Rightarrow y_0y=p(x+x_0)\text{。}\)\\
\(\text{指出这些都可视为“切点弦方程”\text{，}考场上可直接作为解析几何结论使用\text{。}}\)

\noindent\textbf{教师追问：}\(\text{为什么我们不能只记“把 }x^2\text{ 换成 }x_0x\text{”这种口令？}\)\\

\textbf{预期回答：}\(\text{因为容易脱离曲线类型和方程形式\text{，}变式时会失误\text{。}}\)\par

\noindent\textbf{教师追问：}\(\text{抛物线中“唯一公共点”是否一定意味着切线？}\)\\

\textbf{预期回答：}\(\text{不一定\text{，}还要核查是否真满足切线条件\text{。}}\)\par

\paragraph{20分钟：例题研习一}

\(\text{选用真题： }\texttt{2019-全国III卷-理科::21}\text{。}\)\\
\(\text{教学定位：用抛物线切点弦方程体会“极线是定直线结构的来源”\text{。}}\)\\
\(\text{教师先做结构预判：点 }D\text{ 在定直线 }y=-\frac12\text{ 上运动\text{，}故切点弦 }AB\text{ 可能过定点\text{。}}\)\\
\(\text{随后回到高中解法：设 }D=(t,-\frac12)\text{\text{，}对抛物线 }y=\frac{x^2}{2}\text{ 改写为 }x^2=2y\;(p=1)\text{，}\)\\
\(\text{由切点弦方程得 }tx=y-\frac12\text{\text{，}即 }y=tx+\frac12\text{\text{，}从而 }AB\text{ 恒过定点 }\left(0,\frac12\right)\text{。}\)

\noindent\textbf{教师追问：}\(\text{这里的“动对象”是什么？“看起来可能不动”的对象是什么？}\)\\

\textbf{预期回答：}\(\text{动对象是 }D\text{ 和切点弦 }AB\text{\text{；}可能不动的是 }AB\text{ 经过的某个点\text{。}}\)\par

\noindent\textbf{教师追问：}\(\text{把 }y=\frac{x^2}{2}\text{ 改写成 }x^2=2py\text{ 后\text{，}}p\text{ 等于多少？}\)\\

\textbf{预期回答：}\(\text{ }p=1\text{。}\)\par

\noindent\textbf{教师追问：}\(\text{若点 }D=(t,-\frac12)\text{\text{，}代入切点弦公式后得到什么直线方程？}\)\\

\textbf{预期回答：}\(\text{ }tx=y-\frac12\text{\text{，}即 }y=tx+\frac12\text{。}\)\par

\noindent\textbf{教师追问：}\(\text{为什么能立刻读出定点？}\)\\

\textbf{预期回答：}\(\text{因为无论 }t\text{ 怎样变化\text{，}直线都经过 }\left(0,\frac12\right)\text{。}\)\par

\paragraph{5分钟：小结与过渡}

\(\text{第一课时收束：从调和与极线的视角\text{，}我们预判了“动点在定直线上 }\Rightarrow\text{ 切点弦过定点”\text{。}}\)\\
\(\text{第二课时转向椭圆中的“斜率和为定值 }\Rightarrow\text{ 弦过定点”与“切点弦/极线”统一\text{。}}\)

\noindent\textbf{教师追问：}\(\text{今天这一题的书面证明是否需要写交比、调和、La Hire？}\)\\

\textbf{预期回答：}\(\text{不需要\text{，}正式解答用切点弦方程即可\text{。}}\)\par

\noindent\textbf{教师追问：}\(\text{那这些概念在课堂上最大的价值是什么？}\)\\

\textbf{预期回答：}\(\text{帮助我们预判结构、统一认识、减少盲算\text{。}}\)\par

\paragraph{8分钟：复习导入}

\(\text{回顾上课时结论：若点在线上动\text{，}常可猜其切点弦或极线过定点\text{。}}\)\\
\(\text{本节把这一认识迁移到椭圆弦问题\text{。}}\)

\noindent\textbf{教师追问：}\(\text{上一课时中\text{，}哪一道真题体现了“动点在定直线上\text{，}动线过定点”？}\)\\

\textbf{预期回答：}\texttt{2019-全国III卷-理科::21}\par

\noindent\textbf{教师追问：}\(\text{它的正式方法是什么？}\)\\

\textbf{预期回答：}\(\text{切点弦方程或坐标法\text{。}}\)\par

\paragraph{20分钟：例题研习二}

\(\text{选用真题： }\texttt{2017-全国I卷-理科::20}\text{。}\)\\
\(\text{教学定位：把“手电筒模型”与极线视角对应起来\text{。}}\)\\
\(\text{第(2)问中\text{，}直线 }l\text{ 与椭圆交于 }A,B\text{\text{，}且 }P_2A,P_2B\text{ 的斜率和为 }-1\text{。}\)\\
\(\text{课堂先做结构预判：这类“定点 }P_2\text{ 发出两条射线\text{，}斜率和定值”常对应弦 }AB\text{ 过定点\text{。}}\)\\
\(\text{正式处理提示：设过 }P_2(0,1)\text{ 的两直线斜率分别为 }k_1,k_2\text{\text{，}利用题中椭圆方程联立表示交点\text{，}再由 }k_1+k_2=-1\text{ 消元\text{，}得到 }l\text{ 过定点 }(2,-1)\text{。}\)\\
\(\text{这里不在课堂上铺满全部运算\text{，}而突出“设直线参数 }\to\text{ 联立 }\to\text{ 韦达或消元 }\to\text{ 读定点”的规范路线\text{。}}\)

\noindent\textbf{教师追问：}\(\text{这道题如果完全不用射影视角\text{，}最稳妥的做法是什么？}\)\\

\textbf{预期回答：}\(\text{设参数、联立椭圆方程\text{，}用韦达定理或消元证明\text{。}}\)\par

\noindent\textbf{教师追问：}\(\text{为什么“斜率和为定值”容易让我们联想到弦 }AB\text{ 过定点？}\)\\

\textbf{预期回答：}\(\text{因为这是定点向曲线发两束射线的典型结构\text{，}常把两交点关系压缩成弦的线性不变量\text{。}}\)\par

\noindent\textbf{教师追问：}\(\text{在书面解答里\text{，}核心代数工具通常是什么？}\)\\

\textbf{预期回答：}\(\text{联立方程、韦达定理、斜率关系、消元\text{。}}\)\par

\paragraph{12分钟：对照强化}

\(\text{对照真题： }\texttt{2020-全国I卷-理科::20}\text{。}\)\\
\(\text{已知 }P\text{ 在定直线 }x=6\text{ 上\text{，}}PA,PB\text{ 分别再交椭圆于 }C,D\text{\text{，}结论是 }CD\text{ 过定点 }\left(\frac32,0\right)\text{。}\)\\
\(\text{教师不展开整题\text{，}只突出识别：这与上题同属“定点/定直线 }+\text{ 两束射线 }\Rightarrow\text{ 弦的定点性质”\text{。}}\)\\
\(\text{射影视角预判：动点 }P\text{ 在定直线\text{，}连向曲线上两个固定端型点 }A,B\text{ 的构造\text{，}往往诱导一簇弦过某定点\text{。}}\)\\
\(\text{正式方法仍是坐标设点、联立、化简\text{。}}\)

\noindent\textbf{教师追问：}\(\text{它与上一题最相似的结构关键词是什么？}\)\\

\textbf{预期回答：}\(\text{定点或定直线、两条连线、再得一条弦、证明过定点\text{。}}\)\par

\noindent\textbf{教师追问：}\(\text{若你在考场上见到这类题\text{，}第一反应应写什么？}\)\\

\textbf{预期回答：}\(\text{先设参数和直线方程\text{，}联立曲线\text{，}再用韦达定理或消元\text{。}}\)\par

\paragraph{10分钟：误区辨析}

\(\text{结合 }\texttt{2021-天津卷::18}\text{ 与 }\texttt{2021-全国乙卷-理科::21}\text{ 做辨析：}\)\\
\(\text{一是“与曲线只有一个公共点”不一定就能直接判定切线\text{；}}\)\\
\(\text{二是抛物线、椭圆、双曲线的切线/切点弦公式不能混用\text{；}}\)\\
\(\text{三是射影视角不能替代书面证明\text{。}}\)

\noindent\textbf{教师追问：}\(\text{为什么 }\texttt{2021-天津卷::18}\text{ 特别适合提醒“唯一公共点 }\not\Rightarrow\text{ 必为切线”？}\)\\

\textbf{预期回答：}\(\text{因为题目中直线与椭圆有唯一公共点\text{，}需要区分“唯一交点”与“相切”的逻辑\text{。}}\)\par

\noindent\textbf{教师追问：}\(\text{为什么 }\texttt{2021-全国乙卷-理科::21}\text{ 适合用来巩固“切点弦/切线面积问题”的解析处理？}\)\\

\textbf{预期回答：}\(\text{因为它把外点作两条切线、切点、面积最值放在同一题中\text{，}便于回到切线方程与坐标法\text{。}}\)\par

\paragraph{10分钟：当堂归纳}

\(\text{教师组织学生完成本课总结：}\)\\
\(\boxed{\text{结构识别}}\to\boxed{\text{射影视角预判}}\to\boxed{\text{高中方法落地}}\text{。}\)\\
\(\text{并形成三句口令：}\)\\
\(\text{看到外点两切线\text{，}先想切点弦\text{；}}\)\\
\(\text{看到动点在定直线\text{，}先猜极线过定点\text{；}}\)\\
\(\text{看到真题作答\text{，}马上回到坐标与韦达\text{。}}\)

\noindent\textbf{教师追问：}\(\text{本课最重要的不是哪条公式\text{，}而是哪条工作流程？}\)\\

\textbf{预期回答：}\(\text{先识别结构\text{，}再做预判\text{，}最后用高中方法证明\text{。}}\)\par

\noindent\textbf{教师追问：}\(\text{如果要把今天的方法压缩成三个动作\text{，}分别是什么？}\)\\

\textbf{预期回答：}\(\text{识别、翻译、落地\text{。}}\)\par

\paragraph{5分钟：布置作业}

\(\text{分层布置：基础巩固做 }\texttt{2019-全国III卷-理科::21}\text{ 第(1)问与切点弦方程整理\text{；}}\)\\
\(\text{提升训练做 }\texttt{2017-全国I卷-理科::20}\text{ 第(2)问完整证明\text{；}}\)\\
\(\text{拓展迁移做 }\texttt{2020-全国I卷-理科::20}\text{ 第(2)问\text{，}要求写出“结构预判 + 正式证明”双栏笔记\text{。}}\)

\noindent\textbf{教师追问：}\(\text{课后笔记中\text{，}哪一栏可以写射影视角\text{，}哪一栏必须写高考规范解答？}\)\\

\textbf{预期回答：}\(\text{预判栏可写极点极线视角\text{；}正式解答栏必须写坐标法、韦达定理等高中方法\text{。}}\)\par

\subsection*{课堂真题}

\paragraph{2019 年高考数学 全国III卷-理科第21题}

已知曲线 $C: y=\dfrac{x^2}{2}$，$D$ 为直线 $y=-\dfrac{1}{2}$ 上的动点，过 $D$ 作 $C$ 的两条切线，切点分别为 $A$，$B$。
（1）证明：直线 $AB$ 过定点；
（2）若以 $E\left(0,\dfrac{5}{2}\right)$ 为圆心的圆与直线 $AB$ 相切，且切点为线段 $AB$ 的中点，求四边形 $ADBE$ 的面积。

\textbf{答案：}（1）定点为 $\left(0,\dfrac{1}{2}\right)$；（2）四边形 $ADBE$ 的面积为 $3$ 或 $4\sqrt{2}$。\par

\textbf{解答：}（1）设 $D\left(t,-\dfrac{1}{2}\right)$，$A(x_1,y_1)$，则 $x_1^2=2y_1$。由于 $y'=x$，所以切线 $DA$ 的斜率为 $x_1$，故 $\dfrac{y_1+\frac{1}{2}}{x_1-t}=x_1$，整理得 $2tx_1-2y_1+1=0$。设 $B(x_2,y_2)$，同理可得 $2tx_2-2y_2+1=0$。故直线 $AB$ 的方程为 $2tx-2y+1=0$，所以直线 $AB$ 过定点 $\left(0,\dfrac{1}{2}\right)$。
（2）由（1）得直线 $AB$ 的方程为 $y=tx+\dfrac{1}{2}$。与 $y=\dfrac{x^2}{2}$ 联立得 $x^2-2tx-1=0$，于是 $x_1+x_2=2t$，$x_1x_2=-1$，$y_1+y_2=t(x_1+x_2)+1=2t^2+1$，$|AB|=\sqrt{1+t^2}|x_1-x_2|=\sqrt{1+t^2}\times\sqrt{(x_1+x_2)^2-4x_1x_2}=2(t^2+1)$。设 $d_1,d_2$ 分别为点 $D$，$E$ 到直线 $AB$ 的距离，则 $d_1=\sqrt{t^2+1}$，$d_2=\dfrac{2}{\sqrt{t^2+1}}$。因此四边形 $ADBE$ 的面积 $S=\dfrac{1}{2}|AB|(d_1+d_2)=(t^2+3)\sqrt{t^2+1}$。设 $M$ 为线段 $AB$ 的中点，则 $M\left(t,t^2+\dfrac{1}{2}\right)$。由于 $\overrightarrow{EM}\perp\overrightarrow{AB}$，而 $\overrightarrow{EM}=(t,t^2-2)$，$\overrightarrow{AB}$ 与向量 $(1,t)$ 平行，所以 $t+(t^2-2)t=0$，解得 $t=0$ 或 $t=\pm1$。当 $t=0$ 时，$S=3$；当 $t=\pm1$ 时，$S=4\sqrt{2}$。因此，四边形 $ADBE$ 的面积为 $3$ 或 $4\sqrt{2}$。

\paragraph{2017 年高考数学 全国I卷-理科第20题}

已知椭圆 $C:\frac{x^2}{a^2}+\frac{y^2}{b^2}=1$（$a>b>0$），四点 $P_1(1,1)$，$P_2(0,1)$，$P_3(-1,\frac{\sqrt{3}}{2})$，$P_4(1,\frac{\sqrt{3}}{2})$ 中恰有三点在椭圆 $C$ 上．
（1）求 $C$ 的方程；
（2）设直线 $l$ 不经过 $P_2$ 点且与 $C$ 相交于 $A$，$B$ 两点．若直线 $P_2A$ 与直线 $P_2B$ 的斜率的和为 $-1$，证明：$l$ 过定点．

\textbf{答案：}（1）$\frac{x^2}{4}+y^2=1$；（2）证明见解析，定点为 $(2,-1)$\par

\textbf{解答：}解：（1）根据椭圆的对称性，$P_3(-1,\frac{\sqrt{3}}{2})$，$P_4(1,\frac{\sqrt{3}}{2})$ 两点必在椭圆 $C$ 上，又 $P_4$ 的横坐标为 1，\ensuremath{\therefore}椭圆必不过 $P_1(1,1)$，\ensuremath{\therefore}$P_2(0,1)$，$P_3(-1,\frac{\sqrt{3}}{2})$，$P_4(1,\frac{\sqrt{3}}{2})$ 三点在椭圆 $C$ 上．把 $P_2(0,1)$，$P_3(-1,\frac{\sqrt{3}}{2})$ 代入椭圆 $C$，得：$\begin{cases} \frac{0}{a^2}+\frac{1}{b^2}=1 \\ \frac{1}{a^2}+\frac{3}{4b^2}=1 \end{cases}$，解得 $a^2=4$，$b^2=1$，\ensuremath{\therefore}椭圆 $C$ 的方程为 $\frac{x^2}{4}+y^2=1$．
证明：（2）(1)当斜率不存在时，设 $l:x=m$，$A(m,y_A)$，$B(m,-y_A)$，\ensuremath{\because}直线 $P_2A$ 与直线 $P_2B$ 的斜率的和为 $-1$，\ensuremath{\therefore}$\frac{y_A-1}{m-0}+\frac{-y_A-1}{m-0}=\frac{-2}{m}=-1$，解得 $m=2$，此时 $l$ 过椭圆右顶点，不存在两个交点，故不满足．(2)当斜率存在时，设 $l:y=kx+t$，($t\neq1$)，$A(x_1,y_1)$，$B(x_2,y_2)$，联立 $\begin{cases} y=kx+t \\ \frac{x^2}{4}+y^2=1 \end{cases}$，整理，得 $(1+4k^2)x^2+8ktx+4t^2-4=0$，$\Delta>0$，$x_1+x_2=-\frac{8kt}{1+4k^2}$，$x_1x_2=\frac{4t^2-4}{1+4k^2}$，则 $k_{P_2A}+k_{P_2B}=\frac{y_1-1}{x_1}+\frac{y_2-1}{x_2}=\frac{kx_1+t-1}{x_1}+\frac{kx_2+t-1}{x_2}=2k+\frac{t-1}{x_1}+\frac{t-1}{x_2}=2k+(t-1)\frac{x_1+x_2}{x_1x_2}=2k+(t-1)\frac{-\frac{8kt}{1+4k^2}}{\frac{4t^2-4}{1+4k^2}}=2k-\frac{2k(t-1)}{t+1}=\frac{2k(t+1)-2k(t-1)}{t+1}=\frac{4k}{t+1}=-1$，又 $t\neq1$，\ensuremath{\therefore}$t=-2k-1$，此时 $\Delta=64k^2-16(1+4k^2)(t^2-1)$，代入得 $\Delta=-64k$，存在 $k$，使得 $\Delta>0$ 成立，\ensuremath{\therefore}直线 $l$ 的方程为 $y=kx-2k-1$，当 $x=2$ 时，$y=-1$，\ensuremath{\therefore}$l$ 过定点 $(2,-1)$．

\paragraph{2020 年高考数学 全国I卷-理科第20题}

已知$A,B$分别为椭圆$E:\frac{x^2}{a^2}+y^2=1(a>1)$的左、右顶点，$G$为$E$的上顶点，$\overrightarrow{AG}\cdot\overrightarrow{GB}=8$。$P$为直线$x=6$上的动点，$PA$与$E$的另一交点为$C$，$PB$与$E$的另一交点为$D$。
(1)求$E$的方程；
(2)证明：直线$CD$过定点。

\textbf{答案：}(1)$\frac{x^2}{9}+y^2=1$；(2)直线$CD$过定点$\left(\frac{3}{2},0\right)$\par

\textbf{解答：}(1)由椭圆方程得$A(-a,0),B(a,0),G(0,1)$，则$\overrightarrow{AG}=(a,1),\overrightarrow{GB}=(a,-1)$，$\overrightarrow{AG}\cdot\overrightarrow{GB}=a^2-1=8$，解得$a^2=9$，故椭圆方程为$\frac{x^2}{9}+y^2=1$。
(2)设$P(6,y_0)$，则直线$AP$方程为$y=\frac{y_0}{9}(x+3)$，与椭圆联立得$(y_0^2+9)x^2+6y_0^2x+9y_0^2-81=0$，由韦达定理$x_A\cdot x_C=\frac{9y_0^2-81}{y_0^2+9}$，而$x_A=-3$，得$x_C=\frac{-3y_0^2+27}{y_0^2+9}$，代入得$y_C=\frac{6y_0}{y_0^2+9}$，故$C\left(\frac{-3y_0^2+27}{y_0^2+9},\frac{6y_0}{y_0^2+9}\right)$。同理直线$PB$方程$y=\frac{y_0}{3}(x-3)$，与椭圆联立得$D\left(\frac{3y_0^2-3}{y_0^2+1},\frac{-2y_0}{y_0^2+1}\right)$。当$y_0\neq0$时，直线$CD$斜率$k=\frac{\frac{6y_0}{y_0^2+9}-\frac{-2y_0}{y_0^2+1}}{\frac{-3y_0^2+27}{y_0^2+9}-\frac{3y_0^2-3}{y_0^2+1}}=\frac{4y_0}{3(3-y_0^2)}$，直线$CD$方程为$y-\frac{-2y_0}{y_0^2+1}=\frac{4y_0}{3(3-y_0^2)}\left(x-\frac{3y_0^2-3}{y_0^2+1}\right)$，化简得$y=\frac{4y_0}{3(3-y_0^2)}\left(x-\frac{3}{2}\right)$，故过定点$\left(\frac{3}{2},0\right)$。当$y_0=0$时，$C(3,0),D(3,0)$，直线$CD$为$x=3$，也过$\left(\frac{3}{2},0\right)$。综上，直线$CD$过定点$\left(\frac{3}{2},0\right)$。

\subsection*{课后作业}

\begin{enumerate}

\item 2019 年高考数学 全国III卷-理科第21题

\item 2017 年高考数学 全国I卷-理科第20题

\item 2020 年高考数学 全国I卷-理科第20题

\item 2021 年高考数学 全国乙卷-理科第21题

\item 2021 年高考数学 天津卷第18题

\end{enumerate}

\noindent\textbf{作业答案索引：}

\begin{itemize}

\item 2019 年高考数学 全国III卷-理科第21题：（1）定点为 $\left(0,\dfrac{1}{2}\right)$；（2）四边形 $ADBE$ 的面积为 $3$ 或 $4\sqrt{2}$。

\item 2017 年高考数学 全国I卷-理科第20题：（1）$\frac{x^2}{4}+y^2=1$；（2）证明见解析，定点为 $(2,-1)$

\item 2020 年高考数学 全国I卷-理科第20题：(1)$\frac{x^2}{9}+y^2=1$；(2)直线$CD$过定点$\left(\frac{3}{2},0\right)$

\item 2021 年高考数学 全国乙卷-理科第21题：（1）$p=2$；（2）$20\sqrt{5}$

\item 2021 年高考数学 天津卷第18题：（I）$\frac{x^2}{5} + y^2 = 1$；（II）$x - y + \sqrt{6} = 0$

\end{itemize}

\subsection*{板书设计}

\textbf{板书一：主线与边界}\\
\(\text{交比}\rightarrow\text{调和}\rightarrow\text{极点极线}\rightarrow\text{定点定直线}\)\\
\(\text{射影几何：只作结构预判\text{；}书面解答：回到高中方法}\)\\[4pt]
\textbf{板书二：核心概念}\\
\((A,B;C,D)=\frac{\overrightarrow{CA}/\overrightarrow{CB}}{\overrightarrow{DA}/\overrightarrow{DB}},\qquad (A,B;C,D)=-1\)\\
\(X\in \ell_P\iff P\in \ell_X\)\\[4pt]
\textbf{板书三：常用方程}\\
\(\frac{x^2}{a^2}+\frac{y^2}{b^2}=1\Rightarrow \frac{x_0x}{a^2}+\frac{y_0y}{b^2}=1\)\\
\(\frac{x^2}{a^2}-\frac{y^2}{b^2}=1\Rightarrow \frac{x_0x}{a^2}-\frac{y_0y}{b^2}=1\)\\
\(x^2=2py\Rightarrow xx_0=p(y+y_0),\qquad y^2=2px\Rightarrow yy_0=p(x+x_0)\)\\[4pt]
\textbf{板书四：真题模型提炼}\\
\texttt{2019-全国III卷-理科::21}:\; \(D\in y=-\frac12\Rightarrow AB\text{ 过定点 }\left(0,\frac12\right)\)\\
\texttt{2017-全国I卷-理科::20}:\; \(k_{P_2A}+k_{P_2B}=-1\Rightarrow l\text{ 过定点 }(2,-1)\)\\
\texttt{2020-全国I卷-理科::20}:\; \(P\in x=6\Rightarrow CD\text{ 过定点 }\left(\frac32,0\right)\)\\[4pt]
\textbf{板书五：规范流程}\\
\(\text{识别结构}\to\text{预判结论}\to\text{设参联立}\to\text{韦达/消元}\to\text{读出定点或定直线}\)

\subsection*{易错点}

\begin{itemize}

\item \(\text{把射影视角当作正式证明\text{。}纠正：课堂可用来预判\text{，}考卷必须写成坐标法、定义法、韦达定理、点差法或平面几何\text{。}}\)

\item \(\text{机械记忆“把 }x^2\text{ 换成 }x_0x\text{”\text{。}纠正：必须先分清椭圆、双曲线、抛物线及其标准形式\text{。}}\)

\item \(\text{把“唯一公共点”直接等同于“切线”\text{。}纠正：需核查切线定义或判别条件\text{，}参照 }\texttt{2021-天津卷::18}\text{ 的辨析\text{。}}\)

\item \(\text{调和点列中忽视有向线段\text{。}纠正：只要求学生在本轮复习中会识别 }(A,B;C,D)=-1\text{ 的结构\text{，}不做无符号乱算\text{。}}\)

\item \(\text{见到定点题就盲目大算\text{。}纠正：先问“谁在定直线上动”“谁是切点弦/弦”“哪一个对象可能固定”\text{。}}\)

\item \(\text{抛物线切点弦公式混淆 }x^2=2py\text{ 与 }y^2=2px\text{\text{。}纠正：板书并要求学生口头区分\text{。}}\)

\end{itemize}

\subsection*{教师备注}

\textbf{1. 教学定位}\\
\(\text{本模块面向一轮复习\text{，}不追求射影证明训练\text{，}而追求“会识别、会翻译、会落地”\text{。}学生只需形成高频结构意识\text{。}}\)\\
\textbf{2. 教学节奏}\\
\(\text{第一课时务必把 }\texttt{2019-全国III卷-理科::21}\text{ 讲透\text{，}因为它最适合作为“极线 }\to\text{ 切点弦 }\to\text{ 定点”的入门样板\text{。}}\)\\
\(\text{第二课时把 }\texttt{2017-全国I卷-理科::20}\text{ 与 }\texttt{2020-全国I卷-理科::20}\text{ 放在一起对照\text{，}强化“手电筒模型”与“弦过定点”的统一性\text{。}}\)\\
\textbf{3. 板演取舍}\\
\(\text{不建议课堂上完整推导交比不变性\text{；}可只用一句话说明“它是投影下保留下来的比例关系”\text{，}把时间让给真题迁移\text{。}}\)\\
\textbf{4. 规范提醒}\\
\(\text{学生在笔记中可分“结构预判”与“正式作答”两栏：左栏写 }\mathrm{La\ Hire}\text{、极点极线、调和\text{；}右栏必须改写成解析几何步骤\text{。}}\)\\
\textbf{5. 课堂追问建议}\\
\(\text{若学生会公式不会用\text{，}可追问：}\;\text{“谁在动？”“谁可能固定？”“题目给的是点的条件还是线的条件？”“能否先猜再证？”}\)\\
\textbf{6. 评价标准}\\
\(\text{本课达标不看学生是否会说复杂术语\text{，}而看其是否能做到：看到切点弦会写方程\text{，}看到斜率和定值会想到弦的定点\text{，}看到结构后能回到规范代数证明\text{。}}\)

\clearpage

\section{第3课：三种定义与标准方程}

\noindent 建议课时：2课时。

\subsection*{教学目标}

\begin{itemize}

\item 能准确表述并辨认圆锥曲线第一定义：椭圆满足 $|PF_1|+|PF_2|=2a$，双曲线满足 $\bigl||PF_1|-|PF_2|\bigr|=2a$，抛物线满足 $|PF|=d(P,l)$.

\item 能把第一定义中的数量关系转化为标准方程参数关系：椭圆中 $c^2=a^2-b^2$，双曲线中 $c^2=a^2+b^2$，并结合焦点位置写出标准方程.

\item 能用焦点—准线统一定义 $\dfrac{|PF|}{d(P,l)}=e$ 判断曲线类型，并理解 $e<1,e=1,e>1$ 分别对应椭圆、抛物线、双曲线.

\item 能在高中范围内理解第三定义的可用形式：当两定点为顶点时，利用斜率之积为常数判断轨迹类型或求离心率，并能回到坐标法完成规范解答.

\item 能在高考题情境中优先识别定义结构，再选择坐标法、定义法、韦达定理、点差法、向量或平面几何完成求解.

\end{itemize}

\subsection*{重点与难点}

\begin{itemize}

\item 由“距离和为定值”“距离差的绝对值为定值”“点到定点与定直线距离相等”快速识别第一定义.

\item 由 $e=\dfrac{|PF|}{d(P,l)}$ 统一三类圆锥曲线，并能从统一定义退回到标准方程.

\item 第三定义只作结构识别：若与关于中心对称的两定点连线斜率之积为定值，则常对应中心二次曲线；正式求解仍应设点坐标化简.

\item 标准方程书写时必须先判焦点轴，再定分母正负与参数位置.

\end{itemize}

\subsection*{高观点}

从统一视角看，圆锥曲线都可视作同一类二次曲线在不同度量条件下的表现：第一定义强调“到定点、定直线的距离关系”，第二定义用离心率 $e$ 统一分类，第三定义则暴露出中心二次曲线与一组对称定点之间的射影型结构联系。课堂中只用这一视角帮助学生预判“题目中的数量关系大概率对应哪一类曲线”，不把射影结论作为高考正式证明依据。

\subsection*{高中落地}

高中层面的可操作表达是：见到题目中的根式、线段长、点到直线距离、斜率积，先翻译为坐标中的可计算量；若能直接对应定义，就先判曲线类型，再由 $a,b,c,e$ 的关系写标准方程；若出现第三定义型信息，只把它当作“优先怀疑是椭圆或双曲线”的信号，随后必须设点 $P(x,y)$，把斜率积化为方程，或由标准方程反推离心率。

\subsection*{教学流程}

\paragraph{5分钟：导入与目标明示}

教师板书本课主题“三种定义与标准方程”，说明本课核心不是死记方程，而是把题目中的条件识别为轨迹定义，再回到高中方法求解。

\noindent\textbf{教师追问：}看到“距离和为定值”时，你首先想到什么曲线？\\

\textbf{预期回答：}椭圆，但还要比较定值与两定点间距离的大小关系。\par

\noindent\textbf{教师追问：}看到“到定点与到定直线距离相等”时，你首先想到什么曲线？\\

\textbf{预期回答：}抛物线，定点是焦点，定直线是准线。\par

\paragraph{12分钟：第一定义系统梳理}

归纳三类曲线的第一定义，并强调存在性条件：椭圆中 $2a>|F_1F_2|=2c$，即 $a>c$；双曲线中 $0<2a<2c$，即 $0<a<c$；抛物线要求焦点不在准线上。随后把定义与标准方程对应：椭圆 $\dfrac{x^2}{a^2}+\dfrac{y^2}{b^2}=1$ 或 $\dfrac{y^2}{a^2}+\dfrac{x^2}{b^2}=1$，双曲线 $\dfrac{x^2}{a^2}-\dfrac{y^2}{b^2}=1$ 或 $\dfrac{y^2}{a^2}-\dfrac{x^2}{b^2}=1$，抛物线 $y^2=2px,\ y^2=-2px,\ x^2=2py,\ x^2=-2py$.

\noindent\textbf{教师追问：}为什么椭圆里必须有 $2a>2c$？\\

\textbf{预期回答：}由三角形两边之和大于第三边，点到两焦点距离和必须大于焦距。\par

\noindent\textbf{教师追问：}若 $|PF_1|+|PF_2|=|F_1F_2|$，轨迹还是椭圆吗？\\

\textbf{预期回答：}不是，退化为焦点连线间的线段。\par

\noindent\textbf{教师追问：}双曲线中为什么是距离差的绝对值？\\

\textbf{预期回答：}因为左右两支上谁大谁小不固定，但差的绝对值保持常数。\par

\paragraph{10分钟：题例切入：第一定义识别}

引用题目 $\text{2016-全国I卷-理科::20}$ 的第（I）问，只做定义识别训练：题干要求证明 $|EA|+|EB|$ 为定值并求轨迹方程。教师引导学生先不急于列坐标，先判断“和为定值”对应椭圆第一定义，再结合题目答案中的轨迹方程 $\dfrac{x^2}{4}+\dfrac{y^2}{3}=1\ (y\ne 0)$ 反查其焦点与长轴信息。

\noindent\textbf{教师追问：}这一问最先暴露出来的轨迹信号是什么？\\

\textbf{预期回答：}$|EA|+|EB|$ 为定值，对应椭圆第一定义。\par

\noindent\textbf{教师追问：}若最终是椭圆，定点最可能是谁？\\

\textbf{预期回答：}最可能是 $A,B$ 两个定点作为焦点。\par

\noindent\textbf{教师追问：}为什么方程后面还要写 $y\ne 0$？\\

\textbf{预期回答：}因为原题中的轨迹来自构造点 $E$，不一定包含退化或特殊位置，要结合题设去掉不满足几何构造的点。\par

\paragraph{8分钟：抛物线第一定义即时训练}

引用 $\text{2023-北京卷::6}$ 与 $\text{2022-全国乙卷-理科::5}$，突出“抛物线点到焦点距离可直接转为点到准线距离”。对 $y^2=8x$，有 $p=4$，焦点为 $F(2,0)$，准线为 $x=-2$；若点 $M$ 到直线 $x=-3$ 的距离已知，则需先判断该信息能否通过横坐标转化，再与准线距离建立关系。对 $y^2=4x$，焦点为 $F(1,0)$，点 $A$ 在抛物线上且 $|AF|=|BF|$，可把 $|AF|$ 转为点 $A$ 到准线 $x=-1$ 的距离，再与点 $B(3,0)$ 的焦点距离比较。

\noindent\textbf{教师追问：}抛物线中最常用的定义转化是什么？\\

\textbf{预期回答：}把点到焦点距离转化为点到准线距离，或者反过来转化。\par

\noindent\textbf{教师追问：}对 $y^2=4x$，为什么把 $|AF|$ 改写成点到准线的距离会更方便？\\

\textbf{预期回答：}因为准线平行于 $y$ 轴，点到准线的距离可以直接用横坐标表示。\par

\paragraph{10分钟：第一课时小结与当堂检测}

学生口头完成三个判断：若 $|PF_1|+|PF_2|=10,|F_1F_2|=12$，则无椭圆轨迹；若 $\bigl||PF_1|-|PF_2|\bigr|=10,|F_1F_2|=12$，则为双曲线；若 $|PF|=d(P,l)$，则为抛物线。教师强调“先判定，再求参量，再写方程”的解题顺序。

\noindent\textbf{教师追问：}求轨迹类题第一步最忌讳什么？\\

\textbf{预期回答：}一上来就盲目设点硬算，不先识别定义结构。\par

\noindent\textbf{教师追问：}今天关于第一定义，你认为最关键的一句操作性结论是什么？\\

\textbf{预期回答：}先把条件翻译成点到定点、定直线的距离关系。\par

\paragraph{6分钟：第二课时复习衔接}

回顾第一定义后，过渡到第二定义：把“和、差、相等”提升为统一比值 $e=\dfrac{|PF|}{d(P,l)}$。说明该定义的价值在于统一分类，而不是取代标准方程。

\noindent\textbf{教师追问：}第一定义和第二定义的关系是什么？\\

\textbf{预期回答：}第一定义分别描述三类曲线，第二定义用同一形式统一描述。\par

\noindent\textbf{教师追问：}为什么说第二定义更统一？\\

\textbf{预期回答：}因为只看离心率 $e$ 的范围就能区分椭圆、抛物线、双曲线。\par

\paragraph{12分钟：焦点准线统一定义讲解}

板书第二定义：$\dfrac{|PF|}{d(P,l)}=e$。分类为 $0<e<1$ 时椭圆，$e=1$ 时抛物线，$e>1$ 时双曲线。补充标准方程中的离心率：椭圆 $e=\dfrac{c}{a}$，双曲线 $e=\dfrac{c}{a}$ 且 $e>1$。强调高考中常用方向是“已知曲线，反向把焦点距离转成准线距离”。

\noindent\textbf{教师追问：}椭圆和双曲线的离心率公式表面上都写成 $e=\dfrac{c}{a}$，本质区别是什么？\\

\textbf{预期回答：}椭圆中 $c<a$ 所以 $e<1$，双曲线中 $c>a$ 所以 $e>1$。\par

\noindent\textbf{教师追问：}在抛物线题里，第二定义常怎样使用？\\

\textbf{预期回答：}通常是把点到焦点距离与点到准线距离互相代换。\par

\paragraph{8分钟：题例应用：准线转化}

继续利用 $\text{2023-北京卷::6}$、$\text{2025-全国II卷::6}$。后者中已知抛物线 $C:y^2=2px\ (p>0)$，过点 $A$ 作准线垂线，垂足为 $B$，又给出直线 $BF$ 的方程。教师引导学生只抓住一点：因为 $B$ 是点 $A$ 到准线的垂足，所以 $|AB|=d(A,l)=|AF|$，题目实质是在借准线构造等长关系。

\noindent\textbf{教师追问：}在 $\text{2025-全国II卷::6}$ 中，为什么垂足 $B$ 很关键？\\

\textbf{预期回答：}它把点到准线的距离具体化为线段 $AB$。\par

\noindent\textbf{教师追问：}这里真正要用的等量关系是什么？\\

\textbf{预期回答：}$|AF|=|AB|$。\par

\paragraph{12分钟：第三定义的高中化表达}

引入第三定义的高中化表述。若中心在原点、左右顶点为 $A(-a,0),B(a,0)$，设 $P(x,y)$，则
对椭圆 $\dfrac{x^2}{a^2}+\dfrac{y^2}{b^2}=1$，有
$k_{PA}k_{PB}=\dfrac{y}{x+a}\cdot\dfrac{y}{x-a}=\dfrac{y^2}{x^2-a^2}=-\dfrac{b^2}{a^2}$；
对双曲线 $\dfrac{x^2}{a^2}-\dfrac{y^2}{b^2}=1$，有
$k_{PA}k_{PB}=\dfrac{y^2}{x^2-a^2}=\dfrac{b^2}{a^2}$。
据此说明：斜率之积为负常数常对应椭圆，为正常数常对应双曲线。强调这不是额外负担，而是“设点坐标后自然算出来”的结构结论。

\noindent\textbf{教师追问：}为什么第三定义适合放在中心曲线、且两定点常取顶点？\\

\textbf{预期回答：}因为这样坐标表达最简，$x\pm a$ 结构明显，容易化成标准方程。\par

\noindent\textbf{教师追问：}斜率之积的正负分别在提示什么？\\

\textbf{预期回答：}负值常提示椭圆，正值常提示双曲线。\par

\paragraph{10分钟：题例应用：第三定义与离心率}

引用 $\text{2022-全国甲卷-理科::10}$ 与 $\text{2019-全国II卷-理科::21}$。前者中左顶点为 $A$，点 $P,Q$ 关于 $y$ 轴对称，且直线 $AP,AQ$ 的斜率之积为 $\dfrac14$。教师引导学生注意：这是“上、下或左右顶点+斜率积定值”的典型结构，可优先联想到第三定义求离心率。后者则直接给出 $A(-2,0),B(2,0)$ 与斜率积 $-\dfrac12$，可设 $M(x,y)$，由
$\dfrac{y}{x+2}\cdot\dfrac{y}{x-2}=-\dfrac12$
化得
$2y^2=-(x^2-4)$，即 $\dfrac{x^2}{4}+\dfrac{y^2}{2}=1$，从而说明轨迹是椭圆。

\noindent\textbf{教师追问：}在 $\text{2019-全国II卷-理科::21}$ 中，为什么正式解答必须设点坐标推方程？\\

\textbf{预期回答：}因为高考规范答案要落在坐标法与标准方程上，第三定义只能帮助预判。\par

\noindent\textbf{教师追问：}由 $k_{AM}k_{BM}=-\dfrac12$ 直接能看出什么？\\

\textbf{预期回答：}能先预判轨迹更像椭圆，再通过坐标化验证。\par

\paragraph{10分钟：综合提升：第一定义与标准方程联动}

引用 $\text{2021-上海春季卷::19}$ 与 $\text{2019-全国II卷-文科::20}$。前者第（1）问中 $|PA|-|PB|=20$ 千米，焦点位于 $A,B$，先用双曲线第一定义得到 $2a=20$，再由已知站点位置确定 $2c=40$，从而求出 $a=10,c=20,b^2=c^2-a^2=300$，得到标准方程 $\dfrac{x^2}{100}-\dfrac{y^2}{300}=1$。后者则反向从椭圆标准方程出发，利用焦点、离心率和几何条件求参数，帮助学生建立“定义与方程双向互通”的观念。

\noindent\textbf{教师追问：}在 $\text{2021-上海春季卷::19}$ 中，为什么先得到的是 $a,c$ 而不是 $b$？\\

\textbf{预期回答：}因为定义直接给出焦距和实轴参数，$b$ 需由 $b^2=c^2-a^2$ 再求。\par

\noindent\textbf{教师追问：}标准方程和定义之间最重要的双向通道是什么？\\

\textbf{预期回答：}由定义判类型和参数，再写方程；也能由方程读焦点、离心率和定义量。\par

\paragraph{5分钟：课堂总结}

教师总结本课三句主线：一是“见距离先想定义”；二是“见离心率想统一定义”；三是“见斜率积想第三定义预判，再回到坐标法”。提醒学生书写方程时要同时说明曲线名称与必要范围限制。

\noindent\textbf{教师追问：}今天三种定义中，哪一种最适合做统一框架？\\

\textbf{预期回答：}第二定义。\par

\noindent\textbf{教师追问：}哪一种最容易在选择题、填空题中快速得分？\\

\textbf{预期回答：}第一定义和抛物线的焦点准线转化。\par

\noindent\textbf{教师追问：}哪一种最需要注意不能只停留在观察层面？\\

\textbf{预期回答：}第三定义，必须回到坐标法完成规范证明。\par

\subsection*{课堂真题}

\paragraph{2016 年高考数学 全国I卷-理科第20题}

设圆$x^2+y^2+2x-15=0$的圆心为$A$，直线$l$过点$B(1,0)$且与$x$轴不重合，$l$交圆$A$于$C$，$D$两点，过$B$作$AC$的平行线交$AD$于点$E$．
（I）证明$|EA|+|EB|$为定值，并写出点$E$的轨迹方程；
（II）设点$E$的轨迹为曲线$C_1$，直线$l$交$C_1$于$M$，$N$两点，过$B$且与$l$垂直的直线与圆$A$交于$P$，$Q$两点，求四边形$MPNQ$面积的取值范围．

\textbf{答案：}（I）$|EA|+|EB|=4$，轨迹方程为$\frac{x^2}{4}+\frac{y^2}{3}=1$（$y\neq0$）；（II）$[12,8\sqrt{3})$\par

\textbf{解答：}解：（I）证明：圆$x^2+y^2+2x-15=0$即为$(x+1)^2+y^2=16$，可得圆心$A(-1,0)$，半径$r=4$，由$BE\parallel AC$，可得$\angle C=\angle EBD$，由$AC=AD$，可得$\angle D=\angle C$，即为$\angle D=\angle EBD$，即有$EB=ED$，则$|EA|+|EB|=|EA|+|ED|=|AD|=4$，故$E$的轨迹为以$A$，$B$为焦点的椭圆，且有$2a=4$，即$a=2$，$c=1$，$b=\sqrt{a^2-c^2}=\sqrt{3}$，则点$E$的轨迹方程为$\frac{x^2}{4}+\frac{y^2}{3}=1$（$y\neq0$）；（II）椭圆$C_1:\frac{x^2}{4}+\frac{y^2}{3}=1$，设直线$l:x=my+1$，由$PQ\perp l$，设$PQ:y=-m(x-1)$，由$\begin{cases}\frac{x^2}{4}+\frac{y^2}{3}=1\\x=my+1\end{cases}$可得$(3m^2+4)y^2+6my-9=0$，设$M(x_1,y_1)$，$N(x_2,y_2)$，可得$y_1+y_2=-\frac{6m}{3m^2+4}$，$y_1y_2=-\frac{9}{3m^2+4}$，则$|MN|=\sqrt{1+m^2}\cdot|y_1-y_2|=\sqrt{1+m^2}\cdot\sqrt{(y_1+y_2)^2-4y_1y_2}=\sqrt{1+m^2}\cdot\sqrt{\frac{36m^2}{(3m^2+4)^2}+\frac{36}{3m^2+4}}=12\cdot\frac{1+m^2}{3m^2+4}$，$A$到$PQ$的距离为$d=\frac{|0+m(-1-1)|}{\sqrt{1+m^2}}=\frac{|2m|}{\sqrt{1+m^2}}$，$|PQ|=2\sqrt{r^2-d^2}=2\sqrt{16-\frac{4m^2}{1+m^2}}=4\sqrt{\frac{3m^2+4}{1+m^2}}$，则四边形$MPNQ$面积为$S=\frac{1}{2}|PQ|\cdot|MN|=\frac{1}{2}\cdot4\sqrt{\frac{3m^2+4}{1+m^2}}\cdot12\cdot\frac{1+m^2}{3m^2+4}=24\sqrt{\frac{1+m^2}{3m^2+4}}=24\sqrt{\frac{1}{3+\frac{1}{1+m^2}}}$，当$m=0$时，$S$取得最小值$12$，又$\frac{1}{1+m^2}>0$，可得$S<24\sqrt{\frac{1}{3}}=8\sqrt{3}$，即有四边形$MPNQ$面积的取值范围是$[12,8\sqrt{3})$．

\paragraph{2023 年高考数学 北京卷第6题}

已知抛物线 $C: y^2 = 8x$ 的焦点为 $F$ ，点 $M$ 在 $C$ 上．若 $M$ 到直线 $x = -3$ 的距离为 $5$，则 $|MF| =$ （ ）

\textbf{答案：}D\par

\textbf{解答：}因为抛物线 $C: y^2 = 8x$ 的焦点 $F(2,0)$ ，准线方程为 $x = -2$ ，点 $M$ 在 $C$ 上，所以 $M$ 到准线 $x = -2$ 的距离为 $|MF|$，又 $M$ 到直线 $x = -3$ 的距离为 $5$，所以 $|MF| + 1 = 5$，故 $|MF| = 4$.

\paragraph{2022 年高考数学 全国乙卷-理科第5题}

设$F$为抛物线$C: y^2 = 4x$的焦点，点$A$在$C$上，点$B(3, 0)$，若$|AF| = |BF|$，则$|AB| =$（\qquad）

\textbf{答案：}B\par

\textbf{解答：}由题意得，$F(1,0)$，则$|AF| = |BF| = 2$，\\即点$A$到准线$x=-1$的距离为2，所以点$A$的横坐标为$-1+2=1$，\\不妨设点$A$在$x$轴上方，代入得，$A(1,2)$，\\所以$|AB| = \sqrt{(3-1)^2 + (0-2)^2} = 2\sqrt{2}$.

\paragraph{2025 年高考数学 全国II卷第6题}

设抛物线 $C:y^2=2px(p>0)$ 的焦点为 $F$，点 $A$ 在 $C$ 上，过 $A$ 作 $C$ 的准线的垂线，垂足为 $B$，若直线 $BF$ 的方程为 $y=-2x+2$，则 $|AF|=$（ ）

\textbf{答案：}$C$\par

\textbf{解答：}对 $l_{BF}:y=-2x+2$，令 $y=0$，则 $x=1$，所以 $F(1,0)$，$p=2$ 即抛物线 $C:y^2=4x$，故抛物线的准线方程为 $x=-1$，故 $B(-1,4)$，则 $y_A=4$，代入抛物线 $C:y^2=4x$ 得 $x_A=4$，所以 $AF=AB=x_A+\frac{p}{2}=4+1=5$.

\paragraph{2022 年高考数学 全国甲卷-理科第10题}

椭圆 $C:\frac{x^2}{a^2}+\frac{y^2}{b^2}=1(a>b>0)$ 的左顶点为A，点P，Q均在C上，且关于y轴对称. 若直线AP, AQ的斜率之积为 $\frac{1}{4}$，则C的离心率为（ ）

\textbf{答案：}A\par

\textbf{解答：}设$P(x_1,y_1)$，则$Q(-x_1,y_1)$，则由$k_{AP}\cdot k_{AQ}=\frac{1}{4}$得：$k_{AP}\cdot k_{AQ}=\frac{y_1}{x_1+a}\cdot\frac{y_1}{-x_1+a}=\frac{y_1^2}{-x_1^2+a^2}=\frac{1}{4}$. 由$\frac{x_1^2}{a^2}+\frac{y_1^2}{b^2}=1$得$y_1^2=\frac{b^2(a^2-x_1^2)}{a^2}$，代入得$\frac{b^2(a^2-x_1^2)}{a^2(-x_1^2+a^2)}=\frac{1}{4}$，即$\frac{b^2}{a^2}=\frac{1}{4}$，所以椭圆C的离心率$e=\frac{c}{a}=\sqrt{1-\frac{b^2}{a^2}}=\sqrt{1-\frac{1}{4}}=\frac{\sqrt{3}}{2}$. 故选A.

\paragraph{2019 年高考数学 全国II卷-理科第21题}

已知点 $A(-2,0)$，$B(2,0)$，动点 $M(x,y)$ 满足直线 $AM$ 与 $BM$ 的斜率之积为 $-\frac{1}{2}$．记 $M$ 的轨迹为曲线 $C$．
（1）求 $C$ 的方程，并说明 $C$ 是什么曲线；
（2）过坐标原点的直线交 $C$ 于 $P$，$Q$ 两点，点 $P$ 在第一象限，$PE\perp x$ 轴，垂足为 $E$，连结 $QE$ 并延长交 $C$ 于点 $G$．
     （i）证明：$\triangle PQG$ 是直角三角形；
     （ii）求 $\triangle PQG$ 面积的最大值．

\textbf{答案：}（1）$\frac{x^2}{4}+\frac{y^2}{2}=1\ (|x|\neq 2)$，$C$ 为中心在坐标原点，焦点在 $x$ 轴上的椭圆，不含左右顶点；（2）（i）证明见解析；（ii）$\frac{16}{9}$\par

\textbf{解答：}（1）直线 $AM$ 的斜率为 $\frac{y}{x+2}\ (x\neq -2)$，直线 $BM$ 的斜率为 $\frac{y}{x-2}\ (x\neq 2)$，由题意可知 $\frac{y}{x+2}\cdot\frac{y}{x-2}=-\frac{1}{2}$，化简得 $\frac{x^2}{4}+\frac{y^2}{2}=1\ (x\neq\pm 2)$，所以 $C$ 为中心在坐标原点，焦点在 $x$ 轴上的椭圆，不含左右顶点．
（2）（i）设直线 $PQ$ 的斜率为 $k$，则其方程为 $y=kx\ (k>0)$．由 $\begin{cases} y=kx \\ \frac{x^2}{4}+\frac{y^2}{2}=1 \end{cases}$ 得 $x=\pm\frac{2}{\sqrt{1+2k^2}}$．记 $u=\frac{2}{\sqrt{1+2k^2}}$，则 $P(u,uk)$，$Q(-u,-uk)$，$E(u,0)$．于是直线 $QG$ 的斜率为 $\frac{k}{2}$，方程为 $y=\frac{k}{2}(x-u)$．由 $\begin{cases} y=\frac{k}{2}(x-u) \\ \frac{x^2}{4}+\frac{y^2}{2}=1 \end{cases}$ 得 $(2+k^2)x^2-2uk^2x+k^2u^2-8=0$．(1) 设 $G(x_G,y_G)$，则 $-u$ 和 $x_G$ 是方程(1)的解，故 $x_G=\frac{u(3k^2+2)}{2+k^2}$，$y_G=\frac{uk^3}{2+k^2}$．从而直线 $PG$ 的斜率为 $\frac{\frac{uk^3}{2+k^2}-uk}{\frac{u(3k^2+2)}{2+k^2}-u}=-\frac{1}{k}$．所以 $PQ\perp PG$，即 $\triangle PQG$ 是直角三角形．
（ii）由（i）得 $|PQ|=2u\sqrt{1+k^2}$，$|PG|=\frac{2uk\sqrt{1+k^2}}{2+k^2}$，所以 $\triangle PQG$ 的面积 $S=\frac{1}{2}|PQ||PG|=\frac{8k(1+k^2)}{(1+2k^2)(2+k^2)}=\frac{8(\frac{1}{k}+k)}{1+2(\frac{1}{k}+k)^2}$．设 $t=k+\frac{1}{k}$，则由 $k>0$ 得 $t\geq 2$，当且仅当 $k=1$ 时取等号．因为 $S=\frac{8t}{1+2t^2}$ 在 $[2,+\infty)$ 单调递减，所以当 $t=2$，即 $k=1$ 时，$S$ 取得最大值 $\frac{16}{9}$．因此，$\triangle PQG$ 面积的最大值为 $\frac{16}{9}$．

\paragraph{2021 年高考数学 上海春季卷第19题}

（1）团队在 $O$ 点西侧、东侧20千米处设有 $A$、$B$ 两站点，测量距离发现一点 $P$ 满足 $|PA|-|PB|=20$ 千米，可知 $P$ 在 $A$、$B$ 为焦点的双曲线上，以 $O$ 点为原点，东侧为 $x$ 轴正半轴，北侧为 $y$ 轴正半轴，建立平面直角坐标系，$P$ 在北偏东 $60^\circ$ 处，求双曲线标准方程和 $P$ 点坐标．\newline（2）团队又在南侧、北侧15千米处设有 $C$、$D$ 两站点，测量距离发现 $|QA|-|QB|=30$ 千米，$|QC|-|QD|=10$ 千米，求 $|OQ|$（精确到1米）和 $Q$ 点位置（精确到1米，$1^\circ$）

\textbf{答案：}$(1)\ \frac{x^2}{100}-\frac{y^2}{300}=1,\ P(\frac{15\sqrt{2}}{2},\frac{5\sqrt{6}}{2})$；$(2)\ |OQ|\approx 19$ 米，$Q$ 点位置北偏东 $66^\circ$\par

\textbf{解答：}（1）由题意可得 $a=10$，$c=20$，所以 $b^2=300$，所以双曲线的标准方程为 $\frac{x^2}{100}-\frac{y^2}{300}=1$，直线 $OP:y=\frac{\sqrt{3}}{3}x$，联立双曲线方程，可得 $x=\frac{15\sqrt{2}}{2}$，$y=\frac{5\sqrt{6}}{2}$，即点 $P$ 的坐标为 $(\frac{15\sqrt{2}}{2},\frac{5\sqrt{6}}{2})$．\newline（2）(1) $|QA|-|QB|=30$，则 $a=15$，$c=20$，所以 $b^2=175$，双曲线方程为 $\frac{x^2}{225}-\frac{y^2}{175}=1$；(2) $|QC|-|QD|=10$，则 $a=5$，$c=15$，所以 $b^2=200$，双曲线方程为 $\frac{y^2}{25}-\frac{x^2}{200}=1$，两双曲线方程联立，得 $Q(\frac{14400}{47},\frac{2975}{47})$，所以 $|OQ|\approx 19$ 米，$Q$ 点位置北偏东 $66^\circ$．

\paragraph{2019 年高考数学 全国II卷-文科第20题}

已知 $F_1,F_2$ 是椭圆 $C:\dfrac{x^2}{a^2}+\dfrac{y^2}{b^2}=1(a>b>0)$ 的两个焦点，$P$ 为 $C$ 上一点，$O$ 为坐标原点．
（1）若 $\triangle POF_2$ 为等边三角形，求 $C$ 的离心率；
（2）如果存在点 $P$，使得 $PF_1\perp PF_2$，且 $\triangle F_1PF_2$ 的面积等于 $16$，求 $b$ 的值和 $a$ 的取值范围．

\textbf{答案：}（1）$e=\sqrt{3}-1$；（2）$b=4$，$a$ 的取值范围为 $[4\sqrt{2},+\infty)$\par

\textbf{解答：}（1）连结 $PF_1$，由 $\triangle POF_2$ 为等边三角形可知：在 $\triangle F_1PF_2$ 中，$\angle F_1PF_2=90^\circ$，$PF_2=c$，$PF_1=\sqrt{3}c$，于是 $2a=PF_1+PF_2=c+\sqrt{3}c$，故椭圆 $C$ 的离心率为 $e=\dfrac{c}{a}=\dfrac{1}{1+\sqrt{3}}=\sqrt{3}-1$．
（2）由题意可知，满足条件的点 $P(x,y)$ 存在，当且仅当 $\dfrac{1}{2}|y|\cdot2c=16$，$\dfrac{y}{x+c}\cdot\dfrac{y}{x-c}=-1$，$\dfrac{x^2}{a^2}+\dfrac{y^2}{b^2}=1$，即 $c|y|=16$ (1)，$x^2+y^2=c^2$ (2)，$\dfrac{x^2}{a^2}+\dfrac{y^2}{b^2}=1$ (3)，由(2)(3)以及 $a^2=b^2+c^2$ 得 $y^2=\dfrac{b^4}{c^2}$，又由(1)知 $y^2=\dfrac{16^2}{c^2}$，故 $b=4$；由(2)(3)得 $x^2=\dfrac{a^2}{c^2}(c^2-b^2)$，所以 $c^2\ge b^2$，从而 $a^2=b^2+c^2\ge 2b^2=32$，故 $a\ge 4\sqrt{2}$；当 $b=4$，$a\ge 4\sqrt{2}$ 时，存在满足条件的点 $P$．所以 $b=4$，$a$ 的取值范围为 $[4\sqrt{2},+\infty)$．

\subsection*{课后作业}

\begin{enumerate}

\item 2023 年高考数学 北京卷第6题

\item 2022 年高考数学 全国乙卷-理科第5题

\item 2022 年高考数学 全国甲卷-理科第10题

\item 2019 年高考数学 全国II卷-理科第21题

\item 2021 年高考数学 上海春季卷第19题

\end{enumerate}

\noindent\textbf{作业答案索引：}

\begin{itemize}

\item 2023 年高考数学 北京卷第6题：D

\item 2022 年高考数学 全国乙卷-理科第5题：B

\item 2022 年高考数学 全国甲卷-理科第10题：A

\item 2019 年高考数学 全国II卷-理科第21题：（1）$\frac{x^2}{4}+\frac{y^2}{2}=1\ (|x|\neq 2)$，$C$ 为中心在坐标原点，焦点在 $x$ 轴上的椭圆，不含左右顶点；（2）（i）证明见解析；（ii）$\frac{16}{9}$

\item 2021 年高考数学 上海春季卷第19题：$(1)\ \frac{x^2}{100}-\frac{y^2}{300}=1,\ P(\frac{15\sqrt{2}}{2},\frac{5\sqrt{6}}{2})$；$(2)\ |OQ|\approx 19$ 米，$Q$ 点位置北偏东 $66^\circ$

\end{itemize}

\subsection*{板书设计}

一、第一定义
1. 椭圆：$|PF_1|+|PF_2|=2a,\ a>c$
2. 双曲线：$\bigl||PF_1|-|PF_2|\bigr|=2a,\ 0<a<c$
3. 抛物线：$|PF|=d(P,l)$

二、定义 $\to$ 标准方程
1. 椭圆：$\dfrac{x^2}{a^2}+\dfrac{y^2}{b^2}=1,\ c^2=a^2-b^2$
2. 双曲线：$\dfrac{x^2}{a^2}-\dfrac{y^2}{b^2}=1,\ c^2=a^2+b^2$
3. 抛物线：$y^2=2px$ 等四式

三、第二定义
$\dfrac{|PF|}{d(P,l)}=e$
$0<e<1$ 椭圆，$e=1$ 抛物线，$e>1$ 双曲线

四、第三定义的高中化表达
设 $A(-a,0),B(a,0),P(x,y)$
$k_{PA}k_{PB}=\dfrac{y^2}{x^2-a^2}$
椭圆：$k_{PA}k_{PB}=-\dfrac{b^2}{a^2}$
双曲线：$k_{PA}k_{PB}=\dfrac{b^2}{a^2}$

五、题型提醒
1. 见“和、差、等距”先判定义
2. 见“离心率”联想统一定义
3. 见“斜率积”先预判，再设点坐标

\subsection*{易错点}

\begin{itemize}

\item 把椭圆的条件误记为 $2a\ge 2c$。应明确：$2a=2c$ 时退化为线段，不是椭圆。

\item 忽略双曲线定义中的绝对值，只写成 $|PF_1|-|PF_2|=2a$，导致只覆盖一支或符号混乱。

\item 把抛物线中“点到直线的距离”误写成代数差，如对准线 $x=-1$ 直接写成 $x+1$ 而不考虑距离本义。课堂中应先确认相对位置，再写非负距离。

\item 已由定义判出曲线类型，却不会继续求参数。应训练从定义中的常数读出 $a$ 或 $c$，再由 $c^2=a^2\pm b^2$ 求 $b$。

\item 第三定义使用时只凭经验判断，不做坐标验证。正式解答必须设点、写斜率、化简成标准方程或参数关系。

\item 标准方程书写时分不清焦点在 $x$ 轴还是 $y$ 轴，导致分母位置与正负号错误。

\item 求轨迹后只写方程，不说明曲线名称或不补充定义域、特殊点去留等限制条件。

\end{itemize}

\subsection*{教师备注}

1. 本课对“第三定义”的处理尺度要稳：只能作为识别结构、压缩运算、理解离心率的补充工具，不能替代高中课标允许的方法。
2. 讲 $\text{2019-全国II卷-理科::21}$ 时，建议把“预判是椭圆”和“坐标化严格推出椭圆”分成两层，帮助学生体会观察与证明的区别。
3. 讲 $\text{2021-上海春季卷::19}$ 时，要特别提醒单位与建系：题目背景是测量，但数学本质仍是双曲线第一定义。
4. 讲抛物线题时，优先强化“$|PF|=d(P,l)$ 可双向代换”，这是高频得分点。
5. 若课堂时间紧，可把 $\text{2019-全国II卷-文科::20}$ 留作课后拓展，只在课上强调“标准方程与定义双向互通”。
6. 课后任务要求：学生完成题目后，在每题答案旁标注自己调用的是“第一定义/第二定义/第三定义预判+坐标法”，以训练方法意识。

\clearpage

\section{第4课：轨迹、隐圆与几何最值}

\noindent 建议课时：1课时。

\subsection*{教学目标}

\begin{itemize}

\item 理解并会识别两类典型“隐圆”结构：一类是条件可化为圆的判定形式，另一类是圆锥曲线背景下经代数化、配方后出现的圆方程。

\item 掌握把几何条件代数化的基本路径：设点坐标、向量表示、消元配方，并能把轨迹问题转化为标准形式 $x^2+y^2+Dx+Ey+F=0$。

\item 能在高考规范下使用坐标法、定义法、向量、平面几何处理轨迹与最值问题；知道射影几何视角只用于结构预判，不直接进入正式解答。

\item 借助题目 \texttt{2017-全国II卷-理科::20} 认识“椭圆背景下轨迹化圆”，借助题目 \texttt{2025-全国I卷::18} 体会“几何约束代数化后服务于最值”的方法链。

\end{itemize}

\subsection*{重点与难点}

\begin{itemize}

\item 从条件中提取轨迹本质：定比、定长、垂直、投影、共线、参数关系。

\item 由点的构造关系写出坐标：如利用 $\overrightarrow{NP}=\sqrt{2}\overrightarrow{NM}$、点在线段或射线上时的比例表示。

\item 二元二次式识别圆：当方程可化为 $x^2+y^2+Dx+Ey+F=0$ 时，应及时配方得到 $(x-a)^2+(y-b)^2=r^2$。

\item 最值问题中将“动点到曲线的距离”转化为“定点到轨迹/圆锥曲线的距离”或等价代数最值。

\end{itemize}

\subsection*{高观点}

从统一视角看，许多“隐圆”问题本质上是在不同几何对象之间做结构转译：圆锥曲线、投影、比例关系、垂线关系经过适当坐标化后，会显露出同类二次结构。射影几何层面可帮助学生预判：题设若由一条定轴、一个投影点、一个比例放缩构成，则往往会把原曲线上的点变成另一条二次轨迹；当代数化后出现等权的 $x^2,y^2$ 项时，就应警觉“圆”正在显现。但这一视角只用于统摄结构与概念溯源，正式解答必须回到高中课标允许的方法。

\subsection*{高中落地}

落实到高中解题语言，就是三步：第一，\textbf{设坐标}，把几何条件翻译为点的坐标关系或向量关系；第二，\textbf{化关系}，把“垂足、比例、共线、斜率、距离”写成方程；第三，\textbf{识别结果}，若化成 $x^2+y^2+Dx+Ey+F=0$，立即配方写成圆的标准方程，再用圆的圆心、半径、弦、切线或距离最值工具继续求解。

\subsection*{教学流程}

\paragraph{5分钟：导入与目标明确}

教师呈现本节主题：\textbf{轨迹、隐圆与几何最值}。回顾上节圆锥曲线中的轨迹处理经验，提出本节核心判断：题目不说“圆”，不等于结果不是圆；看到投影、比例、垂直、配方，要主动检查是否出现隐含圆结构。

\noindent\textbf{教师追问：}如果题目背景是椭圆，最后的轨迹还有可能是圆吗？\\

\textbf{预期回答：}有可能，因为经过构造、变换或消元后，轨迹类型可能改变。\par

\noindent\textbf{教师追问：}当我们得到一个二元二次方程时，第一眼应检查什么？\\

\textbf{预期回答：}检查 $x^2$ 与 $y^2$ 的系数是否相等，能否配方成圆。\par

\paragraph{8分钟：方法建构}

结合参考方法，教师归纳“隐圆”识别清单：
\[
\text{定比}\ \to\ \text{阿波罗尼斯圆};\quad
\text{直角}\ \to\ \text{直径所对圆};\quad
\text{二次式配方}\ \to\ \text{标准圆方程}.
\]
再强调本节更常见的是第二、三类：圆锥曲线背景下，经坐标化后出现圆。

\noindent\textbf{教师追问：}若轨迹满足 $\dfrac{|PA|}{|PB|}=\lambda\ (\lambda\ne 1)$，它通常是什么？\\

\textbf{预期回答：}阿波罗尼斯圆。\par

\noindent\textbf{教师追问：}若证明某点在以 $AB$ 为直径的圆上，最便捷的判定往往是什么？\\

\textbf{预期回答：}证明 $\angle APB=90^\circ$。\par

\noindent\textbf{教师追问：}若化简后得到 $x^2+y^2-2x+4y-3=0$，下一步做什么？\\

\textbf{预期回答：}配方，化为标准圆方程。\par

\paragraph{15分钟：例题一：轨迹识别与求法}

选用题目 \texttt{2017-全国II卷-理科::20} 的第（1）问进行课堂主例。设椭圆上动点为 $M(x,y)$，则
\[
\frac{x^2}{2}+y^2=1.
\]
过 $M$ 作 $x$ 轴的垂线，垂足为 $N$，故
\[
N(x,0).
\]
由
\[
\overrightarrow{NP}=\sqrt{2}\,\overrightarrow{NM}
\]
得
\[
\overrightarrow{NM}=(0,y),\qquad \overrightarrow{NP}=(0,\sqrt{2}y),
\]
从而
\[
P(x,\sqrt{2}y).
\]
设 $P(X,Y)$，则
\[
X=x,\qquad Y=\sqrt{2}y.
\]
代回椭圆方程：
\[
\frac{X^2}{2}+\frac{Y^2}{2}=1
\]
即
\[
X^2+Y^2=2.
\]
故点 $P$ 的轨迹是圆
\[
x^2+y^2=2.
\]
教师在此强调：原背景是椭圆，但经过“竖直投影 + 纵向伸缩”后，轨迹化为圆，这正是隐圆的典型来源。

\noindent\textbf{教师追问：}为什么 $N$ 的坐标是 $(x,0)$？\\

\textbf{预期回答：}因为 $N$ 是点 $M$ 到 $x$ 轴的垂足，横坐标不变、纵坐标为 $0$。\par

\noindent\textbf{教师追问：}由向量关系怎样最快得到 $P$ 的坐标？\\

\textbf{预期回答：}先写 $\overrightarrow{NM}=(0,y)$，再乘以 $\sqrt{2}$ 得 $\overrightarrow{NP}$，从而得到 $P$。\par

\noindent\textbf{教师追问：}代换时为什么要把 $P$ 的坐标记成新字母 $X,Y$？\\

\textbf{预期回答：}避免把表示椭圆上点的变量与轨迹点变量混淆。\par

\noindent\textbf{教师追问：}这个结果说明了怎样的解题意识？\\

\textbf{预期回答：}题目背景不是结论类型，关键看代数化后的结构。\par

\paragraph{6分钟：例题一提升：结构回看}

对题目 \texttt{2017-全国II卷-理科::20} 第（1）问做方法复盘。该题本质上是把椭圆上的点 $M(x,y)$ 经过映射
\[
(x,y)\mapsto (x,\sqrt{2}y)
\]
送到点 $P$。而原椭圆方程
\[
\frac{x^2}{2}+y^2=1
\]
在新变量下恰好变为
\[
x^2+y^2=2.
\]
教师强调：正式书写仍然用坐标代换，不写“仿射变换”“射影变换”等超纲术语；但教师自己在备课时要知道，这类题的深层结构是“二次曲线在特定线性关系下的像”。

\noindent\textbf{教师追问：}如果把题中 $\sqrt{2}$ 改成一般常数 $k>0$，你预判轨迹还一定是圆吗？\\

\textbf{预期回答：}不一定，通常会得到二次曲线，需要看代回后 $x^2,y^2$ 的系数是否相等。\par

\noindent\textbf{教师追问：}当什么情况下会得到圆？\\

\textbf{预期回答：}代回后 $x^2$ 与 $y^2$ 的系数相等时会得到圆。\par

\paragraph{10分钟：例题二：由几何条件到最值结构}

选用题目 \texttt{2025-全国I卷::18} 的第（2）问作方法示范，不展开全部计算细节，而突出建模链条。
先由第（1）问椭圆方程
\[
\frac{x^2}{9}+y^2=1
\]
明确研究对象。再聚焦第（2）问中条件：点 $R$ 在射线 $AP$ 上，且
\[
AR\cdot AP=3.
\]
设
\[
P(m,n),\qquad A(0,-1).
\]
因为 $R$ 在射线 $AP$ 上，可设
\[
\overrightarrow{AR}=t\overrightarrow{AP}\quad (t>0).
\]
则
\[
AR=t\,AP,
\]
由
\[
AR\cdot AP=3
\]
得
\[
t\,AP^2=3.
\]
故
\[
t=\frac{3}{m^2+(n+1)^2}.
\]
于是
\[
R=A+t\overrightarrow{AP},
\]
这就得到题中第（i）问的坐标表达。接着利用“直线 $OR$ 的斜率为直线 $OP$ 的斜率的 $3$ 倍”，建立 $m,n$ 的约束，再把点 $P$ 的轨迹化简。教师指出：这一问的关键不是盲目求最值，而是先看 $P$ 被限制在怎样的几何对象上；一旦轨迹被识别，最值就转化为点到曲线的距离最值问题。

\noindent\textbf{教师追问：}为什么这里先设 $\overrightarrow{AR}=t\overrightarrow{AP}$ 是自然的？\\

\textbf{预期回答：}因为 $R$ 在射线 $AP$ 上，方向与比例关系都可以用一个正参数表示。\par

\noindent\textbf{教师追问：}条件 $AR\cdot AP=3$ 中的两个量本质上是什么？\\

\textbf{预期回答：}是两段长度的乘积，可以借助比例参数与 $AP^2$ 关联。\par

\noindent\textbf{教师追问：}做最值前，最需要先弄清什么？\\

\textbf{预期回答：}先弄清动点 $P$ 的轨迹或满足的约束方程。\par

\paragraph{7分钟：课堂练习与即时反馈}

学生独立完成两个小任务：
\(1\) 对题目 \texttt{2017-全国II卷-理科::20} 第（1）问，口头复述从 $M$ 到 $P$ 的坐标转化链；
\(2\) 对题目 \texttt{2025-全国I卷::18} 第（2）问，写出由 $R\in \overrightarrow{AP}$ 导致的参数表示，并解释它为什么优于直接设 $R(x,y)$。
教师巡视，选择典型答案展示。

\noindent\textbf{教师追问：}若直接设 $R(x,y)$，会比设比例参数多出什么困难？\\

\textbf{预期回答：}会多出共线、同向、长度关系等多个条件，方程更杂。\par

\noindent\textbf{教师追问：}在轨迹最值问题中，什么时候要警惕变量过多导致计算失控？\\

\textbf{预期回答：}当点的构造本来就带有明显比例或向量关系时，应优先参数化。\par

\paragraph{4分钟：归纳提升与总结}

教师总结本节通法：
\[
\text{读条件}\to\text{找几何关系}\to\text{坐标/向量代数化}\to\text{配方识别轨迹}\to\text{转入最值工具}.
\]
并强调两条课堂规范：
\(1\) 看到椭圆、抛物线、双曲线，不要预设答案仍是该类曲线；
\(2\) 看到最值，不要先求导式硬算，先找轨迹、找圆心、找距离关系。

\noindent\textbf{教师追问：}今天两题共同体现的核心方法是什么？\\

\textbf{预期回答：}把几何条件代数化，再识别轨迹结构，用结构解决后续问题。\par

\noindent\textbf{教师追问：}本节中“射影视角”对我们真正有用的地方是什么？\\

\textbf{预期回答：}帮助预判结构、统一认识来源，但正式解答仍用高中方法。\par

\subsection*{课堂真题}

\paragraph{2017 年高考数学 全国II卷-理科第20题}

设 $O$ 为坐标原点，动点 $M$ 在椭圆 $C:\frac{x^2}{2}+y^2=1$ 上，过 $M$ 作 $x$ 轴的垂线，垂足为 $N$，点 $P$ 满足 $\overrightarrow{NP}=\sqrt{2}\overrightarrow{NM}$．
（1）求点 $P$ 的轨迹方程；
（2）设点 $Q$ 在直线 $x=-3$ 上，且 $\overrightarrow{OP}\cdot\overrightarrow{PQ}=1$．证明：过点 $P$ 且垂直于 $OQ$ 的直线 $l$ 过 $C$ 的左焦点 $F$．

\textbf{答案：}（1）$x^2+y^2=2$；（2）略\par

\textbf{解答：}（1）设 $M(x_0,y_0)$，由题意可得 $N(x_0,0)$，设 $P(x,y)$，由点 $P$ 满足 $\overrightarrow{NP}=\sqrt{2}\overrightarrow{NM}$．可得 $(x-x_0,y)=\sqrt{2}(0,y_0)$，可得 $x-x_0=0$，$y=\sqrt{2}y_0$，即有 $x_0=x$，$y_0=\frac{y}{\sqrt{2}}$，代入椭圆方程 $\frac{x^2}{2}+y^2=1$，可得 $\frac{x^2}{2}+\frac{y^2}{2}=1$，即有点 $P$ 的轨迹方程为圆 $x^2+y^2=2$；
（2）证明：设 $Q(-3,m)$，$P(\sqrt{2}\cos\alpha,\sin\alpha)$（$0\le\alpha<2\pi$），$\overrightarrow{OP}\cdot\overrightarrow{PQ}=1$，可得 $(\sqrt{2}\cos\alpha,\sin\alpha)\cdot(-3-\sqrt{2}\cos\alpha,m-\sin\alpha)=1$，即为 $-3\sqrt{2}\cos\alpha-2\cos^2\alpha+\sqrt{2}m\sin\alpha-2\sin^2\alpha=1$，当 $\alpha=0$ 时，上式不成立，则 $0<\alpha<2\pi$，解得 $m=\frac{3(1+\sqrt{2}\cos\alpha)}{\sqrt{2}\sin\alpha}$，即有 $Q(-3,\frac{3(1+\sqrt{2}\cos\alpha)}{\sqrt{2}\sin\alpha})$，椭圆 $\frac{x^2}{2}+y^2=1$ 的左焦点 $F(-1,0)$，由 $\overrightarrow{PF}\cdot\overrightarrow{OQ}=(-1-\sqrt{2}\cos\alpha,-\sin\alpha)\cdot(-3,\frac{3(1+\sqrt{2}\cos\alpha)}{\sqrt{2}\sin\alpha})=3+3\sqrt{2}\cos\alpha-3(1+\sqrt{2}\cos\alpha)=0$．可得过点 $P$ 且垂直于 $OQ$ 的直线 $l$ 过 $C$ 的左焦点 $F$．

\paragraph{2025 年高考数学 全国I卷第18题}

设椭圆 $C: \frac{x^2}{a^2} + \frac{y^2}{b^2} = 1 \ (a > b > 0)$ 的离心率为 $\frac{2\sqrt{2}}{3}$，下顶点为 $A$，右顶点为 $B$，$|AB| = \sqrt{10}$.

（1）求椭圆的标准方程；

（2）已知动点 $P$ 不在 $y$ 轴上，点 $R$ 在射线 $AP$ 上，且满足 $AR \cdot AP = 3$.

（i）设 $P(m,n)$，求点 $R$ 的坐标（用 $m$，$n$ 表示）；

（ii）设 $O$ 为坐标原点，$M$ 是椭圆上的动点，直线 $OR$ 的斜率为直线 $OP$ 的斜率的 $3$ 倍，求 $|PM|$ 的最大值.

\textbf{答案：}（1）$\frac{x^2}{9} + y^2 = 1$；（2）(i) $\left( \frac{3m}{m^2+(n+1)^2}, \frac{n+2 - m^2 - n^2}{m^2+(n+1)^2} \right)$；(ii) $3(\sqrt{3}+2)$\par

\textbf{解答：}（1）由题知 $A(0,-b)$，$B(a,0)$，$e = \frac{c}{a} = \frac{2\sqrt{2}}{3}$，$|AB| = \sqrt{a^2+b^2} = \sqrt{10}$，又 $c^2 = a^2 - b^2$，解得 $a^2=9$，$b^2=1$，$c^2=8$，故椭圆方程为 $\frac{x^2}{9} + y^2 = 1$.
（2）(i)设 $R(x_0,y_0)$，$A(0,-1)$，由 $A,P,R$ 共线得 $\frac{y_0+1}{x_0} = \frac{n+1}{m}$，由 $AR\cdot AP = 3$ 得 $\sqrt{x_0^2+(y_0+1)^2} \cdot \sqrt{m^2+(n+1)^2} = 3$，联立解得 $x_0 = \frac{3m}{m^2+(n+1)^2}$，$y_0 = \frac{n+2 - m^2 - n^2}{m^2+(n+1)^2}$，故 $R\left(\frac{3m}{m^2+(n+1)^2}, \frac{n+2 - m^2 - n^2}{m^2+(n+1)^2}\right)$.
(ii) $k_{OP} = \frac{n}{m}$，$k_{OR} = \frac{y_0}{x_0} = \frac{n+2 - m^2 - n^2}{3m}$，由 $k_{OR}=3k_{OP}$ 得 $\frac{n+2 - m^2 - n^2}{3m} = \frac{3n}{m}$，化简得 $m^2+n^2+8n-2=0$，即 $m^2+(n+4)^2=18$（$m \neq 0$），所以点 $P$ 在以 $N(0,-4)$ 为圆心，$3\sqrt{2}$ 为半径的圆上. 设 $M(3\cos\theta,\sin\theta)$，则 $|MN|^2 = (3\cos\theta)^2 + (\sin\theta+4)^2 = 9\cos^2\theta + \sin^2\theta + 8\sin\theta + 16 = 8\cos^2\theta + 8\sin\theta + 17 = -8\sin^2\theta + 8\sin\theta + 25 = -8\left(\sin\theta - \frac{1}{2}\right)^2 + 27 \le 27$，当且仅当 $\sin\theta = \frac{1}{2}$ 时取等，故 $|PM|_{\max} = \sqrt{27} + 3\sqrt{2} = 3(\sqrt{3}+2)$.

\subsection*{课后作业}

\begin{enumerate}

\item 2017 年高考数学 全国II卷-理科第20题

\item 2025 年高考数学 全国I卷第18题

\end{enumerate}

\noindent\textbf{作业答案索引：}

\begin{itemize}

\item 2017 年高考数学 全国II卷-理科第20题：（1）$x^2+y^2=2$；（2）略

\item 2025 年高考数学 全国I卷第18题：（1）$\frac{x^2}{9} + y^2 = 1$；（2）(i) $\left( \frac{3m}{m^2+(n+1)^2}, \frac{n+2 - m^2 - n^2}{m^2+(n+1)^2} \right)$；(ii) $3(\sqrt{3}+2)$

\end{itemize}

\subsection*{板书设计}

一、课题：轨迹、隐圆与几何最值

二、隐圆识别清单
\[
\text{定比}\to\text{阿波罗尼斯圆},\quad
\angle APB=90^\circ\to P\text{ 在以 }AB\text{ 为直径的圆上},\quad
x^2+y^2+Dx+Ey+F=0\to\text{圆}
\]

三、例1：\texttt{2017-全国II卷-理科::20}
\[
M(x,y),\ \frac{x^2}{2}+y^2=1
\]
\[
N(x,0),\quad \overrightarrow{NM}=(0,y)
\]
\[
\overrightarrow{NP}=\sqrt{2}\overrightarrow{NM}\Rightarrow P(x,\sqrt{2}y)
\]
设 $P(X,Y)$，则
\[
X=x,\ Y=\sqrt{2}y
\]
代回得
\[
\frac{X^2}{2}+\frac{Y^2}{2}=1\Rightarrow X^2+Y^2=2
\]
结论：
\[
\boxed{x^2+y^2=2}
\]

四、例2：\texttt{2025-全国I卷::18}
\[
C:\frac{x^2}{9}+y^2=1,\quad A(0,-1),\quad P(m,n)
\]
\[
R\in \overrightarrow{AP}\Rightarrow \overrightarrow{AR}=t\overrightarrow{AP}\ (t>0)
\]
\[
AR\cdot AP=3\Rightarrow t\,AP^2=3
\]
\[
AP^2=m^2+(n+1)^2
\]
\[
t=\frac{3}{m^2+(n+1)^2}
\]
\[
R=A+t\overrightarrow{AP}
\]
提示：先定轨迹，再谈最值

五、方法小结
\[
\text{几何条件}\xrightarrow{\text{坐标化}}\text{方程}\xrightarrow{\text{配方/识别}}\text{圆或其他轨迹}\xrightarrow{\text{距离关系}}\text{最值}
\]

\subsection*{易错点}

\begin{itemize}

\item 把题目背景中的曲线类型直接当作结果类型，例如见到椭圆就默认轨迹仍是椭圆，忽视了中间构造对坐标的改变。

\item 由 $\overrightarrow{NP}=\sqrt{2}\overrightarrow{NM}$ 只放大长度，忘记方向相同，导致点 $P$ 坐标写错。

\item 代换时混用同一组字母表示不同点，出现“既表示 $M$ 又表示 $P$”的变量混乱。

\item 看到最值立即进入复杂代数运算，没有先判断动点轨迹，导致计算量失控。

\item 对“点在射线 $AP$ 上”只写共线，不写同向与比例关系，错失最简参数化路径。

\item 把教师补充的统一视角误当正式答题语言，在书面解答中使用超纲术语。

\end{itemize}

\subsection*{教师备注}

1. 本节定位是一轮复习中的“方法模块课”，不是单纯讲题课。两题都服务于同一个主线：\textbf{几何条件代数化后识别隐含轨迹，再转入几何性质或最值}。

2. 题目 \texttt{2017-全国II卷-理科::20} 课堂上重点完成第（1）问，并可用第（2）问口头点拨“轨迹变圆后，许多向量垂直关系会转化为圆与极值或焦点性质的联动”，但不宜在本节全部展开，以免挤压方法总结时间。

3. 题目 \texttt{2025-全国I卷::18} 不建议课堂上把第（2）问算完，而应把重心放在两步：一是从 $R\in \overrightarrow{AP}$ 与 $AR\cdot AP=3$ 提炼出参数表示；二是强调最值问题先找轨迹约束。完整求解可留作课后分层任务。

4. 分层要求：基础学生至少要能独立完成“设点坐标 $\to$ 写向量/比例关系 $\to$ 代入原方程”；中档学生要能自主配方识别圆；拔高学生要能从轨迹进一步处理距离最值。

5. 课堂语言要反复压实三句：
\[
\text{先翻译条件\text{，}再判断结构\text{；}先确定轨迹\text{，}再处理最值\text{；}正式解答只用高中方法\text{。}}
\]

6. 课后任务建议：
第1层：独立规范完成 \texttt{2017-全国II卷-理科::20} 第（1）问与第（2）问的思路整理；
第2层：完成 \texttt{2025-全国I卷::18} 第（2）问的全解，并在旁批中注明每一步使用的是“坐标法、向量、距离公式、斜率关系”中的哪一种。

7. 评价关注点不是最终是否一遍算通，而是学生能否主动说出：
\[
\text{“这个条件适合参数化”},\quad \text{“这个二次式该配方”},\quad \text{“这个最值应先找轨迹”}.
\]

\clearpage

\section{第5课：直线与圆锥曲线：联立、韦达与点差}

\noindent 建议课时：2课时。

\subsection*{教学目标}

\begin{itemize}

\item 理解直线与圆锥曲线相交问题中的三条主线：联立方程得到交点参数关系、用判别式判断位置关系、用韦达定理进行整体代换。

\item 能在椭圆、双曲线、抛物线背景下，熟练完成如下转化：设直线为 $y=kx+m$ 或参数式后，将交点坐标满足的一元二次方程写出，并提取 $x_1+x_2,\ x_1x_2$ 或 $y_1+y_2,\ y_1y_2$。

\item 掌握弦中点问题的点差法：由两点代入曲线方程后作差，建立弦斜率与中点坐标的关系，并能回到高中课标方法写出标准解答。

\item 能识别何时适合使用“超级韦达定理”思想：当研究定点与两交点连线斜率的和、积、垂直、点积等关系时，先构造关于斜率的二次方程，再用韦达定理处理。

\item 形成方法选择意识：选择题优先用结构判断与必要条件，解答题先建系、再联立、后整理，避免无序展开。

\end{itemize}

\subsection*{重点与难点}

\begin{itemize}

\item 联立直线与圆锥曲线方程后，将问题归结为一元二次方程。

\item 判别式 $\Delta$ 的意义：$\Delta>0$ 表示相交于两点，$\Delta=0$ 表示相切，$\Delta<0$ 表示无实交点。

\item 韦达整体代换：把交点信息压缩为 $x_1+x_2,\ x_1x_2$ 或对应的 $y$ 形式，减少直接求根。

\item 点差法的核心公式来自“代入后作差”，本质是把端点信息转为中点与斜率信息。

\item 含定点连线斜率的题型，可借助“超级韦达定理”视角进行结构预判，但正式书写仍回到坐标法与韦达定理。

\end{itemize}

\subsection*{高观点}

从统一视角看，直线与二次曲线相交时，许多结论都表现为“一个一次对象切入一个二次对象，于是产生二次代数关系”。这解释了为什么判别式、韦达定理、点差法会反复出现：它们都在处理同一个二次结构。进一步地，弦中点、极点极线、过定点的弦束等现象，在更高层次上都可看作射影几何中“直线束与二次曲线”的对应关系。但本课中射影几何只用于帮助学生预判结构，例如预判“中点条件最终应变成一次关系”“两条连线斜率最终应满足二次方程”，正式求解一律回到高中允许的方法。

\subsection*{高中落地}

把统一视角落回高中方法，可概括为三句话：第一，见交点先联立；第二，见两交点整体关系先想韦达；第三，见弦中点先想点差。若题目研究的是过定点的两条连线斜率，就把曲线方程改写后转成关于 $k$ 的一元二次方程，再用韦达定理；若题目研究的是弦所在直线或中点轨迹，就对两端点代入曲线方程后作差，再把 $x_1+x_2=2x_0,\ y_1+y_2=2y_0$ 代入。所有步骤都可在坐标法框架内完成。

\subsection*{教学流程}

\paragraph{8分钟：导入与目标定位}

回顾直线与圆锥曲线问题的常见设问：位置关系、弦长、中点、过定点、斜率和积、轨迹。教师板书本课主线：\[\text{联立}\to\text{判别式}\to\text{韦达}\to\text{点差}.\] 说明本课两课时分别解决“整体代换”与“中点弦”。

\noindent\textbf{教师追问：}直线与圆锥曲线相交后，为什么很多题最后都会落到一元二次方程？\\

\textbf{预期回答：}因为把直线方程代入二次曲线方程后，通常会消去一个变量，得到关于另一个变量的一元二次方程。\par

\noindent\textbf{教师追问：}如果题目只关心两交点的和、积、中点，而不要求具体坐标，最该优先想到什么？\\

\textbf{预期回答：}优先用韦达定理，不一定直接求根。\par

\paragraph{12分钟：方法建模一：联立与判别式}

以一般步骤梳理：设直线 $y=kx+m$，与二次曲线联立，得到关于 $x$ 的方程\[Ax^2+Bx+C=0.\] 若交点为 $A(x_1,y_1),B(x_2,y_2)$，则\[x_1+x_2=-\frac{B}{A},\qquad x_1x_2=\frac{C}{A}.\] 再由 $y_i=kx_i+m$ 得到对应的 $y$ 关系。强调：只要题目中出现“中点”“弦长”“斜率和积”“过定点”，都应先把交点整体信息抽出来。

\noindent\textbf{教师追问：}判别式在这类题中除了判断相交，还常常承担什么作用？\\

\textbf{预期回答：}它还能给参数范围，或者保证后续韦达定理可用时有两个实根。\par

\noindent\textbf{教师追问：}为什么很多时候只写 $x_1+x_2,\ x_1x_2$ 就够了？\\

\textbf{预期回答：}因为很多几何量如中点、弦长平方、对称式都能用和与积表示。\par

\paragraph{15分钟：例1选择题快解}

使用题目 ID：\texttt{2023-全国乙卷-理科::11}。目标：训练“中点可行性”的结构判断。设双曲线 $x^2-\dfrac{y^2}{9}=1$ 上两点为 $A(x_1,y_1),B(x_2,y_2)$，中点为 $M(x_0,y_0)$。用点差法预判：若存在斜率存在的弦，则有\[k_{AB}=\frac{9x_0}{y_0}.\] 再结合中点弦的存在性与双曲线区域特征判断选项。课堂不重在穷举，而重在识别：并非任意点都可作某条弦的中点。必要时回到联立设过 $M$ 的弦为 $y-y_0=k(x-x_0)$，代入双曲线，由韦达判断能否得到关于两端点的实根。

\noindent\textbf{教师追问：}双曲线上的两点中点为什么不一定仍落在双曲线上？\\

\textbf{预期回答：}因为双曲线不是凸集，中点的位置受弦方向和分支位置限制。\par

\noindent\textbf{教师追问：}如果一个点是某条弦的中点，最直接的验证策略是什么？\\

\textbf{预期回答：}过该点设弦方程，与双曲线联立，再看是否能得到两个实交点并满足该点是中点。\par

\paragraph{10分钟：小结与第一课时当堂训练}

归纳第一课时结论：\[\text{位置关系看 }\Delta,\qquad \text{整体关系看韦达}.\] 安排口头训练：若直线与椭圆、双曲线、抛物线相交于两点，哪些量优先用和与积表示？学生回答后，教师补充中点坐标公式\[\left(\frac{x_1+x_2}{2},\frac{y_1+y_2}{2}\right).\]

\noindent\textbf{教师追问：}弦长是否也能借助韦达处理？\\

\textbf{预期回答：}能，常先求某一坐标差的平方，再用 $(x_1-x_2)^2=(x_1+x_2)^2-4x_1x_2$。\par

\noindent\textbf{教师追问：}中点问题为什么天然适合整体代换？\\

\textbf{预期回答：}因为中点本身就是两个端点坐标和的一半。\par

\paragraph{8分钟：第二课时导入：从中点到弦}

回顾技巧 \texttt{49}。板书点差法模板：设端点 $A(x_1,y_1),B(x_2,y_2)$，中点 $M(x_0,y_0)$，则\[x_1+x_2=2x_0,\qquad y_1+y_2=2y_0.\] 将 $A,B$ 分别代入曲线方程后作差，利用平方差公式，提取\[\frac{y_1-y_2}{x_1-x_2}=k_{AB}.\]

\noindent\textbf{教师追问：}点差法与直接联立相比，最大的优势是什么？\\

\textbf{预期回答：}它不必先把直线参数完全设出并求根，能直接建立中点与斜率的联系。\par

\noindent\textbf{教师追问：}为什么它特别适合处理中点弦问题？\\

\textbf{预期回答：}因为作差后自然出现 $x_1+x_2,\ y_1+y_2$，正好可以换成中点坐标。\par

\paragraph{15分钟：例2点差法专题}

使用题目 ID：\texttt{2016-全国III卷-文科::20}，聚焦第（II）问的思路训练。抛物线 $C:y^2=2x$ 中，若 $AB$ 为一条弦，中点为 $M(x_0,y_0)$，由点差法可直接得到\[k_{AB}=\frac{1}{y_0}.\] 若题设再给出与平行于 $x$ 轴的两条直线、面积倍数有关的条件，就把这些条件转化为 $A,B$ 的纵坐标关系，最终归结为 $M$ 的坐标关系。教师强调：正式求解时，仍然应从题设中的平行关系和面积关系出发，用坐标法逐步推出中点轨迹；点差法用于降低结构复杂度。

\noindent\textbf{教师追问：}对抛物线 $y^2=2px$，为什么中点弦斜率公式是 $k_{AB}=\dfrac{p}{y_0}$？\\

\textbf{预期回答：}因为两点代入后作差得到 $y_1+y_2=p\dfrac{x_1-x_2}{y_1-y_2}$ 的等价关系，代入 $y_1+y_2=2y_0$ 后可化成该公式。\par

\noindent\textbf{教师追问：}已知中点后，弦方程应怎样写？\\

\textbf{预期回答：}写成 \(y-y_0=k_{AB}(x-x_0)\)。\par

\paragraph{17分钟：例3超级韦达定理视角}

使用题目 ID：\texttt{2017-全国I卷-理科::20} 第（2）问。以点 $P_2(0,1)$ 为定点，直线 $l$ 与椭圆相交于 $A,B$。若研究 $P_2A$ 与 $P_2B$ 的斜率 $k_1,k_2$，则可把椭圆方程改写为关于 $(x-0),(y-1)$ 的形式，并结合不过 $P_2$ 的动直线，得到关于 $k$ 的二次方程，其两根为 $k_1,k_2$。由条件\[k_1+k_2=-1\] 用韦达定理得到动直线参数满足的线性关系，进而证明 $l$ 过定点。教师强调：课堂可用“超级韦达定理”命名帮助记忆，但板书规范写法必须是“构造关于斜率的一元二次方程，再用韦达定理”。

\noindent\textbf{教师追问：}为什么这里不直接设 $A,B$ 再硬算，而要转向研究 $P_2A,P_2B$ 的斜率？\\

\textbf{预期回答：}因为题目条件直接给的是两条连线斜率的和，研究斜率比研究点坐标更贴近条件。\par

\noindent\textbf{教师追问：}所谓‘超级韦达’的本质是什么？\\

\textbf{预期回答：}本质仍是韦达定理，只是未知量从交点坐标换成了连线斜率。\par

\paragraph{10分钟：变式提升与方法分流}

快速对接两个题目 ID：\texttt{2024-天津卷::18} 与 \texttt{2018-全国III卷-理科::16}。前者可用连线斜率或向量点积转化，体现“定点 $T$ 与两交点 $P,Q$”背景下的韦达整体代换；后者可把过焦点弦参数化后联立抛物线，再用\[\overrightarrow{MA}\cdot\overrightarrow{MB}=0\] 或斜率关系处理。教师只做方法分流，不展开整题全解：\[\text{定点+两交点连线}\Rightarrow\text{斜率二次方程};\qquad \text{中点/弦}\Rightarrow\text{点差法}.\]

\noindent\textbf{教师追问：}看到向量点积恒成立，你更愿意先保留端点坐标，还是先压缩成和与积？\\

\textbf{预期回答：}先压缩成和与积，更容易消去端点坐标。\par

\noindent\textbf{教师追问：}看到过焦点的弦，你会先联立还是先用几何性质？\\

\textbf{预期回答：}一般先联立，因为焦点条件常能把直线参数设得更简洁。\par

\paragraph{5分钟：全课总结}

教师带学生完成方法树：\[\begin{array}{c}
\text{直线与圆锥曲线}\\
\downarrow\\
\text{联立}\\
\downarrow\\
\begin{cases}
\Delta & \text{判位置}\\
\text{韦达} & \text{判整体关系}\\
\text{点差} & \text{判中点弦}
\end{cases}
\end{array}\] 强调规范：选择最短路径，但书写必须完整、变量定义清楚、结论回扣题意。

\noindent\textbf{教师追问：}这节课三种工具各自最典型的触发词是什么？\\

\textbf{预期回答：}判别式对应‘相交、相切、参数范围’，韦达对应‘两交点和积、中点、弦长’，点差对应‘中点弦、弦斜率、轨迹’。\par

\noindent\textbf{教师追问：}正式解答中，哪些步骤最容易丢分？\\

\textbf{预期回答：}设元不清、漏写联立方程、韦达使用对象混乱、把结构观察当成证明。\par

\subsection*{课堂真题}

\paragraph{2023 年高考数学 全国乙卷-理科第11题}

设 $A$，$B$ 为双曲线 $x^2-\frac{y^2}{9}=1$ 上两点，下列四个点中，可为线段 $AB$ 中点的是（ ）

\textbf{答案：}$(-1,-4)$\par

\textbf{解答：}设 $A(x_1,y_1),B(x_2,y_2)$，则 $AB$ 的中点 $M(\frac{x_1+x_2}{2},\frac{y_1+y_2}{2})$，可得 $k_{AB}=\frac{y_1-y_2}{x_1-x_2}$，$k=\frac{\frac{y_1+y_2}{2}}{\frac{x_1+x_2}{2}}=\frac{y_1+y_2}{x_1+x_2}$，因为 $A,B$ 在双曲线上，则 $\begin{cases} x_1^2-\frac{y_1^2}{9}=1 \\ x_2^2-\frac{y_2^2}{9}=1 \end{cases}$，两式相减得 $(x_1^2-x_2^2)-\frac{y_1^2-y_2^2}{9}=0$，所以 $k_{AB}\cdot k=\frac{y_1^2-y_2^2}{x_1^2-x_2^2}=9$. 对于选项A：可得 $k=1$，$k_{AB}=9$，则 $AB:y=9x-8$，联立方程 $\begin{cases} y=9x-8 \\ x^2-\frac{y^2}{9}=1 \end{cases}$，消去 $y$ 得 $72x^2-2\times72x+73=0$，此时 $\Delta=(-2\times72)^2-4\times72\times73=-288<0$，所以直线 $AB$ 与双曲线没有交点，故A错误；对于选项B：可得 $k=-2$，$k_{AB}=-\frac{9}{2}$，则 $AB:y=-\frac{9}{2}x-\frac{5}{2}$，联立方程 $\begin{cases} y=-\frac{9}{2}x-\frac{5}{2} \\ x^2-\frac{y^2}{9}=1 \end{cases}$，消去 $y$ 得 $45x^2+2\times45x+61=0$，此时 $\Delta=(2\times45)^2-4\times45\times61=-4\times45\times16<0$，所以直线 $AB$ 与双曲线没有交点，故B错误；对于选项C：可得 $k=3$，$k_{AB}=3$，则 $AB:y=3x$，由双曲线方程可得 $a=1,b=3$，则 $AB:y=3x$ 为双曲线的渐近线，所以直线 $AB$ 与双曲线没有交点，故C错误；对于选项D：$k=4$，$k_{AB}=\frac{9}{4}$，则 $AB:y=\frac{9}{4}x-\frac{7}{4}$，联立方程 $\begin{cases} y=\frac{9}{4}x-\frac{7}{4} \\ x^2-\frac{y^2}{9}=1 \end{cases}$，消去 $y$ 得 $63x^2+126x-193=0$，此时 $\Delta=126^2+4\times63\times193>0$，故直线 $AB$ 与双曲线有两个交点，故D正确.

\paragraph{2016 年高考数学 全国III卷-文科第20题}

已知抛物线 $C: y^2=2x$ 的焦点为 $F$, 平行于 $x$ 轴的两条直线 $l_1$, $l_2$ 分别交 $C$ 于 $A$, $B$ 两点, 交 $C$ 的准线于 $P$, $Q$ 两点.
(I) 若 $F$ 在线段 $AB$ 上, $R$ 是 $PQ$ 的中点, 证明 $AR\parallel FQ$;
(II) 若 $\triangle PQF$ 的面积是 $\triangle ABF$ 的面积的两倍, 求 $AB$ 中点的轨迹方程.

\textbf{答案：}见解析\par

\textbf{解答：}(I) 连接 $RF$, $PF$. 由 $AP=AF$, $BQ=BF$ 及 $AP\parallel BQ$, 得 $\angle AFP+\angle BFQ=90^{\circ}$, 所以 $\angle PFQ=90^{\circ}$. 因为 $R$ 是 $PQ$ 的中点, 所以 $RF=RP=RQ$, 所以 $\triangle PAR\cong\triangle FAR$, 所以 $\angle PAR=\angle FAR$, $\angle PRA=\angle FRA$. 又 $\angle BQF+\angle BFQ=180^{\circ}-\angle QBF=\angle PAF=2\angle PAR$, 所以 $\angle FQB=\angle PAR$, 所以 $\angle PRA=\angle PQF$, 所以 $AR\parallel FQ$.
(II) 设 $A(x_1,y_1)$, $B(x_2,y_2)$, $F(\frac{1}{2},0)$, 准线为 $x=-\frac{1}{2}$. $S_{\triangle PQF}=\frac{1}{2}|PQ|=\frac{1}{2}|y_1-y_2|$. 设直线 $AB$ 与 $x$ 轴交点为 $N$, 则 $S_{\triangle ABF}=\frac{1}{2}|FN||y_1-y_2|$. 因为 $\triangle PQF$ 的面积是 $\triangle ABF$ 的面积的两倍, 所以 $2|FN|=1$, 得 $x_N=1$, 即 $N(1,0)$. 设 $AB$ 中点为 $M(x,y)$, 由 $\begin{cases}y_1^2=2x_1\\ y_2^2=2x_2\end{cases}$ 得 $(y_1-y_2)(y_1+y_2)=2(x_1-x_2)$, 又 $\dfrac{y_1-y_2}{x_1-x_2}=\dfrac{y}{x-1}$, 所以 $y\cdot2y=2$, 即 $y^2=x-1$. 所以 $AB$ 中点轨迹方程为 $y^2=x-1$.

\paragraph{2017 年高考数学 全国I卷-理科第20题}

已知椭圆 $C:\frac{x^2}{a^2}+\frac{y^2}{b^2}=1$（$a>b>0$），四点 $P_1(1,1)$，$P_2(0,1)$，$P_3(-1,\frac{\sqrt{3}}{2})$，$P_4(1,\frac{\sqrt{3}}{2})$ 中恰有三点在椭圆 $C$ 上．
（1）求 $C$ 的方程；
（2）设直线 $l$ 不经过 $P_2$ 点且与 $C$ 相交于 $A$，$B$ 两点．若直线 $P_2A$ 与直线 $P_2B$ 的斜率的和为 $-1$，证明：$l$ 过定点．

\textbf{答案：}（1）$\frac{x^2}{4}+y^2=1$；（2）证明见解析，定点为 $(2,-1)$\par

\textbf{解答：}解：（1）根据椭圆的对称性，$P_3(-1,\frac{\sqrt{3}}{2})$，$P_4(1,\frac{\sqrt{3}}{2})$ 两点必在椭圆 $C$ 上，又 $P_4$ 的横坐标为 1，\ensuremath{\therefore}椭圆必不过 $P_1(1,1)$，\ensuremath{\therefore}$P_2(0,1)$，$P_3(-1,\frac{\sqrt{3}}{2})$，$P_4(1,\frac{\sqrt{3}}{2})$ 三点在椭圆 $C$ 上．把 $P_2(0,1)$，$P_3(-1,\frac{\sqrt{3}}{2})$ 代入椭圆 $C$，得：$\begin{cases} \frac{0}{a^2}+\frac{1}{b^2}=1 \\ \frac{1}{a^2}+\frac{3}{4b^2}=1 \end{cases}$，解得 $a^2=4$，$b^2=1$，\ensuremath{\therefore}椭圆 $C$ 的方程为 $\frac{x^2}{4}+y^2=1$．
证明：（2）(1)当斜率不存在时，设 $l:x=m$，$A(m,y_A)$，$B(m,-y_A)$，\ensuremath{\because}直线 $P_2A$ 与直线 $P_2B$ 的斜率的和为 $-1$，\ensuremath{\therefore}$\frac{y_A-1}{m-0}+\frac{-y_A-1}{m-0}=\frac{-2}{m}=-1$，解得 $m=2$，此时 $l$ 过椭圆右顶点，不存在两个交点，故不满足．(2)当斜率存在时，设 $l:y=kx+t$，($t\neq1$)，$A(x_1,y_1)$，$B(x_2,y_2)$，联立 $\begin{cases} y=kx+t \\ \frac{x^2}{4}+y^2=1 \end{cases}$，整理，得 $(1+4k^2)x^2+8ktx+4t^2-4=0$，$\Delta>0$，$x_1+x_2=-\frac{8kt}{1+4k^2}$，$x_1x_2=\frac{4t^2-4}{1+4k^2}$，则 $k_{P_2A}+k_{P_2B}=\frac{y_1-1}{x_1}+\frac{y_2-1}{x_2}=\frac{kx_1+t-1}{x_1}+\frac{kx_2+t-1}{x_2}=2k+\frac{t-1}{x_1}+\frac{t-1}{x_2}=2k+(t-1)\frac{x_1+x_2}{x_1x_2}=2k+(t-1)\frac{-\frac{8kt}{1+4k^2}}{\frac{4t^2-4}{1+4k^2}}=2k-\frac{2k(t-1)}{t+1}=\frac{2k(t+1)-2k(t-1)}{t+1}=\frac{4k}{t+1}=-1$，又 $t\neq1$，\ensuremath{\therefore}$t=-2k-1$，此时 $\Delta=64k^2-16(1+4k^2)(t^2-1)$，代入得 $\Delta=-64k$，存在 $k$，使得 $\Delta>0$ 成立，\ensuremath{\therefore}直线 $l$ 的方程为 $y=kx-2k-1$，当 $x=2$ 时，$y=-1$，\ensuremath{\therefore}$l$ 过定点 $(2,-1)$．

\paragraph{2024 年高考数学 天津卷第18题}

已知椭圆 $\frac{x^2}{a^2}+\frac{y^2}{b^2}=1(a>b>0)$ 的离心率 $e=\frac{1}{2}$. 左顶点为 $A$, 下顶点为 $B$, $C$ 是线段 $OB$ 的中点, 其中 $S_{\triangle ABC}=\frac{3\sqrt{3}}{2}$.
(1) 求椭圆方程.
(2) 过点 $(0,-\frac{3}{2})$ 的动直线与椭圆有两个交点 $P$, $Q$. 在 $y$ 轴上是否存在点 $T$ 使得 $\overrightarrow{TP}\cdot\overrightarrow{TQ}\leq 0$ 恒成立. 若存在求出这个 $T$ 点纵坐标的取值范围, 若不存在请说明理由.

\textbf{答案：}(1) $\frac{x^2}{12}+\frac{y^2}{9}=1$; (2) 存在 $T(0,t)$ 且 $-3\leq t\leq \frac{3}{2}$, 使得 $\overrightarrow{TP}\cdot\overrightarrow{TQ}\leq 0$ 恒成立.\par

\textbf{解答：}(1) 因为 $e=\frac{1}{2}$, 故 $a=2c$, $b=\sqrt{3}c$. 所以 $A(-2c,0)$, $B(0,-\sqrt{3}c)$, $C(0,-\frac{\sqrt{3}}{2}c)$. $S_{\triangle ABC}=\frac{1}{2}\times 2c\times \frac{\sqrt{3}}{2}c=\frac{\sqrt{3}}{2}c^2=\frac{3\sqrt{3}}{2}$, 解得 $c=\sqrt{3}$, 所以 $a=2\sqrt{3}$, $b=3$, 故椭圆方程为 $\frac{x^2}{12}+\frac{y^2}{9}=1$.
(2) 设过点 $(0,-\frac{3}{2})$ 的动直线方程为 $y=kx-\frac{3}{2}$, $P(x_1,y_1)$, $Q(x_2,y_2)$, $T(0,t)$. 联立 $\begin{cases} 3x^2+4y^2=36 \\ y=kx-\frac{3}{2} \end{cases}$, 得 $(3+4k^2)x^2-12kx-27=0$, $\Delta>0$, $x_1+x_2=\frac{12k}{3+4k^2}$, $x_1x_2=-\frac{27}{3+4k^2}$. $\overrightarrow{TP}=(x_1,y_1-t)$, $\overrightarrow{TQ}=(x_2,y_2-t)$. $\overrightarrow{TP}\cdot\overrightarrow{TQ}=x_1x_2+(y_1-t)(y_2-t)=x_1x_2+(kx_1-\frac{3}{2}-t)(kx_2-\frac{3}{2}-t)=(1+k^2)x_1x_2-k(\frac{3}{2}+t)(x_1+x_2)+(\frac{3}{2}+t)^2$. 代入韦达定理, 得 $\frac{[(3+2t)^2-12t-45]k^2+3(\frac{3}{2}+t)^2-27}{3+4k^2}$. 要使 $\overrightarrow{TP}\cdot\overrightarrow{TQ}\leq 0$ 恒成立, 需 $\begin{cases} (3+2t)^2-12t-45\leq 0 \\ 3(\frac{3}{2}+t)^2-27\leq 0 \end{cases}$, 解得 $-3\leq t\leq \frac{3}{2}$. 当斜率不存在时, 直线为 $x=0$, 与椭圆交于 $(0,3)$, $(0,-3)$, 此时需 $-3\leq t\leq 3$, 结合得 $-3\leq t\leq \frac{3}{2}$. 综上, 存在 $T(0,t)$ 且 $-3\leq t\leq \frac{3}{2}$, 使得 $\overrightarrow{TP}\cdot\overrightarrow{TQ}\leq 0$ 恒成立.

\paragraph{2018 年高考数学 全国III卷-理科第16题}

已知点$M(-1,1)$和抛物线$C:y^2=4x$，过$C$的焦点且斜率为$k$的直线与$C$交于$A$，$B$两点．若$\angle AMB=90^\circ$，则$k=$\underline{\hspace{3cm}}．

\textbf{答案：}$2$\par

\textbf{解答：}焦点$F(1,0)$，设直线$AB:y=k(x-1)$，联立$\begin{cases}y=k(x-1)\\y^2=4x\end{cases}$得$k^2x^2-2(2+k^2)x+k^2=0$，设$A(x_1,y_1)$，$B(x_2,y_2)$，则$x_1+x_2=\frac{2(2+k^2)}{k^2}$，$x_1x_2=1$，$y_1+y_2=k(x_1+x_2-2)=\frac{4}{k}$，$y_1y_2=k^2(x_1-1)(x_2-1)=k^2[x_1x_2-(x_1+x_2)+1]=-4$，$\because \angle AMB=90^\circ$，$\overrightarrow{MA}\cdot\overrightarrow{MB}=0$，即$(x_1+1)(x_2+1)+(y_1-1)(y_2-1)=0$，整理得$x_1x_2+(x_1+x_2)+y_1y_2-(y_1+y_2)+2=0$，代入得$1+\frac{2(2+k^2)}{k^2}-4-\frac{4}{k}+2=0$，解得$k=2$．

\subsection*{课后作业}

\begin{enumerate}

\item 2023 年高考数学 全国乙卷-理科第11题

\item 2016 年高考数学 全国III卷-文科第20题

\item 2017 年高考数学 全国I卷-理科第20题

\item 2024 年高考数学 天津卷第18题

\item 2018 年高考数学 全国III卷-理科第16题

\end{enumerate}

\noindent\textbf{作业答案索引：}

\begin{itemize}

\item 2023 年高考数学 全国乙卷-理科第11题：$(-1,-4)$

\item 2016 年高考数学 全国III卷-文科第20题：见解析

\item 2017 年高考数学 全国I卷-理科第20题：（1）$\frac{x^2}{4}+y^2=1$；（2）证明见解析，定点为 $(2,-1)$

\item 2024 年高考数学 天津卷第18题：(1) $\frac{x^2}{12}+\frac{y^2}{9}=1$; (2) 存在 $T(0,t)$ 且 $-3\leq t\leq \frac{3}{2}$, 使得 $\overrightarrow{TP}\cdot\overrightarrow{TQ}\leq 0$ 恒成立.

\item 2018 年高考数学 全国III卷-理科第16题：$2$

\end{itemize}

\subsection*{板书设计}

一、主线框架\[\text{联立}\to\text{判别式}\to\text{韦达}\to\text{点差}\]
二、联立模板\[\begin{cases}
\text{直线}: y=kx+m\\
\text{曲线}: \text{二次方程}
\end{cases}\Rightarrow Ax^2+Bx+C=0\]\[x_1+x_2=-\frac{B}{A},\qquad x_1x_2=\frac{C}{A}\]
三、判别式\[\Delta=B^2-4AC\]\[\Delta>0:\text{相交于两点},\quad \Delta=0:\text{相切},\quad \Delta<0:\text{无实交点}\]
四、点差法模板\[x_1+x_2=2x_0,\qquad y_1+y_2=2y_0\]\[F(x_1,y_1)=0,\ F(x_2,y_2)=0\Rightarrow F(x_1,y_1)-F(x_2,y_2)=0\]\[\frac{y_1-y_2}{x_1-x_2}=k_{AB}\]
五、常用中点弦公式
椭圆 \[\frac{x^2}{a^2}+\frac{y^2}{b^2}=1:\quad k_{AB}=-\frac{b^2x_0}{a^2y_0}\]
双曲线 \[\frac{x^2}{a^2}-\frac{y^2}{b^2}=1:\quad k_{AB}=\frac{b^2x_0}{a^2y_0}\]
抛物线 \[y^2=2px:\quad k_{AB}=\frac{p}{y_0}\]
六、斜率二次方程视角\[k_{PA},k_{PB}\text{ 是某一元二次方程的两根}\]\[\text{用韦达： }k_{PA}+k_{PB},\ k_{PA}k_{PB}\]

\subsection*{易错点}

\begin{itemize}

\item 把“联立后得到的一元二次方程的根”误认为一定是两交点的纵坐标。应先明确消去的是哪个变量，根对应哪个坐标。

\item 只会用判别式判断有无交点，却不会把 $\Delta\ge 0$ 继续化为参数范围。

\item 已知两交点时直接求根，忽视了用 \(x_1+x_2\)、\(x_1x_2\) 整体代换可以显著简化计算。

\item 点差法中忘记作差后还要代入 \(x_1+x_2=2x_0\)、\(y_1+y_2=2y_0\)，导致无法与中点建立联系。

\item 把点差法得到的关系当作结论直接套用，却忽略适用前提，例如弦斜率存在、分母不为零。

\item 将‘超级韦达定理’当成新定理记忆。应明确它只是构造关于斜率的二次方程后使用韦达定理的技巧，不是脱离坐标法的独立方法。

\item 在正式解答中，只写观察到的结构，不写联立、化简、由韦达得出的具体等式，造成过程分丢失。

\end{itemize}

\subsection*{教师备注}

1. 两课时节奏建议：第一课时重“联立+判别式+韦达”，第二课时重“点差法+斜率二次方程”。
2. 对基础一般的班级，\texttt{2023-全国乙卷-理科::11} 只讲必要条件与排除法，不必做过深推广；对基础较强的班级，可补充“某点可作弦中点的充要条件”探究。
3. 讲授 \texttt{2016-全国III卷-文科::20} 时，不宜直接把中点弦公式端给学生，应先让学生亲自完成一次‘代入后作差’，再抽象成模板。
4. 讲授 \texttt{2017-全国I卷-理科::20} 时，要反复强调：题目条件给的是 \(P_2A\) 与 \(P_2B\) 的斜率之和，因此未知量应改设为斜率，而不是执着于端点坐标。
5. 对 \texttt{2024-天津卷::18}，课堂上不必完整展开全部推导，可作为方法统摄题，突出从 \(\overrightarrow{TP}\cdot\overrightarrow{TQ}\le 0\) 到和积关系的转化意识。
6. 课后任务布置建议：
   (1) 基础巩固：完成 \texttt{2023-全国乙卷-理科::11}，要求写出‘为什么另外三个点不可能’；
   (2) 方法专练：完成 \texttt{2018-全国III卷-理科::16}，要求写出联立方程与垂直条件转化；
   (3) 综合提升：二选一完成 \texttt{2016-全国III卷-文科::20} 或 \texttt{2017-全国I卷-理科::20}；
   (4) 拓展反思：阅读 \texttt{2024-天津卷::18}，标出其中哪些步骤适合使用韦达整体代换。
7. 批改关注点：是否明确设元；是否说明根的对应意义；是否规范写出韦达关系；是否把结构观察落实为可评分的代数步骤。

\clearpage

\section{第6课：焦点三角形：面积与离心率}

\noindent 建议课时：2课时。

\subsection*{教学目标}

\begin{itemize}

\item 会识别圆锥曲线中的焦点三角形模型，并能优先调用定义与三角形面积公式处理问题：椭圆用 $|PF_1|+|PF_2|=2a$，双曲线用 $\bigl||PF_1|-|PF_2|\bigr|=2a$，抛物线关注过焦点弦与顶点构成的三角形。

\item 会在焦点三角形中综合运用余弦定理、正弦定理、面积公式 $S=\dfrac12 ab\sin C$，实现边角互化，并能据此求面积、参数与离心率。

\item 掌握两组核心结论并能在适当条件下直接调用：椭圆中 $S_{\triangle PF_1F_2}=b^2\tan\dfrac\theta2$，双曲线中 $S_{\triangle PF_1F_2}=\dfrac{b^2}{\tan\dfrac\theta2}$，其中 $\theta=\angle F_1PF_2$。

\item 会把离心率转化为焦点三角形中的边角关系：椭圆中 $e=\dfrac{|F_1F_2|}{|PF_1|+|PF_2|}$，双曲线中 $e=\dfrac{|F_1F_2|}{\bigl||PF_1|-|PF_2|\bigr|}$，并结合正弦定理化简。

\item 理解射影几何视角仅用于统一认识“焦点、弦、角、交比不变量诱导的结构预判”，正式求解仍回到高中课标允许的坐标法、定义法、韦达定理、点差法、向量或平面几何。

\end{itemize}

\subsection*{重点与难点}

\begin{itemize}

\item 焦点三角形的第一切入口不是坐标硬算，而是定义 $+$ 面积公式 $+$ 边角关系。

\item 已知 $\angle F_1PF_2=\theta$ 时，优先联想椭圆、双曲线对应的面积表达式。

\item 已知 $PF_1\perp PF_2$ 时，立即转化为 $\theta=90^\circ$，从而简化面积与参数关系。

\item 求离心率时，优先在 $\triangle PF_1F_2$ 中使用正弦定理，把 $|PF_1|,|PF_2|$ 用角表示，再代回 $e$ 的定义式。

\end{itemize}

\subsection*{高观点}

从统一视角看，椭圆、双曲线、抛物线都可放在“焦点-点-弦”这一类二次曲线结构中理解：题目若给出焦点、过焦点弦、顶角、切线等条件，往往预示应先寻找三角形中的边角不变量，而不是立即设点坐标展开。这个视角只用于结构预判与概念溯源；课堂中的正式高考解答一律落回到高中方法，如圆锥曲线定义、余弦定理、正弦定理、面积公式、坐标运算等。

\subsection*{高中落地}

高中可执行的处理流程是：先辨认曲线类型，再写定义；若有角，先用 $S=\dfrac12|PF_1|\,|PF_2|\sin\theta$；若有和差关系，用定义把 $|PF_1|,|PF_2|$ 联系起来；若要求离心率，用 $e=\dfrac ca$，并借助 $|F_1F_2|=2c$ 与定义中的 $2a$ 进行转化；若条件复杂，再用余弦定理或正弦定理补上边角关系。

\subsection*{教学流程}

\paragraph{6分钟：导入与目标揭示}

呈现本课主题：焦点三角形中的“面积”与“离心率”。教师明确两条主线：一是由定义与面积公式求面积，二是由边角关系求离心率。快速回顾：椭圆 $\dfrac{x^2}{a^2}+\dfrac{y^2}{b^2}=1$，双曲线 $\dfrac{x^2}{a^2}-\dfrac{y^2}{b^2}=1$，其中 $c^2=a^2\pm b^2$，离心率分别为 $e=\dfrac ca$。

\noindent\textbf{教师追问：}看到“焦点、曲线上一点、三角形面积”这类信息时，第一反应应写什么？\\

\textbf{预期回答：}先写圆锥曲线定义，再写三角形面积公式 $S=\dfrac12ab\sin C$。\par

\noindent\textbf{教师追问：}为什么这类题通常不宜一开始就设 $P(x,y)$ 全面硬算？\\

\textbf{预期回答：}因为定义直接给出边的和或差，面积又直接依赖两边与夹角，结构比坐标展开更短。\par

\paragraph{12分钟：知识建构一：焦点三角形面积公式}

系统整理三类模型。椭圆中，若 $\theta=\angle F_1PF_2$，则
$|PF_1|+|PF_2|=2a$，
$S_{\triangle PF_1F_2}=\dfrac12|PF_1|\,|PF_2|\sin\theta$，
并可得到 $|PF_1|\,|PF_2|=\dfrac{2b^2}{1+\cos\theta}$，故
$S_{\triangle PF_1F_2}=b^2\tan\dfrac\theta2$。
双曲线中，
$\bigl||PF_1|-|PF_2|\bigr|=2a$，
$|PF_1|\,|PF_2|=\dfrac{2b^2}{1-\cos\theta}$，故
$S_{\triangle PF_1F_2}=\dfrac{b^2}{\tan\dfrac\theta2}$。抛物线中仅作方法提示：过焦点弦与顶点构成三角形时，优先用焦点弦性质。

\noindent\textbf{教师追问：}椭圆与双曲线在面积公式推导中，本质差异出现在哪一步？\\

\textbf{预期回答：}差异在定义：一个是和为 $2a$，一个是差的绝对值为 $2a$，所以后续 $|PF_1|\,|PF_2|$ 的表达式分母分别变成 $1+\cos\theta$ 与 $1-\cos\theta$。\par

\noindent\textbf{教师追问：}若已知 $PF_1\perp PF_2$，对椭圆面积公式有何立刻结论？对双曲线又如何？\\

\textbf{预期回答：}此时 $\theta=90^\circ$。椭圆有 $S=b^2$；双曲线也有 $S=b^2$。\par

\paragraph{10分钟：例题训练一}

使用题目 ID：`2020-全国I卷-文科::11`。教学目标：在双曲线中由点到原点距离切入，借坐标关系得到顶角，再回到面积公式。建议解题路径：由题设双曲线 $x^2-\dfrac{y^2}{3}=1$ 得 $a=1,b^2=3,c=2$。设 $P(x,y)$，由 $|OP|=2$ 得 $x^2+y^2=4$；联立双曲线方程可求出与 $x^2-y^2$ 相关的关系，再转化到 $|PF_1|,|PF_2|$ 或直接用坐标面积公式。总结时强调：本题虽可坐标解，但结构上仍是焦点三角形面积题。

\noindent\textbf{教师追问：}本题先用定义还是先用坐标关系更顺？为什么？\\

\textbf{预期回答：}先用坐标关系更顺，因为直接给了原点距离，容易得到 $x^2+y^2$。\par

\noindent\textbf{教师追问：}已知 $|OP|$ 时，哪一个高中方法最自然？\\

\textbf{预期回答：}设 $P(x,y)$，联立方程求点的位置关系，再求面积。\par

\paragraph{12分钟：例题训练二}

使用题目 ID：`2023-全国甲卷-理科::12`。教学目标：在椭圆中已知 $\cos\angle F_1PF_2$，优先调用焦点三角形积与面积结构，再求 $|OP|$。教师带学生先识别参数：椭圆 $\dfrac{x^2}{9}+\dfrac{y^2}{6}=1$ 中 $a^2=9,b^2=6,c^2=3$。由 $\cos\theta=\dfrac35$，得 $\sin\theta=\dfrac45$，再由
$|PF_1|\,|PF_2|=\dfrac{2b^2}{1+\cos\theta}=\dfrac{12}{1+3/5}=\dfrac{15}{2}$。
之后可结合余弦定理
$(2c)^2=|PF_1|^2+|PF_2|^2-2|PF_1|\,|PF_2|\cos\theta$
求出 $|PF_1|^2+|PF_2|^2$，再由椭圆中
$|PF_1|^2+|PF_2|^2=2(|OP|^2+c^2)$
求 $|OP|$。

\noindent\textbf{教师追问：}本题已知的是角的余弦，为什么不先设点坐标？\\

\textbf{预期回答：}因为角直接对应焦点三角形边角关系，先求 $|PF_1|\,|PF_2|$ 更快。\par

\noindent\textbf{教师追问：}式子 $|PF_1|^2+|PF_2|^2=2(|OP|^2+c^2)$ 是怎么来的？\\

\textbf{预期回答：}可由焦点坐标设为 $(\pm c,0)$，展开两段距离平方后相加得到。\par

\paragraph{5分钟：小结与当堂检测}

归纳第一课时结论：
$\boxed{\text{面积题先定义、再面积、后边角}}$。
设置口答检测：若椭圆中 $PF_1\perp PF_2$，则 $S=b^2$；若双曲线中 $PF_1\perp PF_2$，也有 $S=b^2$。提醒学生：结论相同，但来源不同，不可混淆。

\noindent\textbf{教师追问：}焦点三角形面积题的三步法是什么？\\

\textbf{预期回答：}先辨曲线并写定义，再写面积公式，再用余弦定理或坐标补充边角关系。\par

\noindent\textbf{教师追问：}为什么“直角”条件在这类题中尤其关键？\\

\textbf{预期回答：}因为它会把 $\theta$ 固定为 $90^\circ$，很多公式立刻大幅简化。\par

\paragraph{8分钟：复习承接}

第二课时开头复盘：面积问题与离心率问题并不分离，二者都依赖焦点三角形中的边角关系。教师板书离心率转化：椭圆中
$e=\dfrac{c}{a}=\dfrac{|F_1F_2|}{|PF_1|+|PF_2|}$；
双曲线中
$e=\dfrac{c}{a}=\dfrac{|F_1F_2|}{\bigl||PF_1|-|PF_2|\bigr|}$。

\noindent\textbf{教师追问：}为什么离心率问题也可以放进焦点三角形来做？\\

\textbf{预期回答：}因为 $2c=|F_1F_2|$，而 $2a$ 又能由定义转成焦点三角形两边的和或差。\par

\noindent\textbf{教师追问：}这里的关键转化对象是什么？\\

\textbf{预期回答：}把 $e=\dfrac ca$ 转成关于 $|F_1F_2|,|PF_1|,|PF_2|$ 的式子。\par

\paragraph{12分钟：知识建构二：离心率角度公式}

在 $\triangle PF_1F_2$ 中设 $\angle PF_1F_2=\alpha,\ \angle PF_2F_1=\beta$，则顶角为 $\pi-(\alpha+\beta)$。由正弦定理
$\dfrac{|PF_1|}{\sin\beta}=\dfrac{|PF_2|}{\sin\alpha}=\dfrac{|F_1F_2|}{\sin(\alpha+\beta)}$。
于是椭圆中
$e=\dfrac{\sin(\alpha+\beta)}{\sin\alpha+\sin\beta}=\dfrac{\cos\dfrac{\alpha+\beta}{2}}{\cos\dfrac{\alpha-\beta}{2}}$；
双曲线中
$e=\left|\dfrac{\sin(\alpha+\beta)}{\sin\beta-\sin\alpha}\right|=\left|\dfrac{\sin\dfrac{\alpha+\beta}{2}}{\sin\dfrac{\alpha-\beta}{2}}\right|$。
强调：这些公式不是死记对象，而是由正弦定理即时推导的结果。

\noindent\textbf{教师追问：}为什么顶角写成 $\pi-(\alpha+\beta)$ 后，正弦定理仍能出现 $\sin(\alpha+\beta)$？\\

\textbf{预期回答：}因为 $\sin[\pi-(\alpha+\beta)]=\sin(\alpha+\beta)$。\par

\noindent\textbf{教师追问：}这组公式的本质工具是什么？\\

\textbf{预期回答：}本质是正弦定理加上圆锥曲线定义。\par

\paragraph{13分钟：例题训练三}

使用题目 ID：`2019-全国II卷-文科::20`。本题分两层使用。
第一层聚焦第（1）问：由 $\triangle POF_2$ 为等边三角形，把几何条件转化为焦点三角形中的角，再用正弦定理与椭圆定义求离心率。
第二层聚焦第（2）问：由 $PF_1\perp PF_2$ 得 $\theta=90^\circ$，故面积 $S_{\triangle F_1PF_2}=b^2=16$，得 $b=4$；再讨论存在点 $P$ 的条件，回到椭圆方程及几何位置约束求 $a$ 的范围。教师强调：这里“存在性”不能只凭公式，必须回到椭圆上点的实际可取性判断。

\noindent\textbf{教师追问：}第（1）问中，等边三角形提供的核心信息是边，还是角？\\

\textbf{预期回答：}本质上提供了角与边的等量关系，最终要转化到焦点三角形的角关系。\par

\noindent\textbf{教师追问：}第（2）问中，为什么得到 $b=4$ 后还不能立刻结束？\\

\textbf{预期回答：}因为题目还要求存在点 $P$，所以还要检验参数下图形上确有这样的点。\par

\paragraph{10分钟：例题训练四}

使用题目 ID：`2020-全国III卷-理科::11`。教学目标：在双曲线中把“离心率 $+$ 直角 $+$ 面积”迅速联成参数关系。由 $e=\sqrt5$ 得 $c=\sqrt5 a$，结合 $c^2=a^2+b^2$，得 $b^2=4a^2$。又因 $F_1P\perp F_2P$，所以 $S_{\triangle PF_1F_2}=b^2$。由面积为 $4$ 得 $b^2=4$，进而 $a=1$。本题用来训练学生的“见直角，想 $S=b^2$；见离心率，想 $c^2=a^2+b^2$”的快速联想。

\noindent\textbf{教师追问：}本题最短路径中的第一个关键跳步是什么？\\

\textbf{预期回答：}先由 $PF_1\perp F_2P$ 直接得到面积 $S=b^2$。\par

\noindent\textbf{教师追问：}若学生坚持设点坐标，这样做是否可以？课堂应怎样评价？\\

\textbf{预期回答：}可以，但不经济。应评价为方法合法，但结构意识不足。\par

\paragraph{8分钟：拓展辨析}

使用题目 ID：`2022-全国乙卷-理科::11` 作为辨析材料，不展开完整计算，只作思路点拨：题中有“以实轴为直径的圆”“过焦点作切线”“两支交点”“角的余弦”这些信息，预示应先把切线条件转化为几何角关系，再回到双曲线焦点三角形中的离心率表达。教师提醒：题面结构复杂时，更要先识别核心三角形，而不是被辅助圆牵着走。

\noindent\textbf{教师追问：}这道题表面上多了圆与切线，本质上仍在考什么？\\

\textbf{预期回答：}本质上仍在考双曲线焦点三角形中的角关系与离心率。\par

\noindent\textbf{教师追问：}面对复杂图形，如何避免陷入局部关系的碎算？\\

\textbf{预期回答：}先找主三角形，再把附加条件往主三角形的边角关系上翻译。\par

\paragraph{7分钟：归纳提升与作业布置}

完成整课归纳：
第一，面积题优先顺序为“定义 $\to$ 面积 $\to$ 边角”；
第二，离心率题优先顺序为“$e=\dfrac ca$ $\to$ 焦点边的和差 $\to$ 正弦定理”；
第三，若题目给出直角、等边、切线、圆、原点距离等条件，要先判断它们对应的是哪类边角信息。最后布置分层作业与整理任务。

\noindent\textbf{教师追问：}今天所有题型背后最统一的中枢是什么？\\

\textbf{预期回答：}都是把条件翻译到焦点三角形中，再用定义和边角公式求解。\par

\noindent\textbf{教师追问：}正式考试中，哪些方法是允许且推荐的？\\

\textbf{预期回答：}坐标法、定义法、韦达定理、点差法、向量或平面几何。\par

\subsection*{课堂真题}

\paragraph{2020 年高考数学 全国I卷-文科第11题}

设 $F_1, F_2$ 是双曲线 $C: x^2 - \frac{y^2}{3} = 1$ 的两个焦点，$O$ 为坐标原点，点 $P$ 在 $C$ 上且 $|OP| = 2$, 则 $\triangle PF_1F_2$ 的面积为

\textbf{答案：}$3$\par

\textbf{解答：}由已知，$F_1(-2,0), F_2(2,0)$, 则 $a=1, c=2$, 因为 $|OP|=2 = \frac{1}{2}|F_1F_2|$, 所以点 $P$ 在以 $F_1F_2$ 为直径的圆上，即 $\triangle PF_1F_2$ 是以 $P$ 为直角顶点的直角三角形，故 $|PF_1|^2 + |PF_2|^2 = |F_1F_2|^2 = 16$, 又 $|PF_1|-|PF_2|=2a=2$, 所以 $4 = (|PF_1|-|PF_2|)^2 = |PF_1|^2 + |PF_2|^2 - 2|PF_1||PF_2| = 16 - 2|PF_1||PF_2|$, 解得 $|PF_1||PF_2|=6$, 所以 $S_{\triangle PF_1F_2} = \frac{1}{2}|PF_1||PF_2|=3$.

\paragraph{2023 年高考数学 全国甲卷-理科第12题}

设 $O$ 为坐标原点，$F_1$，$F_2$ 为椭圆 $C: \frac{x^2}{9} + \frac{y^2}{6} = 1$ 的两个焦点，点 $P$ 在 $C$ 上，$\cos\angle F_1PF_2 = \frac{3}{5}$，则 $|OP| = $（         ）

\textbf{答案：}B\par

\textbf{解答：}方法一：设 $\angle F_1PF_2 = 2\theta$，$0<\theta<\frac{\pi}{2}$，所以 $S_{\triangle PF_1F_2} = b^2\tan\frac{\angle F_1PF_2}{2} = b^2\tan\theta$，由 $\cos\angle F_1PF_2 = \cos2\theta = \frac{\cos^2\theta-\sin^2\theta}{\cos^2\theta+\sin^2\theta} = \frac{1-\tan^2\theta}{1+\tan^2\theta} = \frac{3}{5}$，解得：$\tan\theta = \frac{1}{2}$，由椭圆方程可知，$a^2=9$，$b^2=6$，$c^2=a^2-b^2=3$，所以 $S_{\triangle PF_1F_2} = \frac{1}{2}\times F_1F_2\times|y_p| = \frac{1}{2}\times2\sqrt{3}\times|y_p| = 6\times\frac{1}{2}$，解得：$y_p^2=3$，则 $x_p^2 = 9\times\left(1-\frac{3}{6}\right) = \frac{9}{2}$，因此 $|OP| = \sqrt{x_p^2+y_p^2} = \sqrt{3+\frac{9}{2}} = \frac{\sqrt{30}}{2}$．故选：B．
方法二：因为 $|PF_1|+|PF_2|=2a=6$ (1)，$|PF_1|^2+|PF_2|^2-2|PF_1||PF_2|\cos\angle F_1PF_2 = |F_1F_2|^2$，即 $|PF_1|^2+|PF_2|^2 - \frac{6}{5}|PF_1||PF_2| = 12$ (2)，联立(1)(2)，解得：$|PF_1||PF_2| = \frac{15}{2}$，$|PF_1|^2+|PF_2|^2 = 21$，而 $\overrightarrow{PO} = \frac{1}{2}(\overrightarrow{PF_1}+\overrightarrow{PF_2})$，所以 $|\overrightarrow{OP}| = |\overrightarrow{PO}| = \frac{1}{2}\sqrt{|\overrightarrow{PF_1}|^2+|\overrightarrow{PF_2}|^2+2\overrightarrow{PF_1}\cdot\overrightarrow{PF_2}} = \frac{1}{2}\sqrt{21+2\times\frac{3}{5}\times\frac{15}{2}} = \frac{\sqrt{30}}{2}$．故选：B．
方法三：因为 $|PF_1|+|PF_2|=2a=6$ (1)，$|PF_1|^2+|PF_2|^2-2|PF_1||PF_2|\cos\angle F_1PF_2 = |F_1F_2|^2$，即 $|PF_1|^2+|PF_2|^2 - \frac{6}{5}|PF_1||PF_2| = 12$ (2)，联立(1)(2)，解得：$|PF_1|^2+|PF_2|^2 = 21$，由中线定理可知，$2(|OP|^2+|OF_1|^2) = |PF_1|^2+|PF_2|^2 = 42$，易知 $|F_1F_2| = 2\sqrt{3}$，解得：$|OP| = \frac{\sqrt{30}}{2}$．故选：B．

\paragraph{2019 年高考数学 全国II卷-文科第20题}

已知 $F_1,F_2$ 是椭圆 $C:\dfrac{x^2}{a^2}+\dfrac{y^2}{b^2}=1(a>b>0)$ 的两个焦点，$P$ 为 $C$ 上一点，$O$ 为坐标原点．
（1）若 $\triangle POF_2$ 为等边三角形，求 $C$ 的离心率；
（2）如果存在点 $P$，使得 $PF_1\perp PF_2$，且 $\triangle F_1PF_2$ 的面积等于 $16$，求 $b$ 的值和 $a$ 的取值范围．

\textbf{答案：}（1）$e=\sqrt{3}-1$；（2）$b=4$，$a$ 的取值范围为 $[4\sqrt{2},+\infty)$\par

\textbf{解答：}（1）连结 $PF_1$，由 $\triangle POF_2$ 为等边三角形可知：在 $\triangle F_1PF_2$ 中，$\angle F_1PF_2=90^\circ$，$PF_2=c$，$PF_1=\sqrt{3}c$，于是 $2a=PF_1+PF_2=c+\sqrt{3}c$，故椭圆 $C$ 的离心率为 $e=\dfrac{c}{a}=\dfrac{1}{1+\sqrt{3}}=\sqrt{3}-1$．
（2）由题意可知，满足条件的点 $P(x,y)$ 存在，当且仅当 $\dfrac{1}{2}|y|\cdot2c=16$，$\dfrac{y}{x+c}\cdot\dfrac{y}{x-c}=-1$，$\dfrac{x^2}{a^2}+\dfrac{y^2}{b^2}=1$，即 $c|y|=16$ (1)，$x^2+y^2=c^2$ (2)，$\dfrac{x^2}{a^2}+\dfrac{y^2}{b^2}=1$ (3)，由(2)(3)以及 $a^2=b^2+c^2$ 得 $y^2=\dfrac{b^4}{c^2}$，又由(1)知 $y^2=\dfrac{16^2}{c^2}$，故 $b=4$；由(2)(3)得 $x^2=\dfrac{a^2}{c^2}(c^2-b^2)$，所以 $c^2\ge b^2$，从而 $a^2=b^2+c^2\ge 2b^2=32$，故 $a\ge 4\sqrt{2}$；当 $b=4$，$a\ge 4\sqrt{2}$ 时，存在满足条件的点 $P$．所以 $b=4$，$a$ 的取值范围为 $[4\sqrt{2},+\infty)$．

\paragraph{2020 年高考数学 全国III卷-理科第11题}

设双曲线 $C:\dfrac{x^2}{a^2}-\dfrac{y^2}{b^2}=1\ (a>0,b>0)$ 的左、右焦点分别为 $F_1,F_2$，离心率为 $\sqrt{5}$。$P$ 是 $C$ 上一点，且 $F_1P\perp F_2P$。若 $\triangle PF_1F_2$ 的面积为 $4$，则 $a=$

\textbf{答案：}$A$\par

\textbf{解答：}因为 $e=\dfrac{c}{a}=\sqrt{5}$，所以 $c=\sqrt{5}a$。根据双曲线定义 $|PF_1|-|PF_2|=2a$。$\triangle PF_1F_2$ 面积 $\dfrac12|PF_1|\cdot|PF_2|=4$，得 $|PF_1|\cdot|PF_2|=8$。又 $F_1P\perp F_2P$，所以 $|PF_1|^2+|PF_2|^2=(2c)^2=20a^2$。而 $(|PF_1|-|PF_2|)^2=4a^2$，即 $|PF_1|^2+|PF_2|^2-2|PF_1|\cdot|PF_2|=4a^2$，所以 $20a^2-16=4a^2$，得 $16a^2=16$，$a=1$。故选 $A$。

\paragraph{2022 年高考数学 全国乙卷-理科第11题}

双曲线$C$的两个焦点为$F_1, F_2$，以$C$的实轴为直径的圆记为$D$，过$F_1$作$D$的切线与$C$的两支交于$M, N$两点，且$\cos\angle F_1NF_2 = \frac{3}{5}$，则$C$的离心率为（\qquad）

\textbf{答案：}C\par

\textbf{解答：}解：依题意不妨设双曲线焦点在$x$轴，设过$F_1$作圆$D$的切线切点为$G$，\\所以$OG \perp NF_1$，因为$\cos\angle F_1NF_2 = \frac{3}{5} > 0$，所以$N$在双曲线的右支，\\所以$OG = a$，$OF_1 = c$，$GF_1 = b$，设$\angle F_1NF_2 = \alpha$，$\angle F_2F_1N = \beta$，\\由$\cos\angle F_1NF_2 = \frac{3}{5}$，即$\cos\alpha = \frac{3}{5}$，则$\sin\alpha = \frac{4}{5}$，$\sin\beta = \frac{a}{c}$，$\cos\beta = \frac{b}{c}$，\\在$\triangle F_2F_1N$中，$\sin\angle F_1F_2N = \sin(\pi-\alpha-\beta) = \sin(\alpha+\beta)$\\$= \sin\alpha\cos\beta + \cos\alpha\sin\beta = \frac{4}{5}\cdot\frac{b}{c} + \frac{3}{5}\cdot\frac{a}{c} = \frac{3a+4b}{5c}$，\\由正弦定理得$\frac{2c}{\sin\alpha} = \frac{NF_2}{\sin\beta} = \frac{NF_1}{\sin\angle F_1F_2N} = \frac{5c}{2}$，\\所以$NF_1 = \frac{5c}{2}\sin\angle F_1F_2N = \frac{5c}{2}\cdot\frac{3a+4b}{5c} = \frac{3a+4b}{2}$，$NF_2 = \frac{5c}{2}\sin\beta = \frac{5c}{2}\cdot\frac{a}{c} = \frac{5a}{2}$\\又$NF_1 - NF_2 = \frac{3a+4b}{2} - \frac{5a}{2} = \frac{4b-2a}{2} = 2a$，\\所以$2b = 3a$，即$\frac{b}{a} = \frac{3}{2}$，\\所以双曲线的离心率$e = \frac{c}{a} = \sqrt{1+\frac{b^2}{a^2}} = \frac{\sqrt{13}}{2}$.

\subsection*{课后作业}

\begin{enumerate}

\item 2019 年高考数学 全国II卷-文科第20题

\item 2022 年高考数学 全国乙卷-理科第11题

\item 2023 年高考数学 全国甲卷-理科第12题

\end{enumerate}

\noindent\textbf{作业答案索引：}

\begin{itemize}

\item 2019 年高考数学 全国II卷-文科第20题：（1）$e=\sqrt{3}-1$；（2）$b=4$，$a$ 的取值范围为 $[4\sqrt{2},+\infty)$

\item 2022 年高考数学 全国乙卷-理科第11题：C

\item 2023 年高考数学 全国甲卷-理科第12题：B

\end{itemize}

\subsection*{板书设计}

一、课题：焦点三角形：面积与离心率

二、核心定义
1. 椭圆：$|PF_1|+|PF_2|=2a,\ c^2=a^2-b^2,\ e=\dfrac ca$
2. 双曲线：$\bigl||PF_1|-|PF_2|\bigr|=2a,\ c^2=a^2+b^2,\ e=\dfrac ca$

三、面积主公式
$S_{\triangle PF_1F_2}=\dfrac12|PF_1|\,|PF_2|\sin\theta$
其中 $\theta=\angle F_1PF_2$

四、结论整理
1. 椭圆：$|PF_1|\,|PF_2|=\dfrac{2b^2}{1+\cos\theta}$
   $S=b^2\tan\dfrac\theta2$
2. 双曲线：$|PF_1|\,|PF_2|=\dfrac{2b^2}{1-\cos\theta}$
   $S=\dfrac{b^2}{\tan\dfrac\theta2}$
3. 若 $PF_1\perp PF_2$，则 $\theta=90^\circ$，从而 $S=b^2$

五、离心率转化
1. 椭圆：$e=\dfrac{|F_1F_2|}{|PF_1|+|PF_2|}$
2. 双曲线：$e=\dfrac{|F_1F_2|}{\bigl||PF_1|-|PF_2|\bigr|}$

六、正弦定理模板
设底角为 $\alpha,\beta$，则
$\dfrac{|PF_1|}{\sin\beta}=\dfrac{|PF_2|}{\sin\alpha}=\dfrac{|F_1F_2|}{\sin(\alpha+\beta)}$

七、方法流程
面积题：定义 $\to$ 面积 $\to$ 余弦/坐标
离心率题：$e=\dfrac ca$ $\to$ 焦点边和差 $\to$ 正弦定理

\subsection*{易错点}

\begin{itemize}

\item 把椭圆与双曲线的定义混写，尤其把双曲线误写成 $|PF_1|+|PF_2|=2a$。纠正方式：反复口述“椭圆看和，双曲线看差”。

\item 见到焦点三角形就机械套 $S=b^2\tan\dfrac\theta2$ 或 $S=\dfrac{b^2}{\tan\dfrac\theta2}$，却没有先确认曲线类型。纠正方式：任何公式前先标注“椭/双”。

\item 将 $PF_1\perp PF_2$ 只理解为“能用勾股”，没有意识到它直接给出 $\theta=90^\circ$，从而面积可直接化简为 $b^2$。纠正方式：训练条件反射式识别。

\item 求离心率时只盯住 $e=\dfrac ca$，却不会把 $2c$ 与 $|F_1F_2|$、把 $2a$ 与定义中的和差联系。纠正方式：要求学生每次先写“$2c=|F_1F_2|$”。

\item 复杂图形中被圆、切线、原点、等边等附加元素牵走，没有回到主三角形。纠正方式：先问“哪一个三角形真正连着焦点与所求参数？”

\item 正式解答中过多使用超纲叙述。纠正方式：允许用统一视角做预判，但书写时必须落在坐标法、定义法、余弦定理、正弦定理、向量或平面几何上。

\end{itemize}

\subsection*{教师备注}

1. 本模块适合一轮复习中段，前提是学生已经完成椭圆、双曲线定义与基本量 $a,b,c,e$ 的系统复习。
2. 第一课时重点放在“面积题的统一识别与简算”，第二课时重点放在“离心率的边角化”。
3. 对 `2020-全国I卷-文科::11` 与 `2023-全国甲卷-理科::12`，课堂要有意识对比两种入口：坐标入口与结构入口，帮助学生建立方法选择意识。
4. 对 `2019-全国II卷-文科::20`，建议只完整推进其中一问，另一问做关键步骤点拨，否则容易挤压归纳时间。
5. 对学有余力学生，可要求他们自行推导椭圆、双曲线中 $|PF_1|\,|PF_2|$ 的表达式；对基础薄弱学生，则要求先能熟练使用 $PF_1\perp PF_2\Rightarrow S=b^2$ 这一高频结论。
6. 课堂评价标准应突出“是否先抓住定义与主三角形”，而不只看计算是否做出来。
7. 射影视角只用于教师备课时统一结构与预判题型，不进入学生正式答题表述。

\clearpage

\section{第7课：焦点弦、焦比弦与焦半径}

\noindent 建议课时：2课时。

\subsection*{教学目标}

\begin{itemize}

\item 理解并会辨认焦点弦、焦半径、焦比弦三类对象，能从题干中的“过焦点”“斜率或倾斜角”“长度比”“通径”“准线”迅速识别解题入口。

\item 掌握抛物线焦点弦的核心结论：对 $y^2=2px\,(p>0)$，若过焦点的弦与抛物线交于 $M,N$，倾斜角为 $\alpha$，则 $|MN|=\dfrac{2p}{\sin^2\alpha}$，并能结合 $k=\tan\alpha$ 完成转化。

\item 掌握椭圆 $\dfrac{x^2}{a^2}+\dfrac{y^2}{b^2}=1\,(a>b>0)$ 中过右焦点 $F(c,0)$ 的焦点弦结论：若倾斜角为 $\alpha$，则 $|AF|=\dfrac{b^2}{a+c\cos\alpha}$，$|BF|=\dfrac{b^2}{a-c\cos\alpha}$，并能据此求弦长、比值与离心率。

\item 会把焦比弦信息统一写成 $\lambda=\dfrac{|AF|}{|BF|}$，并在高中允许的方法框架内，通过定义、坐标法、韦达定理、点差法、向量或平面几何完成求解。

\item 形成“先结构判断，后规范求解”的意识：射影视角仅用于识别共线、分比、交比不变所提示的结构，不直接作为高考书写解答。

\end{itemize}

\subsection*{重点与难点}

\begin{itemize}

\item 焦点弦的识别与分类：椭圆、抛物线中的焦点弦长度公式及适用条件。

\item 由斜率到倾斜角的转化：$k=\tan\alpha$，$\sin^2\alpha=\dfrac{k^2}{1+k^2}$，$\cos^2\alpha=\dfrac{1}{1+k^2}$。

\item 焦半径、焦比弦与离心率之间的数量联系，特别是 $\lambda=\dfrac{|AF|}{|BF|}$ 的统一表达。

\item 过焦点且垂直于坐标轴的特殊弦，即通径：抛物线通径长为 $2p$，椭圆通径长为 $\dfrac{2b^2}{a}$。

\item 正式解答时的回归路径：公式结论若已知可直接调用；若需论证，则回到联立方程、根与系数关系、定义法与几何关系。

\end{itemize}

\subsection*{高观点}

从统一视角看，许多“焦点弦—长度比—方向角”问题，本质上都在研究一条过定点的截线与二次曲线的两个交点所形成的分比结构。射影视角只帮助学生预判：当题干出现“过焦点”“两交点”“长度比”时，往往应先寻找分比或参数化入口；但高考作答必须回到高中方法，如设过焦点直线方程、联立曲线方程、用韦达定理得到 $x_1+x_2,\,x_1x_2$ 或 $y_1+y_2,\,y_1y_2$，再求弦长与比值。

\subsection*{高中落地}

可执行的高中化策略为：先设过焦点直线为 $y=k(x-c)$ 或 $x=my+\dfrac p2$ 等标准形式；再与曲线联立得到关于一个变量的二次方程；用韦达定理提取两交点的和与积；最后借助 $|AB|=\sqrt{1+k^2}\,|x_1-x_2|$ 或焦点定义 $|PF|=e\cdot d(P,l)$ 完成求值。若题目给出长度比，则统一设 $\lambda=\dfrac{|AF|}{|BF|}$，先列比例关系，再与定义或弦长公式配合收尾。

\subsection*{教学流程}

\paragraph{8分钟：导入与目标建立}

教师呈现本讲核心对象：焦点弦、焦半径、焦比弦。引导学生回忆：抛物线中“过焦点的弦”为什么比一般弦更值得优先处理；椭圆中“过焦点且垂直于长轴”的特殊弦是什么。明确本课两条主线：一是长度公式，二是比例关系。

\noindent\textbf{教师追问：}什么叫焦点弦？\\

\textbf{预期回答：}过焦点并与圆锥曲线交于两点的弦叫焦点弦。\par

\noindent\textbf{教师追问：}抛物线 $y^2=2px$ 的焦点坐标是什么？通径长是多少？\\

\textbf{预期回答：}焦点为 $\left(\dfrac p2,0\right)$，通径长为 $2p$。\par

\noindent\textbf{教师追问：}如果题目中出现“过焦点”“交于两点”“长度比”，你首先应想到哪类结构？\\

\textbf{预期回答：}应优先想到焦点弦、焦半径或焦比弦结构，而不是立刻硬算两个交点坐标。\par

\paragraph{12分钟：知识梳理一：一般弦长到焦点弦长}

板书一般弦长公式：若直线 $y=kx+m$ 与曲线交于 $A(x_1,y_1),B(x_2,y_2)$，则 $|AB|=\sqrt{1+k^2}\,|x_1-x_2|$。进一步强调：若该弦还过焦点，则可使用特殊公式。对抛物线 $y^2=2px$，过焦点弦倾斜角为 $\alpha$ 时，$|MN|=\dfrac{2p}{\sin^2\alpha}$。

\noindent\textbf{教师追问：}一般弦长公式的本质来自什么？\\

\textbf{预期回答：}来自两点间距离公式与直线斜率关系。\par

\noindent\textbf{教师追问：}若已知斜率 $k$，怎样把它变成公式中的 $\alpha$？\\

\textbf{预期回答：}由 $k=\tan\alpha$，再用 $\sin^2\alpha=\dfrac{k^2}{1+k^2}$ 或 $\cos^2\alpha=\dfrac{1}{1+k^2}$。\par

\noindent\textbf{教师追问：}为什么焦点弦问题通常不必先求出两个交点的完整坐标？\\

\textbf{预期回答：}因为过焦点带来额外结构，常能直接得到弦长或焦半径关系。\par

\paragraph{20分钟：例题研讨一}

使用题目 ID：$\text{2017-全国I卷-理科::10}$。目标：通过“互相垂直的两条过焦点直线”建立两个焦点弦长度之和的表达式，求最小值。教师不改写题干，只组织学生按题目原意分析。解题主线：设 $l_1,l_2$ 的倾斜角分别为 $\alpha,\,\alpha+\dfrac\pi2$。由抛物线 $y^2=4x$，可知 $p=2$，故 $|AB|=\dfrac{4}{\sin^2\alpha}$，$|DE|=\dfrac{4}{\cos^2\alpha}$。于是
$|AB|+|DE|=4\left(\dfrac{1}{\sin^2\alpha}+\dfrac{1}{\cos^2\alpha}\right)=\dfrac{4}{\sin^2\alpha\cos^2\alpha}\ge 16$。

\noindent\textbf{教师追问：}由 $y^2=4x$ 可读出 $p$ 的值是多少？\\

\textbf{预期回答：}$p=2$。\par

\noindent\textbf{教师追问：}为什么第二条弦的公式里会出现 $\cos^2\alpha$？\\

\textbf{预期回答：}因为两条直线互相垂直，若一条倾斜角为 $\alpha$，另一条可视为 $\alpha+\dfrac\pi2$，故 $\sin^2\left(\alpha+\dfrac\pi2\right)=\cos^2\alpha$。\par

\noindent\textbf{教师追问：}最后的最值可转化为哪个基本不等式问题？\\

\textbf{预期回答：}转化为 $\sin^2\alpha\cos^2\alpha\le \dfrac14$，从而得到最小值。\par

\paragraph{10分钟：小结与规范表达}

教师示范本题的规范书写：先由标准式确定焦点与参数，再调用焦点弦公式，最后用三角恒等变形求最值。强调高考书写中不能只写“由结论知”，应标明结论适用于“过焦点的抛物线弦”。

\noindent\textbf{教师追问：}本题若不用焦点弦公式，能否做？\\

\textbf{预期回答：}能做，但会更繁琐，需要设直线方程、联立并用韦达定理。\par

\noindent\textbf{教师追问：}什么时候不建议强行使用坐标联立硬算？\\

\textbf{预期回答：}当题目已经明确给出过焦点、方向或垂直关系时，应优先利用结构。\par

\paragraph{5分钟：当堂检测一}

口头检测：若抛物线 $y^2=2px$ 的焦点弦与对称轴夹角越接近 $\dfrac\pi2$，则弦长如何变化；若接近 $0$，图形上意味着什么。

\noindent\textbf{教师追问：}由 $|MN|=\dfrac{2p}{\sin^2\alpha}$，当 $\alpha\to \dfrac\pi2$ 时，$|MN|$ 怎样？\\

\textbf{预期回答：}$|MN|\to 2p$，即趋近通径长。\par

\noindent\textbf{教师追问：}当 $\alpha\to 0$ 时，为什么弦长会变大？\\

\textbf{预期回答：}因为直线越来越接近对称轴，交点相距越来越远，公式中分母趋近 $0$。\par

\paragraph{15分钟：知识梳理二：椭圆中的焦半径与焦比弦}

进入第二课时。整理椭圆 $\dfrac{x^2}{a^2}+\dfrac{y^2}{b^2}=1$ 中过右焦点 $F(c,0)$ 的焦点弦关系：
$|AF|=\dfrac{b^2}{a+c\cos\alpha}$，$|BF|=\dfrac{b^2}{a-c\cos\alpha}$。
进一步得到
$\lambda=\dfrac{|AF|}{|BF|}=\dfrac{a-c\cos\alpha}{a+c\cos\alpha}$。
令 $e=\dfrac ca$，可化为
$\left|\cos\alpha\right|=\dfrac{1}{e}\left|\dfrac{\lambda-1}{\lambda+1}\right|$。强调：这里的结构判断可借助统一视角，但正式求解应回到定义与代数变形。

\noindent\textbf{教师追问：}在横轴椭圆中，为什么比例关系最终对应 $\cos\alpha$ 而不是 $\sin\alpha$？\\

\textbf{预期回答：}因为焦点在 $x$ 轴上，方向与长轴的投影关系自然落在 $\cos\alpha$ 上。\par

\noindent\textbf{教师追问：}若 $\lambda=1$，说明这条焦点弦具有什么位置特征？\\

\textbf{预期回答：}说明 $|AF|=|BF|$，即焦点是弦的中点，对应弦垂直于长轴，是通径。\par

\paragraph{20分钟：例题研讨二}

使用题目 ID：$\text{2019-全国I卷-理科::10}$。教学目标：从“长度比 $+\,$另一长度关系”反推椭圆参数。分析路径：已知 $F_1(-1,0),F_2(1,0)$，故 $c=1$。设 $\lambda=\dfrac{|AF_2|}{|F_2B|}=2$。由横轴椭圆焦比弦关系得
$\left|\cos\theta\right|=\dfrac{1}{e}\left|\dfrac{\lambda-1}{\lambda+1}\right|=\dfrac{1}{3e}$。
另一方面，题设还有 $|AB|=|BF_1|$，应引导学生转回高中方法：由椭圆定义，点 $B$ 满足 $|BF_1|+|BF_2|=2a$，而 $|AF_2|=2|F_2B|$，故 $|AB|=|AF_2|+|F_2B|=3|F_2B|$。联立 $|AB|=|BF_1|$ 与 $|BF_1|+|BF_2|=2a$，即可逐步得到参数关系，再确定椭圆方程。课堂上不替代原题选项，只强调推理链条。

\noindent\textbf{教师追问：}本题最先读出的确定量是什么？\\

\textbf{预期回答：}$c=1$，以及比例 $\lambda=2$。\par

\noindent\textbf{教师追问：}为什么可以把 $|AB|$ 写成 $|AF_2|+|F_2B|$？\\

\textbf{预期回答：}因为焦点 $F_2$ 在线段 $AB$ 上，且 $A,B$ 是过 $F_2$ 的直线与椭圆的两个交点。\par

\noindent\textbf{教师追问：}点 $B$ 在椭圆上时，哪条定义关系一定成立？\\

\textbf{预期回答：}$|BF_1|+|BF_2|=2a$。\par

\paragraph{12分钟：例题研讨三}

使用题目 ID：$\text{2020-全国II卷-理科::19}$。本题用于强化“通径长度与离心率”的联系。题干中“过 $F$ 且与 $x$ 轴垂直的直线”意味着这是一条通径所在直线。椭圆部分有 $|AB|=\dfrac{2b^2}{a}$；抛物线部分，若 $y^2=2px$，则过焦点且垂直于轴的弦长为 $2p$。又因两曲线焦点重合，且椭圆中心与抛物线顶点重合，可建立 $c=\dfrac p2$。由 $|CD|=\dfrac43|AB|$ 得
$2p=\dfrac43\cdot \dfrac{2b^2}{a}$，再结合 $c^2=a^2-b^2$ 与 $c=\dfrac p2$ 求离心率。第二问用于课上点拨，不作全解展开，保留给学生课后整理。

\noindent\textbf{教师追问：}为什么这里首先抓“与 $x$ 轴垂直且过焦点”这一条件？\\

\textbf{预期回答：}因为这直接指向通径，不必先联立求交点。\par

\noindent\textbf{教师追问：}椭圆与抛物线在这条直线上的弦分别是什么特殊弦？\\

\textbf{预期回答：}都是通径。\par

\noindent\textbf{教师追问：}由焦点重合可得到哪条参数关系？\\

\textbf{预期回答：}$c=\dfrac p2$。\par

\paragraph{10分钟：迁移训练与多选题点拨}

使用题目 ID：$\text{2023-新高考II卷::10}$ 与 $\text{2025-全国I卷::10}$ 进行思路点拨，不展开完整解答。前者抓“过焦点的已知斜率直线”，先由 $k=-\sqrt3$ 得 $\sin^2\alpha=\dfrac34$，从而迅速得到弦长或相关量；后者抓“过焦点的一条弦及其垂线、准线垂足”，优先建立焦点弦长度、焦半径与到准线距离的对应，再回到定义法判断选项。

\noindent\textbf{教师追问：}面对选择题，什么情况下可以先算结构量再判断选项？\\

\textbf{预期回答：}当选项考查的是弦长、比值、位置关系等可由结构快速得到的量时。\par

\noindent\textbf{教师追问：}在抛物线中，点到焦点距离与点到准线距离有什么关系？\\

\textbf{预期回答：}二者相等，即 $|PF|=d(P,l)$。\par

\paragraph{8分钟：课堂总结}

归纳本讲三条执行口令：
$1.$ 见“过焦点”先判断能否转入焦点弦；
$2.$ 见“斜率”先转成 $\sin^2\alpha,\cos^2\alpha$；
$3.$ 见“长度比”统一设 $\lambda=\dfrac{|AF|}{|BF|}$。
再次强调：射影视角只用于看出“分比结构”，解答必须落在高中方法上。

\noindent\textbf{教师追问：}今天最稳定的统一变量是什么？\\

\textbf{预期回答：}$\alpha$、$k$、$\lambda$ 这三个量最稳定，其中长度比常统一为 $\lambda=\dfrac{|AF|}{|BF|}$。\par

\noindent\textbf{教师追问：}为什么课堂上要反复强调‘回到高中方法’？\\

\textbf{预期回答：}因为高考评分依据的是课标内方法与规范表达。\par

\subsection*{课堂真题}

\paragraph{2017 年高考数学 全国I卷-理科第10题}

已知 $F$ 为抛物线 $C:y^2=4x$ 的焦点，过 $F$ 作两条互相垂直的直线 $l_1$，$l_2$，直线 $l_1$ 与 $C$ 交于 $A$、$B$ 两点，直线 $l_2$ 与 $C$ 交于 $D$、$E$ 两点，则 $|AB|+|DE|$ 的最小值为（ ）

\textbf{答案：}A\par

\textbf{解答：}方法一：如图，$l_1\perp l_2$，要使 $|AB|+|DE|$ 最小，则 $A$ 与 $D$，$B$，$E$ 关于 $x$ 轴对称，即直线 $DE$ 的斜率为 1，又直线 $l_2$ 过点 $(1,0)$，则直线 $l_2$ 的方程为 $y=x-1$，联立 $\begin{cases} y^2=4x \\ y=x-1 \end{cases}$，则 $y^2-4y-4=0$，\ensuremath{\therefore}$y_1+y_2=4$，$y_1y_2=-4$，\ensuremath{\therefore}$|DE|=\sqrt{1+1^2}\cdot|y_1-y_2|=\sqrt{2}\times\sqrt{4^2+16}=8$，\ensuremath{\therefore}$|AB|+|DE|$ 的最小值为 $2|DE|=16$．方法二：设直线 $l_1$ 的倾斜角为 $\theta$，则 $l_2$ 的倾斜角为 $\frac{\pi}{2}+\theta$，根据焦点弦长公式可得 $|AB|=\frac{2p}{\sin^2\theta}=\frac{4}{\sin^2\theta}$，$|DE|=\frac{4}{\sin^2(\frac{\pi}{2}+\theta)}=\frac{4}{\cos^2\theta}$，\ensuremath{\therefore}$|AB|+|DE|=\frac{4}{\sin^2\theta}+\frac{4}{\cos^2\theta}=\frac{4}{\sin^2\theta\cos^2\theta}=\frac{16}{\sin^2 2\theta}$，\ensuremath{\because}$0<\sin^2 2\theta\leq 1$，\ensuremath{\therefore}当 $\theta=45^\circ$ 时，$|AB|+|DE|$ 最小，最小为 16．故选：A．

\paragraph{2019 年高考数学 全国I卷-理科第10题}

已知椭圆 $C$ 的焦点为 $F_1(-1, 0)$, $F_2(1, 0)$，过 $F_2$ 的直线与 $C$ 交于 $A$, $B$ 两点．若 $|AF_2| = 2|F_2B|$, $|AB| = |BF_1|$, 则 $C$ 的方程为

\textbf{答案：}$B$\par

\textbf{解答：}如图，由已知可设 $|F_2B| = n$，则 $|AF_2| = 2n$, $|BF_1| = |AB| = 3n$，由椭圆的定义有 $2a = |BF_1| + |BF_2| = 4n$, $\therefore |AF_1| = 2a - |AF_2| = 2n$．在 $\triangle AF_1B$ 中，由余弦定理推论得 $\cos \angle F_1AB = \frac{4n^2 + 9n^2 - 9n^2}{2 \cdot 2n \cdot 3n} = \frac{1}{3}$．在 $\triangle AF_1F_2$ 中，由余弦定理得 $4n^2 + 4n^2 - 2 \cdot 2n \cdot 2n \cdot \frac{1}{3} = 4$，解得 $n = \frac{\sqrt{3}}{2}$．$\therefore 2a = 4n = 2\sqrt{3}$, $\therefore a = \sqrt{3}$, $\therefore b^2 = a^2 - c^2 = 3 - 1 = 2$, $\therefore$ 所求椭圆方程为 $\frac{x^2}{3} + \frac{y^2}{2} = 1$，故选 $B$．

\paragraph{2020 年高考数学 全国II卷-理科第19题}

已知椭圆$C_1:\dfrac{x^2}{a^2}+\dfrac{y^2}{b^2}=1(a>b>0)$的右焦点$F$与抛物线$C_2$的焦点重合，$C_1$的中心与$C_2$的顶点重合.过$F$且与$x$轴垂直的直线交$C_1$于$A,B$两点，交$C_2$于$C,D$两点，且$|CD|=\dfrac43|AB|$.
（1）求$C_1$的离心率；
（2）设$M$是$C_1$与$C_2$的公共点，若$|MF|=5$，求$C_1$与$C_2$的标准方程.

\textbf{答案：}（1）$\dfrac12$；（2）$C_1:\dfrac{x^2}{36}+\dfrac{y^2}{27}=1$，$C_2:y^2=12x$\par

\textbf{解答：}（1）$\because F(c,0)$，$AB\perp x$轴且与椭圆$C_1$相交于$A,B$两点，则直线$AB$的方程为$x=c$，联立$\begin{cases}x=c\\\dfrac{x^2}{a^2}+\dfrac{y^2}{b^2}=1\end{cases}$，解得$y=\pm\dfrac{b^2}{a}$，则$|AB|=\dfrac{2b^2}{a}$，抛物线$C_2$的方程为$y^2=4cx$，联立$\begin{cases}x=c\\y^2=4cx\end{cases}$，解得$y=\pm2c$，$\therefore |CD|=4c$，$\because CD=\dfrac43AB$，即$4c=\dfrac43\times\dfrac{2b^2}{a}$，$\therefore 2b^2=3ac$，即$2(a^2-c^2)=3ac$，即$2c^2+3ac-2a^2=0$，即$2e^2+3e-2=0$，$\because 0<e<1$，解得$e=\dfrac12$.
（2）由（1）知$a=2c$，$b=\sqrt3c$，椭圆$C_1$的方程为$\dfrac{x^2}{4c^2}+\dfrac{y^2}{3c^2}=1$，联立$\begin{cases}y^2=4cx\\\dfrac{x^2}{4c^2}+\dfrac{y^2}{3c^2}=1\end{cases}$，消去$y$并整理得$3x^2+16cx-12c^2=0$，解得$x=\dfrac23c$或$x=-6c$（舍去），由抛物线的定义可得$|MF|=\dfrac23c+c=\dfrac{5c}{3}=5$，解得$c=3$，$\therefore$曲线$C_1$的标准方程为$\dfrac{x^2}{36}+\dfrac{y^2}{27}=1$，曲线$C_2$的标准方程为$y^2=12x$.

\paragraph{2023 年高考数学 新高考II卷第10题}

设$O$为坐标原点，直线$y=-\sqrt{3}(x-1)$过抛物线$C:y^{2}=2px(p>0)$的焦点，且与$C$交于$M$，$N$两点，$l$为$C$的准线，则（ \quad ）

\textbf{答案：}AC\par

\textbf{解答：}A选项：直线$y=-\sqrt{3}(x-1)$过点$(1,0)$，所以抛物线$C:y^{2}=2px(p>0)$的焦点$F(1,0)$，所以$\dfrac{p}{2}=1$，$p=2$，$2p=4$，则A正确，且抛物线$C$的方程为$y^{2}=4x$.
B选项：设$M(x_{1},y_{1})$，$N(x_{2},y_{2})$，由$\begin{cases}y=-\sqrt{3}(x-1)\\ y^{2}=4x\end{cases}$消去$y$并化简得$3x^{2}-10x+3=(x-3)(3x-1)=0$，解得$x_{1}=3$，$x_{2}=\dfrac{1}{3}$，所以$|MN|=x_{1}+x_{2}+p=3+\dfrac{1}{3}+2=\dfrac{16}{3}$，B错误.
C选项：设$MN$的中点为$A$，$M,N,A$到直线$l$的距离分别为$d_{1},d_{2},d$，因为$d=\dfrac{1}{2}(d_{1}+d_{2})=\dfrac{1}{2}(|MF|+|NF|)=\dfrac{1}{2}|MN|$，即$A$到直线$l$的距离等于$|MN|$的一半，所以以$MN$为直径的圆与直线$l$相切，C正确.
D选项：直线$y=-\sqrt{3}(x-1)$即$\sqrt{3}x+y-\sqrt{3}=0$，$O$到直线的距离为$d=\dfrac{\sqrt{3}}{2}$，所以三角形$OMN$的面积为$\dfrac{1}{2}\times\dfrac{16}{3}\times\dfrac{\sqrt{3}}{2}=\dfrac{4\sqrt{3}}{3}$，由上述分析可知$y_{1}=-\sqrt{3}(3-1)=-2\sqrt{3}$，$y_{2}=-\sqrt{3}\left(\dfrac{1}{3}-1\right)=\dfrac{2\sqrt{3}}{3}$，所以$|OM|=\sqrt{3^{2}+(-2\sqrt{3})^{2}}=\sqrt{21}$，$|ON|=\sqrt{\left(\dfrac{1}{3}\right)^{2}+\left(\dfrac{2\sqrt{3}}{3}\right)^{2}}=\dfrac{\sqrt{13}}{3}$，所以三角形$OMN$不是等腰三角形，D错误.

\paragraph{2025 年高考数学 全国I卷第10题}

设抛物线 $C: y^2 = 6x$ 的焦点为 $F$，过 $F$ 的直线交 $C$ 于 $A$、$B$，过 $F$ 且垂直于 $AB$ 的直线交 $l: x = -\frac{3}{2}$ 于 $E$，过点 $A$ 作准线 $l$ 的垂线，垂足为 $D$，则（ ）

\textbf{答案：}ACD\par

\textbf{解答：}法一：对于 A，抛物线 $y^2=6x$，$p=3$，准线方程 $x=-\frac{3}{2}$，焦点 $F\left(\frac{3}{2},0\right)$，由抛物线定义，$|AD| = |AF|$，故 A 正确；对于 B，过点 $B$ 作准线垂线交于点 $P$，易证 $\triangle ADE \cong \triangle AFE$，$\triangle BEP \cong \triangle BEF$，则 $\angle AED = \angle AEF$，$\angle BEP = \angle BEF$，又 $\angle AED+\angle AEF+\angle BEP+\angle BEF=180^\circ$，所以 $\angle AEF+\angle BEF=90^\circ$，即 $\angle AEB=90^\circ$，$AB$ 为斜边，则 $|AE| < |AB|$，故 B 错误；对于 C，设直线 $AB$ 方程为 $x = my + \frac{3}{2}$，与 $y^2=6x$ 联立得 $y^2-6my-9=0$，$y_1+y_2=6m$，$y_1y_2=-9$，$|AB| = x_1+x_2+p = m(y_1+y_2)+3+3 = 6m^2+6 \ge 6$，当且仅当 $m=0$ 时取等号，故 C 正确；对于 D，在 $\triangle ABE$ 与 $\triangle AEF$ 中，$\angle BAE = \angle EAF$，所以 $\triangle ABE \sim \triangle AEF$，则 $\frac{AE}{AB} = \frac{AF}{AE}$，即 $AE^2 = AF \cdot AB$，同理 $BE^2 = BF \cdot AB$，$AF \cdot BF = (x_1+\frac{3}{2})(x_2+\frac{3}{2}) = (my_1+3)(my_2+3) = m^2 y_1 y_2 + 3m(y_1+y_2)+9 = -9m^2+18m^2+9 = 9(m^2+1)$，$AB = 6(m^2+1)$，则 $AE^2 \cdot BE^2 = BF \cdot AF \cdot AB^2 = 9(m^2+1) \times 36(m^2+1)^2$，所以 $AE \cdot BE = 3(m^2+1)^{\frac{1}{2}} \times 6(m^2+1) = 18(m^2+1)^{\frac{3}{2}} \ge 18$，故 D 正确. 故选：ACD. 法二：略.

\subsection*{课后作业}

\begin{enumerate}

\item 2020 年高考数学 全国II卷-理科第19题

\item 2023 年高考数学 新高考II卷第10题

\item 2025 年高考数学 全国I卷第10题

\end{enumerate}

\noindent\textbf{作业答案索引：}

\begin{itemize}

\item 2020 年高考数学 全国II卷-理科第19题：（1）$\dfrac12$；（2）$C_1:\dfrac{x^2}{36}+\dfrac{y^2}{27}=1$，$C_2:y^2=12x$

\item 2023 年高考数学 新高考II卷第10题：AC

\item 2025 年高考数学 全国I卷第10题：ACD

\end{itemize}

\subsection*{板书设计}

一、概念与识别
$1.$ 焦点弦：过焦点的弦
$2.$ 焦半径：曲线上点到焦点的距离
$3.$ 焦比弦：$\lambda=\dfrac{|AF|}{|BF|}$

二、一般弦长公式
若 $A(x_1,y_1),B(x_2,y_2)$ 在直线 $y=kx+m$ 上，
\[
|AB|=\sqrt{1+k^2}\,|x_1-x_2|
\]

三、抛物线 $y^2=2px$
焦点：$F\left(\dfrac p2,0\right)$
通径长：$2p$
过焦点弦：
\[
|MN|=\frac{2p}{\sin^2\alpha}
\]
斜率转化：
\[
\sin^2\alpha=\frac{k^2}{1+k^2},\qquad \cos^2\alpha=\frac{1}{1+k^2}
\]

四、椭圆 $\dfrac{x^2}{a^2}+\dfrac{y^2}{b^2}=1$
$c^2=a^2-b^2,\ e=\dfrac ca$
过右焦点 $F(c,0)$ 的焦点弦：
\[
|AF|=\frac{b^2}{a+c\cos\alpha},\qquad |BF|=\frac{b^2}{a-c\cos\alpha}
\]
\[
\lambda=\frac{|AF|}{|BF|}=\frac{a-c\cos\alpha}{a+c\cos\alpha}
\]
\[
\left|\cos\alpha\right|=\frac{1}{e}\left|\frac{\lambda-1}{\lambda+1}\right|
\]
通径长：
\[
\frac{2b^2}{a}
\]

五、例题结论提炼
$\text{2017-全国I卷-理科::10}$：
\[
|AB|+|DE|=\frac{4}{\sin^2\alpha}+\frac{4}{\cos^2\alpha}\ge 16
\]

六、作答提醒
结构判断可快，书写必须规范：参数确定 $\to$ 公式调用条件 $\to$ 代数化简。

\subsection*{易错点}

\begin{itemize}

\item 把“过焦点的弦”与“以焦点为中点的弦”混同。实际上只有特殊情形，如椭圆通径，焦点才是弦中点。

\item 已知斜率 $k$ 后直接把 $k$ 当成 $\sin\alpha$ 或 $\cos\alpha$ 使用。正确转化是 $k=\tan\alpha$。

\item 使用椭圆焦比弦关系时忽视曲线方向，横轴型通常对应 $\cos\theta$，纵轴型通常对应 $\sin\theta$。

\item 在长度关系中漏掉焦点在线段内部还是在线段延长线上这一位置判断，导致 $|AB|$ 的分拆错误。

\item 把射影视角当作正式证明。课堂中它只能帮助识别“分比结构”，高考解答仍需写出坐标法、定义法或韦达定理过程。

\item 抛物线中只记住 $|MN|=\dfrac{2p}{\sin^2\alpha}$，却忘记其适用前提是“弦过焦点”。

\item 遇到通径不会立刻识别：抛物线中‘过焦点且垂直于对称轴’，椭圆中‘过焦点且垂直于长轴’，都应优先想到通径长度。

\end{itemize}

\subsection*{教师备注}

本模块适合一轮复习中的专题推进，建议按“公式识别 $\to$ 高考例题 $\to$ 统一变量 $\to$ 规范表达”组织。教学时要把握两个边界：其一，不把焦点弦公式讲成死记结论，必须让学生知道它们常可由定义、联立、韦达定理推出；其二，不让统一视角越界为非课标证明。课堂提问尽量围绕三个触发词展开：‘过焦点’、‘斜率或倾斜角’、‘长度比’。对于基础较弱的班级，可弱化比例—离心率公式的推导，保留其应用；对于学有余力的班级，可补充参数化意识：抛物线 $y^2=2px$ 上可设点为 $\left(\dfrac p2 t^2,pt\right)$，两点构成焦点弦时常有 $t_1t_2=-1$，但正式书写时仍建议回到课标内可接受的联立与韦达关系。课后批改重点看三点：是否先判结构，是否正确完成 $k\leftrightarrow \alpha$ 转化，是否在长度分拆时判断了焦点位置。

\clearpage

\section{第8课：切线、极线与阿基米德三角形}

\noindent 建议课时：2课时。

\subsection*{教学目标}

\begin{itemize}

\item 掌握圆锥曲线中三类基础结论的高中化表达：切线方程、切点弦方程、由定点引切线所得切点弦性质。

\item 会在椭圆、抛物线背景下识别并使用“手电筒模型”：若定点 $P$ 与两交点 $A,B$ 的连线满足斜率和或斜率积为定值，则优先预判弦 $AB$ 的“过定点”或“定斜率”性质。

\item 理解阿基米德三角形的基本结构：弦 $AB$ 与两端点切线交点构成的三角形，并能把“切线条件”转化为“交点轨迹、垂直关系、面积关系”。

\item 能区分“唯一公共点”与“相切”这两个不同命题，形成严格的判定意识。

\item 在正式解答中，能回到高中课标允许的方法，综合使用坐标法、定义法、韦达定理、点差法、向量或平面几何完成证明与计算。

\end{itemize}

\subsection*{重点与难点}

\begin{itemize}

\item 椭圆、抛物线的点式切线与切点弦公式的识别与使用。

\item 由“斜率和为定值”预判“弦过定点”的结构意识。

\item 阿基米德三角形中“弦、两切线、交点”的整体处理。

\item 切线存在性与特殊位置的补检验。

\end{itemize}

\subsection*{高观点}

从统一视角看，圆锥曲线中的“切线”“切点弦”“由外点引两条切线所连两切点”可以放在同一组对偶关系中理解：点与线之间存在对应，外点对应一条极线，切点弦正是该外点的极线；而两点连成一条弦、两切线交于一点，也构成一种对偶交换。这个视角的课堂作用仅限于帮助学生做结构预判：一旦题目同时出现“外点”“两条切线”“两切点连线”，就应想到可能存在固定直线；一旦出现“弦的两端点处切线交于一点”，就应把该点与弦作为整体研究。正式解题不使用射影几何证明，只把它当作概念溯源与结构导航。

\subsection*{高中落地}

回到高中方法，可把上述统一视角翻译为三条可操作规则：$1.$ 若已知切点 $P(x_0,y_0)$，优先直接写点式切线；$2.$ 若已知曲线外一点 $P(x_0,y_0)$ 向曲线引两条切线，切点为 $A,B$，优先写切点弦 $AB$ 的方程；$3.$ 若定点 $P$ 与弦端点 $A,B$ 的连线斜率满足 $k_{PA}+k_{PB}$ 或 $k_{PA}k_{PB}$ 为定值，则设动直线、联立曲线、用韦达定理把条件转成弦 $AB$ 的方程，从而证明“过定点”或“定斜率”。阿基米德三角形问题则优先把面积写成 $S_{\triangle ABQ}=\dfrac{1}{2}|AB|\cdot d(Q,AB)$，并结合切线交点的坐标或准线性质求解。

\subsection*{教学流程}

\paragraph{8分钟：导入与目标建立}

教师出示本课主题关键词：$\text{切线方程}$、$\text{切点弦}$、$\text{极线预判}$、$\text{阿基米德三角形}$。先不进入题目细节，只组织学生回忆：圆锥曲线综合题中，哪些条件一出现就提示我们“不要逐条算两条切线”，而要寻找整体结构。

\noindent\textbf{教师追问：}看到“过曲线外一点作两条切线，切点为 $A,B$”时，你第一反应应是什么？\\

\textbf{预期回答：}应优先想到切点弦，可能存在固定直线或易写出的直线方程。\par

\noindent\textbf{教师追问：}看到“直线与曲线交于 $A,B$，且 $P A, P B$ 的斜率和为定值”时，你会预判哪类结论？\\

\textbf{预期回答：}常可预判弦 $AB$ 过定点，或弦 $AB$ 斜率为定值。\par

\noindent\textbf{教师追问：}看到“过 $A,B$ 分别作切线，两切线交于一点”时，为什么不能把三条直线分开处理？\\

\textbf{预期回答：}因为这是一个整体结构，即阿基米德三角形，面积、垂直、轨迹常通过整体关系求出。\par

\paragraph{14分钟：基础公式梳理}

教师板书并带学生口述两组公式。椭圆 $\dfrac{x^2}{a^2}+\dfrac{y^2}{b^2}=1$ 在点 $P(x_0,y_0)$ 处的切线为 $\dfrac{x_0x}{a^2}+\dfrac{y_0y}{b^2}=1$；抛物线 $y^2=2px$ 的点式切线为 $y_0y=p(x+x_0)$，对应地对 $x^2=2py$ 可写成 $xx_0=p(y+y_0)$。同时强调：若曲线外一点为 $P(x_0,y_0)$，其所引两条切线的切点连线，也常可由同样的式子直接给出。再强调判定误区：直线与圆锥曲线只有一个公共点，不一定是切线。

\noindent\textbf{教师追问：}已知切点而不是已知斜率时，为什么不建议先设斜率再判别式为零？\\

\textbf{预期回答：}因为点式切线公式更直接，计算量小，也更稳定。\par

\noindent\textbf{教师追问：}对抛物线 $x^2=2py$，点式切线应如何写？\\

\textbf{预期回答：}$xx_0=p(y+y_0)$。\par

\noindent\textbf{教师追问：}“只有一个公共点”为什么不能直接推出“相切”？\\

\textbf{预期回答：}因为某些特殊位置也可能只有一个公共点，但并不满足切线定义中的重合极限方向条件，必须具体核查。\par

\paragraph{18分钟：例1结构识别：手电筒模型}

选用题目 $\texttt{2017-全国I卷-理科::20}$。本环节只围绕第 $(2)$ 问的结构识别与证明路径组织课堂，不改写原题。教师引导学生观察：定点为 $P_2$，动直线 $l$ 与椭圆交于 $A,B$，且 $P_2A,P_2B$ 的斜率和为 $-1$。这正对应技巧 $47$ 的“手电筒模型”。先做结构预判：弦 $AB$ 即直线 $l$ 很可能过定点。随后用高中方法完成：设 $l:y=kx+m$，与椭圆联立，设两交点坐标为 $A(x_1,y_1),B(x_2,y_2)$，将 $k_{P_2A}+k_{P_2B}=-1$ 化为关于 $x_1,x_2,y_1,y_2$ 的关系，再借助韦达定理消元，最终整理出直线恒过定点的结论。

\noindent\textbf{教师追问：}这道题里谁是“手电筒模型”的定点？谁是那条动弦？\\

\textbf{预期回答：}定点是 $P_2$，动弦是 $AB$，也就是直线 $l$。\par

\noindent\textbf{教师追问：}条件 $k_{P_2A}+k_{P_2B}=-1$ 更可能导向“定斜率”还是“过定点”？\\

\textbf{预期回答：}更可能导向“过定点”。\par

\noindent\textbf{教师追问：}若设 $l:y=kx+m$，联立椭圆后，韦达定理能提供哪两类信息？\\

\textbf{预期回答：}能提供两交点横坐标的和与积，从而把斜率和条件转成 $k,m$ 之间的关系。\par

\paragraph{8分钟：例1结论回收与方法命名}

教师收束为一条板书结论：当题中出现“定点 $P$ + 弦端点连线斜率和为定值”时，优先预判弦过定点，但证明必须回到“设直线 $+$ 联立 $+$ 韦达 $+$ 验证定点”。此处点明该题对应定点是已知答案中的固定点，但课堂不把答案当作起点，而强调如何从结构走向证明。

\noindent\textbf{教师追问：}为什么这里的“极线预判”只能作为导航，而不能直接代替证明？\\

\textbf{预期回答：}因为高考正式解答仍需用坐标法、韦达定理等高中方法写清推导。\par

\noindent\textbf{教师追问：}如果把这类题机械记成公式，有什么风险？\\

\textbf{预期回答：}一旦曲线类型、定点位置或条件形式略变，就容易误套公式。\par

\paragraph{15分钟：例2切点弦：由外点到固定直线}

选用题目 $\texttt{2019-全国III卷-理科::21}$ 与 $\texttt{2019-全国III卷-文科::21}$ 的共同结构，只讨论共同的第 $(1)$ 问类型：过外点 $D$ 作抛物线 $C$ 的两条切线，切点为 $A,B$，证明直线 $AB$ 过定点。教师引导学生把“过外点作两条切线”直接翻译成“切点弦”。对抛物线 $y=\dfrac{x^2}{2}$ 可改写为 $x^2=2y$，此时 $p=1$，若外点 $D(x_0,-\dfrac12)$，则切点弦可直接写成 $xx_0=y-\dfrac12$。于是学生立即读出：无论 $x_0$ 如何变化，点 $\left(0,\dfrac12\right)$ 恒在直线 $AB$ 上。

\noindent\textbf{教师追问：}把 $y=\dfrac{x^2}{2}$ 改写成标准式后，参数 $p$ 是多少？\\

\textbf{预期回答：}$x^2=2y$，所以 $p=1$。\par

\noindent\textbf{教师追问：}由外点 $D(x_0,-\dfrac12)$ 向抛物线引切线，其切点弦为什么能直接写？\\

\textbf{预期回答：}因为这是标准的抛物线切点弦公式 $xx_0=p(y+y_0)$ 的直接应用。\par

\noindent\textbf{教师追问：}定点如何从直线方程中快速读出？\\

\textbf{预期回答：}把直线方程中 $x=0$ 代入，可得 $y=\dfrac12$，所以定点是 $\left(0,\dfrac12\right)$。\par

\paragraph{7分钟：第一课时小结与即时检测}

教师组织学生用一句话分别概括两种常见入口：$1.$ 已知切点，用点式切线；$2.$ 已知外点引两切线，用切点弦；$3.$ 已知定点与两交点连线斜率和或积为定值，用手电筒模型预判，再回到联立与韦达。即时检测采用口答，不另给新题。

\noindent\textbf{教师追问：}若题目给的是切点，核心公式是什么？\\

\textbf{预期回答：}点式切线公式。\par

\noindent\textbf{教师追问：}若题目给的是曲线外一点和两条切线，核心公式是什么？\\

\textbf{预期回答：}切点弦公式。\par

\noindent\textbf{教师追问：}若题目给的是斜率和为定值，核心思想是什么？\\

\textbf{预期回答：}先结构预判，再设直线联立，用韦达证明。\par

\paragraph{10分钟：第二课时导入：阿基米德三角形}

教师给出名称解释：若弦 $AB$ 的两个端点处切线相交于 $Q$，则 $\triangle ABQ$ 称为阿基米德三角形；对抛物线情形同理。强调其课堂价值不在记名词，而在于把“弦、两切线、交点”视为一个整体。若再叠加“弦过焦点”或“面积最值”等条件，往往可迅速组织解法。

\noindent\textbf{教师追问：}为什么阿基米德三角形本质上是一个“整体模型”？\\

\textbf{预期回答：}因为弦、两切线和交点之间不是孤立关系，而是互相制约的整体结构。\par

\noindent\textbf{教师追问：}遇到面积问题时，最先应写什么形式？\\

\textbf{预期回答：}先写成 $\dfrac12\times \text{底} \times \text{高}$。\par

\paragraph{18分钟：例3抛物线切线与面积最值}

选用题目 $\texttt{2021-全国乙卷-理科::21}$。本环节聚焦第 $(2)$ 问的思路训练：点 $P$ 在圆 $M$ 上，$PA,PB$ 是抛物线 $C$ 的两条切线，$A,B$ 是切点，求 $\triangle PAB$ 面积最大值。教师先引导学生辨型：这不是单纯切线方程题，而是“外点引两切线 $+$ 阿基米德三角形面积最值”。解法上可先利用抛物线切点参数或切线式，求出切点弦 $AB$ 与点 $P$ 的距离，再写面积。由于 $P$ 受圆约束，最后转化为点到直线距离或某代数式的最值问题。课堂中强调：这里的“极线”只帮助我们想到 $AB$ 与 $P$ 一一对应，正式推导仍用坐标法与最值方法完成。

\noindent\textbf{教师追问：}本题中三角形的三个顶点分别由哪三个元素给出？\\

\textbf{预期回答：}顶点是外点 $P$ 和两个切点 $A,B$。\par

\noindent\textbf{教师追问：}若直接分别求两条切线，会出现什么问题？\\

\textbf{预期回答：}参数多、代数运算重，结构不清晰。\par

\noindent\textbf{教师追问：}面积最值为什么更适合先抓 $AB$ 这条整体直线？\\

\textbf{预期回答：}因为面积可写为 $\dfrac12|AB|\cdot d(P,AB)$，而 $AB$ 恰是切点弦，整体处理更自然。\par

\paragraph{12分钟：例4补充视角：唯一公共点不等于相切}

选用题目 $\texttt{2021-天津卷::18}$ 的第 $(\mathrm{II})$ 问作提醒性训练。教师不展开全解，只聚焦题眼：题目说“直线 $l$ 与椭圆有唯一的公共点 $M$”，学生很容易立刻写“故 $l$ 为切线”。本题可用来专门纠偏：在圆锥曲线中，“唯一公共点”与“相切”不能机械等同，必须结合具体位置与几何条件核查。再结合题中 $MP\parallel BF$ 等条件，说明正确做法是综合坐标关系判断，而不是偷换概念。

\noindent\textbf{教师追问：}本题最容易发生的逻辑跳步是什么？\\

\textbf{预期回答：}把“唯一公共点”直接当成“相切”。\par

\noindent\textbf{教师追问：}为什么这类表述在圆里常成立，到圆锥曲线里就要更谨慎？\\

\textbf{预期回答：}因为圆的对称性更强，而圆锥曲线中某些特殊方向的直线也可能只交于一个点，必须具体分析。\par

\paragraph{12分钟：综合提升：从手电筒到阿基米德三角形}

教师串联题目 $\texttt{2020-全国I卷-理科::20}$。本题中 $P$ 在线 $x=6$ 上运动，$PA,PB$ 与椭圆再交于 $C,D$，证明直线 $CD$ 过定点。课堂不重做全题，而是用于横向比较：它仍是“定点 $P$ 发出两束射线”的手电筒模型，只是已知的不是斜率和，而是由构型本身诱导出的代数关系。学生需要形成迁移：不必拘泥于“斜率和为常数”的字面形式，只要是“定点 $+$ 两条连线 $+$ 两交点成弦”，就可以优先猜测弦的定点性质。

\noindent\textbf{教师追问：}这道题和 $\texttt{2017-全国I卷-理科::20}$ 的共性结构是什么？\\

\textbf{预期回答：}都有定点向曲线上两点连线，再研究这两点所成弦的性质。\par

\noindent\textbf{教师追问：}它与直接套技巧 $47$ 的差别又在哪里？\\

\textbf{预期回答：}这里未直接给出斜率和为定值，需要通过联立与构型关系推出弦的定点性质。\par

\paragraph{6分钟：课堂总结与作业布置}

教师带学生归纳本课四句工作语句：$1.$ 已知切点，先写点式切线；$2.$ 已知外点引两切线，先想切点弦；$3.$ 已知定点与两交点连线的斜率和或积有约束，先想手电筒模型；$4.$ 已知弦两端点处切线相交，先整体看作阿基米德三角形。最后布置分层作业。

\noindent\textbf{教师追问：}今天最值得你带走的一条“先看结构，再下笔”的经验是什么？\\

\textbf{预期回答：}不要见切线就逐条硬算，而要先识别切点弦、手电筒模型或阿基米德三角形。\par

\subsection*{课堂真题}

\paragraph{2017 年高考数学 全国I卷-理科第20题}

已知椭圆 $C:\frac{x^2}{a^2}+\frac{y^2}{b^2}=1$（$a>b>0$），四点 $P_1(1,1)$，$P_2(0,1)$，$P_3(-1,\frac{\sqrt{3}}{2})$，$P_4(1,\frac{\sqrt{3}}{2})$ 中恰有三点在椭圆 $C$ 上．
（1）求 $C$ 的方程；
（2）设直线 $l$ 不经过 $P_2$ 点且与 $C$ 相交于 $A$，$B$ 两点．若直线 $P_2A$ 与直线 $P_2B$ 的斜率的和为 $-1$，证明：$l$ 过定点．

\textbf{答案：}（1）$\frac{x^2}{4}+y^2=1$；（2）证明见解析，定点为 $(2,-1)$\par

\textbf{解答：}解：（1）根据椭圆的对称性，$P_3(-1,\frac{\sqrt{3}}{2})$，$P_4(1,\frac{\sqrt{3}}{2})$ 两点必在椭圆 $C$ 上，又 $P_4$ 的横坐标为 1，\ensuremath{\therefore}椭圆必不过 $P_1(1,1)$，\ensuremath{\therefore}$P_2(0,1)$，$P_3(-1,\frac{\sqrt{3}}{2})$，$P_4(1,\frac{\sqrt{3}}{2})$ 三点在椭圆 $C$ 上．把 $P_2(0,1)$，$P_3(-1,\frac{\sqrt{3}}{2})$ 代入椭圆 $C$，得：$\begin{cases} \frac{0}{a^2}+\frac{1}{b^2}=1 \\ \frac{1}{a^2}+\frac{3}{4b^2}=1 \end{cases}$，解得 $a^2=4$，$b^2=1$，\ensuremath{\therefore}椭圆 $C$ 的方程为 $\frac{x^2}{4}+y^2=1$．
证明：（2）(1)当斜率不存在时，设 $l:x=m$，$A(m,y_A)$，$B(m,-y_A)$，\ensuremath{\because}直线 $P_2A$ 与直线 $P_2B$ 的斜率的和为 $-1$，\ensuremath{\therefore}$\frac{y_A-1}{m-0}+\frac{-y_A-1}{m-0}=\frac{-2}{m}=-1$，解得 $m=2$，此时 $l$ 过椭圆右顶点，不存在两个交点，故不满足．(2)当斜率存在时，设 $l:y=kx+t$，($t\neq1$)，$A(x_1,y_1)$，$B(x_2,y_2)$，联立 $\begin{cases} y=kx+t \\ \frac{x^2}{4}+y^2=1 \end{cases}$，整理，得 $(1+4k^2)x^2+8ktx+4t^2-4=0$，$\Delta>0$，$x_1+x_2=-\frac{8kt}{1+4k^2}$，$x_1x_2=\frac{4t^2-4}{1+4k^2}$，则 $k_{P_2A}+k_{P_2B}=\frac{y_1-1}{x_1}+\frac{y_2-1}{x_2}=\frac{kx_1+t-1}{x_1}+\frac{kx_2+t-1}{x_2}=2k+\frac{t-1}{x_1}+\frac{t-1}{x_2}=2k+(t-1)\frac{x_1+x_2}{x_1x_2}=2k+(t-1)\frac{-\frac{8kt}{1+4k^2}}{\frac{4t^2-4}{1+4k^2}}=2k-\frac{2k(t-1)}{t+1}=\frac{2k(t+1)-2k(t-1)}{t+1}=\frac{4k}{t+1}=-1$，又 $t\neq1$，\ensuremath{\therefore}$t=-2k-1$，此时 $\Delta=64k^2-16(1+4k^2)(t^2-1)$，代入得 $\Delta=-64k$，存在 $k$，使得 $\Delta>0$ 成立，\ensuremath{\therefore}直线 $l$ 的方程为 $y=kx-2k-1$，当 $x=2$ 时，$y=-1$，\ensuremath{\therefore}$l$ 过定点 $(2,-1)$．

\paragraph{2019 年高考数学 全国III卷-理科第21题}

已知曲线 $C: y=\dfrac{x^2}{2}$，$D$ 为直线 $y=-\dfrac{1}{2}$ 上的动点，过 $D$ 作 $C$ 的两条切线，切点分别为 $A$，$B$。
（1）证明：直线 $AB$ 过定点；
（2）若以 $E\left(0,\dfrac{5}{2}\right)$ 为圆心的圆与直线 $AB$ 相切，且切点为线段 $AB$ 的中点，求四边形 $ADBE$ 的面积。

\textbf{答案：}（1）定点为 $\left(0,\dfrac{1}{2}\right)$；（2）四边形 $ADBE$ 的面积为 $3$ 或 $4\sqrt{2}$。\par

\textbf{解答：}（1）设 $D\left(t,-\dfrac{1}{2}\right)$，$A(x_1,y_1)$，则 $x_1^2=2y_1$。由于 $y'=x$，所以切线 $DA$ 的斜率为 $x_1$，故 $\dfrac{y_1+\frac{1}{2}}{x_1-t}=x_1$，整理得 $2tx_1-2y_1+1=0$。设 $B(x_2,y_2)$，同理可得 $2tx_2-2y_2+1=0$。故直线 $AB$ 的方程为 $2tx-2y+1=0$，所以直线 $AB$ 过定点 $\left(0,\dfrac{1}{2}\right)$。
（2）由（1）得直线 $AB$ 的方程为 $y=tx+\dfrac{1}{2}$。与 $y=\dfrac{x^2}{2}$ 联立得 $x^2-2tx-1=0$，于是 $x_1+x_2=2t$，$x_1x_2=-1$，$y_1+y_2=t(x_1+x_2)+1=2t^2+1$，$|AB|=\sqrt{1+t^2}|x_1-x_2|=\sqrt{1+t^2}\times\sqrt{(x_1+x_2)^2-4x_1x_2}=2(t^2+1)$。设 $d_1,d_2$ 分别为点 $D$，$E$ 到直线 $AB$ 的距离，则 $d_1=\sqrt{t^2+1}$，$d_2=\dfrac{2}{\sqrt{t^2+1}}$。因此四边形 $ADBE$ 的面积 $S=\dfrac{1}{2}|AB|(d_1+d_2)=(t^2+3)\sqrt{t^2+1}$。设 $M$ 为线段 $AB$ 的中点，则 $M\left(t,t^2+\dfrac{1}{2}\right)$。由于 $\overrightarrow{EM}\perp\overrightarrow{AB}$，而 $\overrightarrow{EM}=(t,t^2-2)$，$\overrightarrow{AB}$ 与向量 $(1,t)$ 平行，所以 $t+(t^2-2)t=0$，解得 $t=0$ 或 $t=\pm1$。当 $t=0$ 时，$S=3$；当 $t=\pm1$ 时，$S=4\sqrt{2}$。因此，四边形 $ADBE$ 的面积为 $3$ 或 $4\sqrt{2}$。

\paragraph{2021 年高考数学 全国乙卷-理科第21题}

已知抛物线$C:x^2=2py$（$p>0$）的焦点为$F$，且$F$与圆$M:x^2+(y+4)^2=1$上点的距离的最小值为4。
（1）求$p$；
（2）若点$P$在$M$上，$PA$，$PB$是$C$的两条切线，$A$，$B$是切点，求$\triangle PAB$面积的最大值。

\textbf{答案：}（1）$p=2$；（2）$20\sqrt{5}$\par

\textbf{解答：}（1）焦点$F(0,\frac{p}{2})$到圆$x^2+(y+4)^2=1$的最短距离为$\frac{p}{2}+3=4$，所以$p=2$。
（2）抛物线$y=\frac14x^2$，设$A(x_1,y_1)$，$B(x_2,y_2)$，$P(x_0,y_0)$，得$l_{PA}: y=\frac12x_1(x-x_1)+y_1=\frac12x_1x-\frac14x_1^2=\frac12x_1x-y_1$，$l_{PB}: y=\frac12x_2x-y_2$，且$x_0^2=-y_0^2-8y_0-15$。$l_{PA}$，$l_{PB}$都过点$P(x_0,y_0)$，则$\begin{cases} y_0=\frac12x_1x_0-y_1 \\ y_0=\frac12x_2x_0-y_2 \end{cases}$，故$l_{AB}: y_0=\frac12x_0x-y$，即$y=\frac12x_0x-y_0$。联立$\begin{cases} y=\frac12x_0x-y_0 \\ x^2=4y \end{cases}$，得$x^2-2x_0x+4y_0=0$，$\Delta=4x_0^2-16y_0$。所以$|AB|=\sqrt{1+\frac{x_0^2}{4}}\cdot\sqrt{4x_0^2-16y_0}=\sqrt{4+x_0^2}\cdot\sqrt{x_0^2-4y_0}$，$d_{P\to AB}=\frac{|x_0^2-4y_0|}{\sqrt{x_0^2+4}}$，所以$S_{\triangle PAB}=\frac12|AB|\cdot d_{P\to AB}=\frac12\sqrt{x_0^2-4y_0}\cdot\sqrt{x_0^2-4y_0}=\frac12(x_0^2-4y_0)^{\frac32}=\frac12(-y_0^2-12y_0-15)^{\frac32}$。而$y_0\in[-5,-3]$，故当$y_0=-5$时，$S_{\triangle PAB}$达到最大，最大值为$20\sqrt{5}$。

\paragraph{2021 年高考数学 天津卷第18题}

已知椭圆 $\frac{x^2}{a^2} + \frac{y^2}{b^2} = 1$ ($a > b > 0$) 的右焦点为 $F$，上顶点为 $B$，离心率为 $\frac{2\sqrt{5}}{5}$，且 $|BF| = \sqrt{5}$．
（I）求椭圆的方程；
（II）直线 $l$ 与椭圆有唯一的公共点 $M$，与 $y$ 轴的正半轴交于点 $N$，过 $N$ 与 $BF$ 垂直的直线交 $x$ 轴于点 $P$．若 $MP // BF$，求直线 $l$ 的方程．

\textbf{答案：}（I）$\frac{x^2}{5} + y^2 = 1$；（II）$x - y + \sqrt{6} = 0$\par

\textbf{解答：}（I）$F(c,0)$，$B(0,b)$，$|BF| = \sqrt{c^2+b^2} = a = \sqrt{5}$，离心率 $e = \frac{c}{a} = \frac{2\sqrt{5}}{5}$，得 $c=2$，$b = \sqrt{a^2-c^2} = 1$，所以椭圆方程为 $\frac{x^2}{5} + y^2 = 1$；
（II）设 $M(x_0,y_0)$ 为椭圆上一点，则椭圆在 $M$ 处的切线方程为 $\frac{x_0 x}{5} + y_0 y = 1$，令 $x=0$ 得 $N\left(0,\frac{1}{y_0}\right)$，$BF$ 斜率为 $k_{BF} = -\frac{b}{c} = -\frac{1}{2}$，则 $PN$ 方程为 $y = 2x + \frac{1}{y_0}$，令 $y=0$ 得 $P\left(-\frac{1}{2y_0},0\right)$，$k_{MP} = \frac{y_0}{x_0+\frac{1}{2y_0}} = \frac{2y_0^2}{2x_0 y_0+1}$，由 $MP // BF$ 得 $\frac{2y_0^2}{2x_0 y_0+1} = -\frac{1}{2}$，整理得 $(x_0+5y_0)^2=0$，所以 $x_0 = -5y_0$，代入椭圆方程得 $\frac{25y_0^2}{5}+y_0^2=6y_0^2=1$，$y_0 = \frac{\sqrt{6}}{6}$，$x_0 = -\frac{5\sqrt{6}}{6}$，所以直线 $l$ 方程为 $\frac{-\frac{5\sqrt{6}}{6}x}{5} + \frac{\sqrt{6}}{6}y = 1$，即 $x - y + \sqrt{6}=0$.

\paragraph{2020 年高考数学 全国I卷-理科第20题}

已知$A,B$分别为椭圆$E:\frac{x^2}{a^2}+y^2=1(a>1)$的左、右顶点，$G$为$E$的上顶点，$\overrightarrow{AG}\cdot\overrightarrow{GB}=8$。$P$为直线$x=6$上的动点，$PA$与$E$的另一交点为$C$，$PB$与$E$的另一交点为$D$。
(1)求$E$的方程；
(2)证明：直线$CD$过定点。

\textbf{答案：}(1)$\frac{x^2}{9}+y^2=1$；(2)直线$CD$过定点$\left(\frac{3}{2},0\right)$\par

\textbf{解答：}(1)由椭圆方程得$A(-a,0),B(a,0),G(0,1)$，则$\overrightarrow{AG}=(a,1),\overrightarrow{GB}=(a,-1)$，$\overrightarrow{AG}\cdot\overrightarrow{GB}=a^2-1=8$，解得$a^2=9$，故椭圆方程为$\frac{x^2}{9}+y^2=1$。
(2)设$P(6,y_0)$，则直线$AP$方程为$y=\frac{y_0}{9}(x+3)$，与椭圆联立得$(y_0^2+9)x^2+6y_0^2x+9y_0^2-81=0$，由韦达定理$x_A\cdot x_C=\frac{9y_0^2-81}{y_0^2+9}$，而$x_A=-3$，得$x_C=\frac{-3y_0^2+27}{y_0^2+9}$，代入得$y_C=\frac{6y_0}{y_0^2+9}$，故$C\left(\frac{-3y_0^2+27}{y_0^2+9},\frac{6y_0}{y_0^2+9}\right)$。同理直线$PB$方程$y=\frac{y_0}{3}(x-3)$，与椭圆联立得$D\left(\frac{3y_0^2-3}{y_0^2+1},\frac{-2y_0}{y_0^2+1}\right)$。当$y_0\neq0$时，直线$CD$斜率$k=\frac{\frac{6y_0}{y_0^2+9}-\frac{-2y_0}{y_0^2+1}}{\frac{-3y_0^2+27}{y_0^2+9}-\frac{3y_0^2-3}{y_0^2+1}}=\frac{4y_0}{3(3-y_0^2)}$，直线$CD$方程为$y-\frac{-2y_0}{y_0^2+1}=\frac{4y_0}{3(3-y_0^2)}\left(x-\frac{3y_0^2-3}{y_0^2+1}\right)$，化简得$y=\frac{4y_0}{3(3-y_0^2)}\left(x-\frac{3}{2}\right)$，故过定点$\left(\frac{3}{2},0\right)$。当$y_0=0$时，$C(3,0),D(3,0)$，直线$CD$为$x=3$，也过$\left(\frac{3}{2},0\right)$。综上，直线$CD$过定点$\left(\frac{3}{2},0\right)$。

\subsection*{课后作业}

\begin{enumerate}

\item 2017 年高考数学 全国I卷-理科第20题

\item 2019 年高考数学 全国III卷-文科第21题

\item 2020 年高考数学 全国I卷-理科第20题

\item 2021 年高考数学 全国乙卷-理科第21题

\end{enumerate}

\noindent\textbf{作业答案索引：}

\begin{itemize}

\item 2017 年高考数学 全国I卷-理科第20题：（1）$\frac{x^2}{4}+y^2=1$；（2）证明见解析，定点为 $(2,-1)$

\item 2019 年高考数学 全国III卷-文科第21题：(1)见详解; (2) $x^2+(y-\frac{5}{2})^2=4$或$x^2+(y-\frac{5}{2})^2=2$

\item 2020 年高考数学 全国I卷-理科第20题：(1)$\frac{x^2}{9}+y^2=1$；(2)直线$CD$过定点$\left(\frac{3}{2},0\right)$

\item 2021 年高考数学 全国乙卷-理科第21题：（1）$p=2$；（2）$20\sqrt{5}$

\end{itemize}

\subsection*{板书设计}

一、核心公式
1. 椭圆：若 $P(x_0,y_0)\in \dfrac{x^2}{a^2}+\dfrac{y^2}{b^2}=1$，则切线
\[
\frac{x_0x}{a^2}+\frac{y_0y}{b^2}=1.
\]
2. 抛物线：若 $P(x_0,y_0)\in y^2=2px$，则切线
\[
y_0y=p(x+x_0).
\]
若 $P(x_0,y_0)\in x^2=2py$，则切线
\[
xx_0=p(y+y_0).
\]
3. 切点弦：外点 $P(x_0,y_0)$ 向曲线引两条切线，切点连线常由同式直接给出。

二、结构模型
1. 手电筒模型：定点 $P$，动弦 $AB$，若 $k_{PA}+k_{PB}$ 或 $k_{PA}k_{PB}$ 为定值
\[
\Rightarrow \text{优先预判 }AB\text{ 过定点或定斜率}.
\]
证明流程：设直线 $\to$ 联立曲线 $\to$ 韦达定理 $\to$ 化弦方程。

2. 阿基米德三角形：弦 $AB$ 与两端点切线交于 $Q$
\[
S_{\triangle ABQ}=\frac12|AB|\cdot d(Q,AB).
\]

三、典型结论示意
1. $\texttt{2017-全国I卷-理科::20}$：斜率和定值 $\Rightarrow l$ 过定点。
2. $\texttt{2019-全国III卷-理科::21}$：外点引切线 $\Rightarrow AB$ 为切点弦，过定点。
3. $\texttt{2021-全国乙卷-理科::21}$：切线交点与弦构成阿基米德三角形，面积最值走“底高”。

四、警示
\[
\text{唯一公共点 }\not\Rightarrow \text{ 必为切线}.
\]

\subsection*{易错点}

\begin{itemize}

\item 把“射影中的极线”当作高考可直接引用的证明工具。纠正：只能用于结构预判，正式书写必须回到坐标法、定义法、韦达定理、点差法、向量或平面几何。

\item 已知切点仍然先设斜率、联立、判别式为零。纠正：优先用点式切线，减少运算层级。

\item 看到外点引两条切线，却分别求两条切线方程，不先想切点弦。纠正：切点连线常是最省力的中介量。

\item 把“斜率和为定值”机械背成某个现成公式。纠正：先识别模型，再设直线、联立、用韦达推导，避免错套。

\item 把“唯一公共点”等同于“相切”。纠正：必须检查是否真满足切线条件，尤其是在双曲线、抛物线和特殊方向直线中。

\item 阿基米德三角形问题中只盯两条切线，不把 $\triangle ABQ$ 作为整体。纠正：面积、垂直、轨迹常从整体关系切入更快。

\end{itemize}

\subsection*{教师备注}

1. 两课时建议各 $45$ 分钟，总时长 $90$ 分钟。第一课时以“切线方程、切点弦、手电筒模型”为主，第二课时以“阿基米德三角形与综合辨析”为主。
2. 对技巧 $47$ 的处理要把握尺度：允许教师说“这在更高观点下与极线有关”，但板书证明不得使用射影语言。
3. 对技巧 $48$ 的处理要突出“整体图形意识”，不要把它讲成单纯名词记忆。阿基米德三角形的课堂价值主要是组织解题，而非增加定义负担。
4. 若学生基础偏弱，$\texttt{2021-天津卷::18}$ 只作逻辑纠偏，不要求完整展开。
5. 若学生基础较强，可补充比较：$\texttt{2017-全国I卷-理科::20}$ 与 $\texttt{2020-全国I卷-理科::20}$ 同属“弦过定点”主线，但前者显性给出斜率和条件，后者需要从构型中自行提炼。
6. 作业分层建议：基础层完成 $\texttt{2019-全国III卷-文科::21}$ 与 $\texttt{2017-全国I卷-理科::20}$；提高层完成 $\texttt{2020-全国I卷-理科::20}$ 与 $\texttt{2021-全国乙卷-理科::21}$。
7. 课堂追问要始终围绕“先辨型，再选工具”：是点式切线，还是切点弦，还是手电筒模型，还是阿基米德三角形。这样能把零散题型压缩为少数稳定入口。

\clearpage

\section{第9课：渐近线、定点定值与非对称处理}

\noindent 建议课时：2课时。

\subsection*{教学目标}

\begin{itemize}

\item 理解并会运用双曲线渐近线斜率与离心率的关系：对横轴型双曲线 $\dfrac{x^2}{a^2}-\dfrac{y^2}{b^2}=1$，会由 $\tan\theta=\dfrac{b}{a}$ 推得 $e=\sqrt{1+\tan^2\theta}=\dfrac{1}{\cos\theta}$，也会由 $e$ 反求渐近线方程 $y=\pm\sqrt{e^2-1}\,x$。

\item 掌握“超级韦达”处理定点、斜率和、斜率积问题的基本流程：设过定点外的动直线，转化为关于 $k=\dfrac{y-y_0}{x-x_0}$ 的一元二次方程，再用韦达定理处理 $k_1+k_2,\ k_1k_2$。

\item 能识别圆锥曲线题中的非对称结构，理解何时不能直接硬套对称式，能通过配凑、半代换、保留一边消元等方式，把非对称式转化为可由韦达定理控制的形式。

\item 在解题中形成“先结构判断，再选择工具”的复习意识：选择坐标法、定义法、韦达定理、点差法、向量或平面几何进行正式作答。

\end{itemize}

\subsection*{重点与难点}

\begin{itemize}

\item 双曲线渐近线倾角与离心率的互化：$\tan\theta \leftrightarrow \dfrac{b}{a} \leftrightarrow e$。

\item 钝角倾斜角的处理：题给倾斜角若大于 $90^\circ$，要先转化到对应正斜率渐近线的锐角倾斜角。

\item “超级韦达”的前提条件：动直线不过所选定点，且两条连线斜率存在。

\item 定点问题常见判据：把结论写成一般直线 $ux+vy+w=0$ 或参数方程后，消去参数，得到与参数无关的点或直线。

\item 非对称结构的识别：目标式若不是 $x_1+x_2,\ x_1x_2,\ x_1^2+x_2^2$ 一类对称式，就要主动改造，不能机械联立后分别代入。

\end{itemize}

\subsection*{高观点}

从统一视角看，圆锥曲线与一束直线相交时，常出现“二次对象与一次对象的对应”。射影视角只用于帮助学生预判结构：一条动直线与圆锥曲线交于两点，围绕某个固定点作连线后，往往会出现关于两条斜率 $k_1,k_2$ 的二次关系；若图形中出现两端点角色不对称，则结果通常不宜直接用对称式处理。这种视角仅用于课堂中的结构识别与方法选择，不作为正式高考书写依据。

\subsection*{高中落地}

正式解题时回到高中课标方法：双曲线渐近线问题用坐标法与三角函数关系；定点问题用联立、设点、设斜率、韦达定理或向量点积；非对称问题用联立后的根与系数关系、配凑、半代换、点差法或平面几何比例关系完成推导。

\subsection*{教学流程}

\paragraph{8分钟：导入与目标建立}

教师展示本课关键词：$\text{渐近线倾角}\to\text{离心率}$，$\text{斜率二次方程}\to\text{定点}$，$\text{非对称式}\to\text{结构改造}$。先口头回顾：双曲线中 $a,b,c,e$ 的关系，椭圆或双曲线与直线联立后为什么自然出现韦达定理。

\noindent\textbf{教师追问：}双曲线 $\dfrac{x^2}{a^2}-\dfrac{y^2}{b^2}=1$ 的渐近线方程是什么？\\

\textbf{预期回答：}$y=\pm\dfrac{b}{a}x$。\par

\noindent\textbf{教师追问：}离心率与 $a,b$ 有什么关系？\\

\textbf{预期回答：}$e=\dfrac{c}{a}=\sqrt{1+\dfrac{b^2}{a^2}}$。\par

\noindent\textbf{教师追问：}若题目最后问“过定点”，你们通常先盯什么量？\\

\textbf{预期回答：}先盯参数是否能消掉；常从直线设法、交点坐标关系、斜率和斜率积入手。\par

\paragraph{18分钟：模块一：渐近线倾角与离心率}

以题目 \texttt{2019-全国I卷-文科::10}、\texttt{2020-全国III卷-文科::14}、\texttt{2021-新高考II卷::13} 串联训练。先归纳：对横轴型双曲线，正斜率渐近线倾斜角设为 $\theta\in\left(0,\dfrac{\pi}{2}\right)$，则
\[
\tan\theta=\frac{b}{a},\qquad e=\sqrt{1+\tan^2\theta}=\frac{1}{\cos\theta}.
\]
再强调：若已给某条渐近线倾斜角为钝角，例如 $130^\circ$，其斜率为负，对应另一条正斜率渐近线的倾斜角为 $50^\circ$，故应取 $\theta=50^\circ$，再代入 $e=\dfrac{1}{\cos50^\circ}$。若已知渐近线 $y=\sqrt2x$，则 $\dfrac{b}{a}=\sqrt2$，从而
\[
e=\sqrt{1+2}=\sqrt3.
\]
若已知 $e=2$，则
\[
\frac{b}{a}=\sqrt{e^2-1}=\sqrt3,
\]
故渐近线为 $y=\pm\sqrt3x$。

\noindent\textbf{教师追问：}为什么题目给出 $130^\circ$ 时不能直接把 $\cos130^\circ$ 代入 $e=\dfrac{1}{\cos\theta}$？\\

\textbf{预期回答：}因为公式中的 $\theta$ 指正斜率渐近线对应的锐角倾斜角；直接代入会出现负值，不符合离心率大于 $1$。\par

\noindent\textbf{教师追问：}如果渐近线方程直接给成 $y=mx$，最短路径是什么？\\

\textbf{预期回答：}先读出 $m=\dfrac{b}{a}$，再算 $e=\sqrt{1+m^2}$。\par

\noindent\textbf{教师追问：}若已知离心率，怎样最快写出渐近线？\\

\textbf{预期回答：}先求 $\dfrac{b}{a}=\sqrt{e^2-1}$，再写 $y=\pm\dfrac{b}{a}x$。\par

\paragraph{14分钟：小结与当堂快练}

教师组织口答快练，不展开新题，只围绕三道题的结构变式：已知斜率、已知倾角、已知离心率三种入口如何互相转换。要求学生用统一链条表示：
\[
\text{倾角}\ \Longleftrightarrow\ \text{斜率}\ \Longleftrightarrow\ \frac{b}{a}\ \Longleftrightarrow\ e.
\]
强调这是本课第一个“速判型”工具。

\noindent\textbf{教师追问：}这类选择填空题最容易丢分的步骤是什么？\\

\textbf{预期回答：}倾斜角与斜率的对应处理，尤其是钝角转锐角。\par

\noindent\textbf{教师追问：}如果双曲线换成纵轴型，哪里要调整？\\

\textbf{预期回答：}渐近线斜率应改为 $\pm\dfrac{a}{b}$，不能仍写成 $\pm\dfrac{b}{a}$。\par

\paragraph{5分钟：第一课时收束}

教师用一句话收束第一课时：凡是渐近线、倾角、离心率互求，核心都不是记答案，而是抓住 $\dfrac{b}{a}$ 这一中介量。布置学生在草稿纸上补写三题的最简规范表达。

\noindent\textbf{教师追问：}本节最关键的中介量是什么？\\

\textbf{预期回答：}$\dfrac{b}{a}$。\par

\paragraph{7分钟：第二课时导入：从斜率和到定点}

回顾椭圆、双曲线、抛物线与动直线相交时的典型问题：若题中出现“过某固定点的两条连线的斜率和、斜率积、垂直、点积恒成立”，往往可设两条连线斜率分别为 $k_1,k_2$，转化成二次方程根的关系。

\noindent\textbf{教师追问：}为什么很多定点题不直接设交点坐标，而要优先设斜率？\\

\textbf{预期回答：}因为目标常是 $k_1+k_2$ 或 $k_1k_2$，设斜率更贴近结论，韦达更直接。\par

\noindent\textbf{教师追问：}当题中出现 $PA\perp PB$ 时，斜率语言如何落地？\\

\textbf{预期回答：}若两条斜率都存在，则有 $k_1k_2=-1$。\par

\paragraph{20分钟：模块二：超级韦达与定点}

以 \texttt{2017-全国I卷-理科::20} 为主例讲解，以 \texttt{2018-全国III卷-理科::16}、\texttt{2024-天津卷::18} 作结构呼应。
对 \texttt{2017-全国I卷-理科::20} 的第 $(2)$ 问，课堂重点不在重复完整解析，而在提炼流程：
(1) 取定点 $P_2(0,1)$；
(2) 设不过 $P_2$ 的动直线 $l$ 与椭圆交于 $A,B$；
(3) 设 $P_2A,P_2B$ 的斜率为 $k_1,k_2$；
(4) 通过把椭圆方程改写成关于 $(x-0),(y-1)$ 的形式，并结合直线方程，得到关于 $k$ 的一元二次方程；
(5) 由已知 $k_1+k_2=-1$，推出动直线参数满足线性关系，从而证明 $l$ 过定点 $(2,-1)$。
教师板演时突出“先得到二次方程，再用韦达，不先忙着求点坐标”。
对 \texttt{2018-全国III卷-理科::16}，指出其本质同样是斜率积结构：$\angle AMB=90^\circ$ 对应两条连线斜率乘积为 $-1$。对 \texttt{2024-天津卷::18}，说明向量点积条件
\[
\overrightarrow{TP}\cdot\overrightarrow{TQ}\le 0
\]
也常可转化为关于两根和、积的不等式，这是“超级韦达”的延伸使用。

\noindent\textbf{教师追问：}在 \texttt{2017-全国I卷-理科::20} 中，为什么要选 $P_2(0,1)$ 作为斜率参照点？\\

\textbf{预期回答：}因为题设直接研究 $P_2A$ 与 $P_2B$ 的斜率和，选这个点最贴近条件。\par

\noindent\textbf{教师追问：}“超级韦达”和普通韦达最大的区别是什么？\\

\textbf{预期回答：}普通韦达常用于交点横纵坐标；这里是把几何中的两条连线斜率构造成二次方程的两根。\par

\noindent\textbf{教师追问：}\texttt{2018-全国III卷-理科::16} 中，直角条件最自然的代数翻译是什么？\\

\textbf{预期回答：}若两条斜率存在，则斜率积为 $-1$。\par

\paragraph{18分钟：模块三：非对称处理}

以 \texttt{2022-全国乙卷-理科::20}、\texttt{2023-新高考II卷::21} 为主，说明“非对称”并不等于“复杂计算”，而是提醒我们不能只盯 $x_1+x_2,\ x_1x_2$。
在 \texttt{2022-全国乙卷-理科::20} 中，点 $H$ 由
\[
\overrightarrow{MT}=\overrightarrow{TH}
\]
定义，说明 $T$ 是 $MH$ 的中点，故 $H$ 的坐标表达天然带有“由 $M$ 出发、沿某方向再延长”的不对称性；再与 $N$ 连线时，$M,N$ 扮演的角色并不对等。这类题如果一开始就试图把 $M,N$ 完全对称化，往往走远路。较好的思路是：联立得到 $M,N$ 的根与系数关系，同时只对涉及 $M$ 的部分先做代换，保留涉及 $N$ 的部分，最后在线方程中整体消参，得到定点 $(0,-2)$。
在 \texttt{2023-新高考II卷::21} 中，$P$ 是 $MA_1$ 与 $NA_2$ 的交点，由于 $A_1,A_2$ 分居左右，式子天然带有左右不对称。教学中强调：
\[
\text{非对称目标} \Rightarrow \text{先找“谁和谁不对等”} \Rightarrow \text{按角色分步代换}\text{。}
\]
这比一味追求把所有字母都写成根与系数的对称式更有效。
若学生追问“如何落笔”，教师统一回答：仍然是坐标法、直线方程、联立消元、韦达定理，只是代换顺序要有策略。

\noindent\textbf{教师追问：}什么样的目标式可以判断为非对称结构？\\

\textbf{预期回答：}当两个交点前的系数、位置或参与方式不同，如只出现一个点的截距、某一边被中点或倍长关系特殊处理时，往往是非对称结构。\par

\noindent\textbf{教师追问：}在 \texttt{2022-全国乙卷-理科::20} 中，为什么 $M,N$ 不宜一上来完全等价处理？\\

\textbf{预期回答：}因为点 $H$ 的构造直接依赖于 $M$ 和点 $T$，而 $N$ 只在最后连线中出现，角色不同。\par

\noindent\textbf{教师追问：}遇到非对称题时，第一步最重要的判断是什么？\\

\textbf{预期回答：}先看对象之间是否对等，决定是整体对称处理还是局部代换、半代换。\par

\paragraph{10分钟：综合归纳与方法分流}

教师将本课三类问题整理成判断流程：
\[
\text{见渐近线}\Rightarrow\frac{b}{a}\Rightarrow e;
\qquad
\text{见斜率和/积}\Rightarrow k_1,k_2\text{二次方程}\Rightarrow\text{韦达};
\qquad
\text{见角色不对等}\Rightarrow\text{非对称处理}.
\]
进一步提醒：正式高考书写中，不写“射影”字样，不写结构猜测过程，只写课标内合法推导。

\noindent\textbf{教师追问：}本课三种题型的首要识别信号分别是什么？\\

\textbf{预期回答：}渐近线题看 $\dfrac{b}{a}$；定点题看斜率和、斜率积或点积；非对称题看对象角色是否对等。\par

\noindent\textbf{教师追问：}如果同一道题同时出现定点与非对称，该先做哪一层判断？\\

\textbf{预期回答：}先判断结构是否对称，再决定用普通韦达还是非对称代换。\par

\subsection*{课堂真题}

\paragraph{2019 年高考数学 全国I卷-文科第10题}

双曲线 $C: \frac{x^2}{a^2} - \frac{y^2}{b^2} = 1$ ($a>0,b>0$) 的一条渐近线的倾斜角为 130°，则 $C$ 的离心率为

\textbf{答案：}D\par

\textbf{解答：}由已知可得 $-\frac{b}{a} = \tan130°$，$\Rightarrow \frac{b}{a} = \tan50°$，$\Rightarrow e = \frac{c}{a} = \sqrt{1+\left(\frac{b}{a}\right)^2} = \sqrt{1+\tan^2 50°} = \sqrt{1+\frac{\sin^2 50°}{\cos^2 50°}} = \sqrt{\frac{\sin^2 50°+\cos^2 50°}{\cos^2 50°}} = \frac{1}{\cos50°}$，故选 D。

\paragraph{2020 年高考数学 全国III卷-文科第14题}

设双曲线 $C:\frac{x^2}{a^2}-\frac{y^2}{b^2}=1\ (a>0,b>0)$ 的一条渐近线为 $y=\sqrt{2}x$，则 $C$ 的离心率为

\textbf{答案：}$\sqrt{3}$\par

\textbf{解答：}由双曲线方程可得其焦点在 $x$ 轴上，因为其一条渐近线为 $y=\sqrt{2}x$，所以 $\frac{b}{a}=\sqrt{2}$，$e=\frac{c}{a}=\sqrt{1+\frac{b^2}{a^2}}=\sqrt{3}$.

\paragraph{2021 年高考数学 新高考II卷第13题}

已知双曲线$\frac{x^2}{a^2}-\frac{y^2}{b^2}=1(a>0,b>0)$的离心率为$2$，则该双曲线的渐近线方程为\underline{\hspace{3cm}}

\textbf{答案：}$y=\pm\sqrt{3}x$\par

\textbf{解答：}因为双曲线$\frac{x^2}{a^2}-\frac{y^2}{b^2}=1(a>0,b>0)$的离心率为$2$，所以$e^2=\frac{c^2}{a^2}=\frac{a^2+b^2}{a^2}=2$，所以$\frac{b^2}{a^2}=3$，所以该双曲线的渐近线方程为$y=\pm\frac{b}{a}x=\pm\sqrt{3}x$.故答案为：$y=\pm\sqrt{3}x$.

\paragraph{2017 年高考数学 全国I卷-理科第20题}

已知椭圆 $C:\frac{x^2}{a^2}+\frac{y^2}{b^2}=1$（$a>b>0$），四点 $P_1(1,1)$，$P_2(0,1)$，$P_3(-1,\frac{\sqrt{3}}{2})$，$P_4(1,\frac{\sqrt{3}}{2})$ 中恰有三点在椭圆 $C$ 上．
（1）求 $C$ 的方程；
（2）设直线 $l$ 不经过 $P_2$ 点且与 $C$ 相交于 $A$，$B$ 两点．若直线 $P_2A$ 与直线 $P_2B$ 的斜率的和为 $-1$，证明：$l$ 过定点．

\textbf{答案：}（1）$\frac{x^2}{4}+y^2=1$；（2）证明见解析，定点为 $(2,-1)$\par

\textbf{解答：}解：（1）根据椭圆的对称性，$P_3(-1,\frac{\sqrt{3}}{2})$，$P_4(1,\frac{\sqrt{3}}{2})$ 两点必在椭圆 $C$ 上，又 $P_4$ 的横坐标为 1，\ensuremath{\therefore}椭圆必不过 $P_1(1,1)$，\ensuremath{\therefore}$P_2(0,1)$，$P_3(-1,\frac{\sqrt{3}}{2})$，$P_4(1,\frac{\sqrt{3}}{2})$ 三点在椭圆 $C$ 上．把 $P_2(0,1)$，$P_3(-1,\frac{\sqrt{3}}{2})$ 代入椭圆 $C$，得：$\begin{cases} \frac{0}{a^2}+\frac{1}{b^2}=1 \\ \frac{1}{a^2}+\frac{3}{4b^2}=1 \end{cases}$，解得 $a^2=4$，$b^2=1$，\ensuremath{\therefore}椭圆 $C$ 的方程为 $\frac{x^2}{4}+y^2=1$．
证明：（2）(1)当斜率不存在时，设 $l:x=m$，$A(m,y_A)$，$B(m,-y_A)$，\ensuremath{\because}直线 $P_2A$ 与直线 $P_2B$ 的斜率的和为 $-1$，\ensuremath{\therefore}$\frac{y_A-1}{m-0}+\frac{-y_A-1}{m-0}=\frac{-2}{m}=-1$，解得 $m=2$，此时 $l$ 过椭圆右顶点，不存在两个交点，故不满足．(2)当斜率存在时，设 $l:y=kx+t$，($t\neq1$)，$A(x_1,y_1)$，$B(x_2,y_2)$，联立 $\begin{cases} y=kx+t \\ \frac{x^2}{4}+y^2=1 \end{cases}$，整理，得 $(1+4k^2)x^2+8ktx+4t^2-4=0$，$\Delta>0$，$x_1+x_2=-\frac{8kt}{1+4k^2}$，$x_1x_2=\frac{4t^2-4}{1+4k^2}$，则 $k_{P_2A}+k_{P_2B}=\frac{y_1-1}{x_1}+\frac{y_2-1}{x_2}=\frac{kx_1+t-1}{x_1}+\frac{kx_2+t-1}{x_2}=2k+\frac{t-1}{x_1}+\frac{t-1}{x_2}=2k+(t-1)\frac{x_1+x_2}{x_1x_2}=2k+(t-1)\frac{-\frac{8kt}{1+4k^2}}{\frac{4t^2-4}{1+4k^2}}=2k-\frac{2k(t-1)}{t+1}=\frac{2k(t+1)-2k(t-1)}{t+1}=\frac{4k}{t+1}=-1$，又 $t\neq1$，\ensuremath{\therefore}$t=-2k-1$，此时 $\Delta=64k^2-16(1+4k^2)(t^2-1)$，代入得 $\Delta=-64k$，存在 $k$，使得 $\Delta>0$ 成立，\ensuremath{\therefore}直线 $l$ 的方程为 $y=kx-2k-1$，当 $x=2$ 时，$y=-1$，\ensuremath{\therefore}$l$ 过定点 $(2,-1)$．

\paragraph{2018 年高考数学 全国III卷-理科第16题}

已知点$M(-1,1)$和抛物线$C:y^2=4x$，过$C$的焦点且斜率为$k$的直线与$C$交于$A$，$B$两点．若$\angle AMB=90^\circ$，则$k=$\underline{\hspace{3cm}}．

\textbf{答案：}$2$\par

\textbf{解答：}焦点$F(1,0)$，设直线$AB:y=k(x-1)$，联立$\begin{cases}y=k(x-1)\\y^2=4x\end{cases}$得$k^2x^2-2(2+k^2)x+k^2=0$，设$A(x_1,y_1)$，$B(x_2,y_2)$，则$x_1+x_2=\frac{2(2+k^2)}{k^2}$，$x_1x_2=1$，$y_1+y_2=k(x_1+x_2-2)=\frac{4}{k}$，$y_1y_2=k^2(x_1-1)(x_2-1)=k^2[x_1x_2-(x_1+x_2)+1]=-4$，$\because \angle AMB=90^\circ$，$\overrightarrow{MA}\cdot\overrightarrow{MB}=0$，即$(x_1+1)(x_2+1)+(y_1-1)(y_2-1)=0$，整理得$x_1x_2+(x_1+x_2)+y_1y_2-(y_1+y_2)+2=0$，代入得$1+\frac{2(2+k^2)}{k^2}-4-\frac{4}{k}+2=0$，解得$k=2$．

\paragraph{2024 年高考数学 天津卷第18题}

已知椭圆 $\frac{x^2}{a^2}+\frac{y^2}{b^2}=1(a>b>0)$ 的离心率 $e=\frac{1}{2}$. 左顶点为 $A$, 下顶点为 $B$, $C$ 是线段 $OB$ 的中点, 其中 $S_{\triangle ABC}=\frac{3\sqrt{3}}{2}$.
(1) 求椭圆方程.
(2) 过点 $(0,-\frac{3}{2})$ 的动直线与椭圆有两个交点 $P$, $Q$. 在 $y$ 轴上是否存在点 $T$ 使得 $\overrightarrow{TP}\cdot\overrightarrow{TQ}\leq 0$ 恒成立. 若存在求出这个 $T$ 点纵坐标的取值范围, 若不存在请说明理由.

\textbf{答案：}(1) $\frac{x^2}{12}+\frac{y^2}{9}=1$; (2) 存在 $T(0,t)$ 且 $-3\leq t\leq \frac{3}{2}$, 使得 $\overrightarrow{TP}\cdot\overrightarrow{TQ}\leq 0$ 恒成立.\par

\textbf{解答：}(1) 因为 $e=\frac{1}{2}$, 故 $a=2c$, $b=\sqrt{3}c$. 所以 $A(-2c,0)$, $B(0,-\sqrt{3}c)$, $C(0,-\frac{\sqrt{3}}{2}c)$. $S_{\triangle ABC}=\frac{1}{2}\times 2c\times \frac{\sqrt{3}}{2}c=\frac{\sqrt{3}}{2}c^2=\frac{3\sqrt{3}}{2}$, 解得 $c=\sqrt{3}$, 所以 $a=2\sqrt{3}$, $b=3$, 故椭圆方程为 $\frac{x^2}{12}+\frac{y^2}{9}=1$.
(2) 设过点 $(0,-\frac{3}{2})$ 的动直线方程为 $y=kx-\frac{3}{2}$, $P(x_1,y_1)$, $Q(x_2,y_2)$, $T(0,t)$. 联立 $\begin{cases} 3x^2+4y^2=36 \\ y=kx-\frac{3}{2} \end{cases}$, 得 $(3+4k^2)x^2-12kx-27=0$, $\Delta>0$, $x_1+x_2=\frac{12k}{3+4k^2}$, $x_1x_2=-\frac{27}{3+4k^2}$. $\overrightarrow{TP}=(x_1,y_1-t)$, $\overrightarrow{TQ}=(x_2,y_2-t)$. $\overrightarrow{TP}\cdot\overrightarrow{TQ}=x_1x_2+(y_1-t)(y_2-t)=x_1x_2+(kx_1-\frac{3}{2}-t)(kx_2-\frac{3}{2}-t)=(1+k^2)x_1x_2-k(\frac{3}{2}+t)(x_1+x_2)+(\frac{3}{2}+t)^2$. 代入韦达定理, 得 $\frac{[(3+2t)^2-12t-45]k^2+3(\frac{3}{2}+t)^2-27}{3+4k^2}$. 要使 $\overrightarrow{TP}\cdot\overrightarrow{TQ}\leq 0$ 恒成立, 需 $\begin{cases} (3+2t)^2-12t-45\leq 0 \\ 3(\frac{3}{2}+t)^2-27\leq 0 \end{cases}$, 解得 $-3\leq t\leq \frac{3}{2}$. 当斜率不存在时, 直线为 $x=0$, 与椭圆交于 $(0,3)$, $(0,-3)$, 此时需 $-3\leq t\leq 3$, 结合得 $-3\leq t\leq \frac{3}{2}$. 综上, 存在 $T(0,t)$ 且 $-3\leq t\leq \frac{3}{2}$, 使得 $\overrightarrow{TP}\cdot\overrightarrow{TQ}\leq 0$ 恒成立.

\paragraph{2022 年高考数学 全国乙卷-理科第20题}

已知椭圆$E$的中心为坐标原点，对称轴为$x$轴、$y$轴，且过$A(0,-2), B\left(\frac{3}{2},-1\right)$两点．
(1)求$E$的方程；
(2)设过点$P(1,-2)$的直线交$E$于$M,N$两点，过$M$且平行于$x$轴的直线与线段$AB$交于点$T$，点$H$满足$\overrightarrow{MT}=\overrightarrow{TH}$．证明：直线$HN$过定点．

\textbf{答案：}(1) $\frac{y^2}{4}+\frac{x^2}{3}=1$
(2) $(0,-2)$\par

\textbf{解答：}【小问1详解】\\解：设椭圆E的方程为$mx^2+ny^2=1$，过$A(0,-2), B\left(\frac{3}{2},-1\right)$，\\则$\begin{cases} 4n=1 \\ \frac{9}{4}m+n=1 \end{cases}$，解得$m=\frac{1}{3}, n=\frac{1}{4}$，\\所以椭圆E的方程为：$\frac{y^2}{4}+\frac{x^2}{3}=1$.\\【小问2详解】\\$A(0,-2), B(\frac{3}{2},-1)$，所以$AB: y+2 = \frac{2}{3}x$，\\(1)若过点$P(1,-2)$的直线斜率不存在，直线$x=1$.代入$\frac{x^2}{3}+\frac{y^2}{4}=1$，\\可得$M\left(1,\frac{2\sqrt{6}}{3}\right), N\left(1,-\frac{2\sqrt{6}}{3}\right)$，代入AB方程$y=\frac{2}{3}x-2$，可得\\$T\left(\sqrt{6}+3,\frac{2\sqrt{6}}{3}\right)$，由$\overrightarrow{MT}=\overrightarrow{TH}$得到$H\left(2\sqrt{6}+5,\frac{2\sqrt{6}}{3}\right)$.求得HN方程：$y=\left(2-\frac{2\sqrt{6}}{3}\right)x-2$，过点$(0,-2)$.\\(2)若过点$P(1,-2)$的直线斜率存在，设$kx-y-(k+2)=0, M(x_1,y_1), N(x_2,y_2)$.\\联立$\begin{cases} kx-y-(k+2)=0 \\ \frac{x^2}{3}+\frac{y^2}{4}=1 \end{cases}$，得$(3k^2+4)x^2 - 6k(2+k)x + 3k(k+4)=0$，\\可得$\begin{cases} x_1+x_2 = \frac{6k(2+k)}{3k^2+4} \\ x_1x_2 = \frac{3k(4+k)}{3k^2+4} \end{cases}$，$\begin{cases} y_1+y_2 = \frac{-8(2+k)}{3k^2+4} \\ y_1y_2 = \frac{4(4+4k-2k^2)}{3k^2+4} \end{cases}$，\\且$x_1y_2+x_2y_1 = \frac{-24k}{3k^2+4}$ (*)\\联立$\begin{cases} y=y_1 \\ y=\frac{2}{3}x-2 \end{cases}$，可得$T\left(\frac{3y_1}{2}+3, y_1\right), H(3y_1+6-x_1, y_1)$.\\可求得此时$HN: y-y_2 = \frac{y_1-y_2}{3y_1+6-x_1-x_2}(x-x_2)$，\\将$(0,-2)$代入整理得$2(x_1+x_2)-6(y_1+y_2)+x_1y_2+x_2y_1-3y_1y_2-12=0$，\\将(*)代入，得$24k+12k^2+96+48k-24k-48-48k+24k^2-36k^2-48=0$，显然成立，\\综上，可得直线HN过定点$(0,-2)$.

\paragraph{2023 年高考数学 新高考II卷第21题}

已知双曲线$C$的中心为坐标原点，左焦点为$(-2\sqrt{5},0)$，离心率为$\sqrt{5}$．
（1）求$C$的方程；
（2）记$C$的左、右顶点分别为$A_{1}$，$A_{2}$，过点$(-4,0)$的直线与$C$的左支交于$M$，$N$两点，$M$在第二象限，直线$MA_{1}$与$NA_{2}$交于点$P$．证明:点$P$在定直线上.

\textbf{答案：}（1）$\dfrac{x^{2}}{4}-\dfrac{y^{2}}{16}=1$；（2）证明见解析.\par

\textbf{解答：}（1）设双曲线方程为$\dfrac{x^{2}}{a^{2}}-\dfrac{y^{2}}{b^{2}}=1(a>0,b>0)$，由焦点坐标可知$c=2\sqrt{5}$，则由$e=\dfrac{c}{a}=\sqrt{5}$可得$a=2$，$b=\sqrt{c^{2}-a^{2}}=4$，双曲线方程为$\dfrac{x^{2}}{4}-\dfrac{y^{2}}{16}=1$.
（2）由(1)可得$A_{1}(-2,0)$，$A_{2}(2,0)$，设$M(x_{1},y_{1})$，$N(x_{2},y_{2})$，显然直线的斜率不为$0$，所以设直线$MN$的方程为$x=my-4$，且$-\dfrac{1}{2}<m<\dfrac{1}{2}$，与$\dfrac{x^{2}}{4}-\dfrac{y^{2}}{16}=1$联立可得$(4m^{2}-1)y^{2}-32my+48=0$，且$\Delta=64(4m^{2}+3)>0$，则$y_{1}+y_{2}=\dfrac{32m}{4m^{2}-1}$，$y_{1}y_{2}=\dfrac{48}{4m^{2}-1}$，直线$MA_{1}$的方程为$y=\dfrac{y_{1}}{x_{1}+2}(x+2)$，直线$NA_{2}$的方程为$y=\dfrac{y_{2}}{x_{2}-2}(x-2)$，联立直线$MA_{1}$与直线$NA_{2}$的方程可得：$\dfrac{x+2}{x-2}=\dfrac{y_{2}(x_{1}+2)}{y_{1}(x_{2}-2)}=\dfrac{y_{2}(my_{1}-2)}{y_{1}(my_{2}-6)}=\dfrac{my_{1}y_{2}-2y_{1}}{my_{1}y_{2}-6y_{1}}=\dfrac{m\cdot\dfrac{48}{4m^{2}-1}-2y_{1}}{m\cdot\dfrac{48}{4m^{2}-1}-6y_{1}}=\dfrac{\dfrac{48m}{4m^{2}-1}-2y_{1}}{\dfrac{48m}{4m^{2}-1}-6y_{1}}=\dfrac{\dfrac{48m}{4m^{2}-1}-2y_{1}}{\dfrac{48m}{4m^{2}-1}-6y_{1}}$，由$y_{1}+y_{2}=\dfrac{32m}{4m^{2}-1}$可得$\dfrac{48m}{4m^{2}-1}=2\cdot\dfrac{16m}{4m^{2}-1}+\dfrac{16m}{4m^{2}-1}=\dfrac{1}{2}(y_{1}+y_{2})+\dfrac{16m}{4m^{2}-1}$，代入得$\dfrac{x+2}{x-2}=\dfrac{\dfrac{1}{2}(y_{1}+y_{2})+\dfrac{16m}{4m^{2}-1}-2y_{1}}{\dfrac{1}{2}(y_{1}+y_{2})+\dfrac{16m}{4m^{2}-1}-6y_{1}}=\dfrac{-\dfrac{3}{2}y_{1}+\dfrac{1}{2}y_{2}+\dfrac{16m}{4m^{2}-1}}{-\dfrac{11}{2}y_{1}+\dfrac{1}{2}y_{2}+\dfrac{16m}{4m^{2}-1}}$，又$y_{2}=\dfrac{32m}{4m^{2}-1}-y_{1}$，代入得$\dfrac{x+2}{x-2}=\dfrac{-\dfrac{3}{2}y_{1}+\dfrac{1}{2}(\dfrac{32m}{4m^{2}-1}-y_{1})+\dfrac{16m}{4m^{2}-1}}{-\dfrac{11}{2}y_{1}+\dfrac{1}{2}(\dfrac{32m}{4m^{2}-1}-y_{1})+\dfrac{16m}{4m^{2}-1}}=\dfrac{-2y_{1}+\dfrac{32m}{4m^{2}-1}}{-6y_{1}+\dfrac{32m}{4m^{2}-1}}=\dfrac{2(-y_{1}+\dfrac{16m}{4m^{2}-1})}{6(-y_{1}+\dfrac{16m}{4m^{2}-1})}=\dfrac{1}{3}$，即$\dfrac{x+2}{x-2}=\dfrac{1}{3}$，解得$x=-4$，故点$P$在定直线$x=-4$上.

\subsection*{课后作业}

\begin{enumerate}

\item 2019 年高考数学 全国I卷-文科第10题

\item 2020 年高考数学 全国III卷-文科第14题

\item 2021 年高考数学 新高考II卷第13题

\item 2017 年高考数学 全国I卷-理科第20题

\item 2018 年高考数学 全国III卷-理科第16题

\item 2024 年高考数学 天津卷第18题

\item 2022 年高考数学 全国乙卷-理科第20题

\item 2023 年高考数学 新高考II卷第21题

\end{enumerate}

\noindent\textbf{作业答案索引：}

\begin{itemize}

\item 2019 年高考数学 全国I卷-文科第10题：D

\item 2020 年高考数学 全国III卷-文科第14题：$\sqrt{3}$

\item 2021 年高考数学 新高考II卷第13题：$y=\pm\sqrt{3}x$

\item 2017 年高考数学 全国I卷-理科第20题：（1）$\frac{x^2}{4}+y^2=1$；（2）证明见解析，定点为 $(2,-1)$

\item 2018 年高考数学 全国III卷-理科第16题：$2$

\item 2024 年高考数学 天津卷第18题：(1) $\frac{x^2}{12}+\frac{y^2}{9}=1$; (2) 存在 $T(0,t)$ 且 $-3\leq t\leq \frac{3}{2}$, 使得 $\overrightarrow{TP}\cdot\overrightarrow{TQ}\leq 0$ 恒成立.

\item 2022 年高考数学 全国乙卷-理科第20题：(1) $\frac{y^2}{4}+\frac{x^2}{3}=1$
(2) $(0,-2)$

\item 2023 年高考数学 新高考II卷第21题：（1）$\dfrac{x^{2}}{4}-\dfrac{y^{2}}{16}=1$；（2）证明见解析.

\end{itemize}

\subsection*{板书设计}

板书一区：课题与总框架
\[
\text{第9讲：渐近线、定点定值与非对称处理}
\]
\[
\text{识别} \to \text{转化} \to \text{规范表达}
\]

板书二区：渐近线模块
\[
\frac{x^2}{a^2}-\frac{y^2}{b^2}=1\Rightarrow y=\pm\frac{b}{a}x
\]
\[
\tan\theta=\frac{b}{a},\qquad e=\sqrt{1+\frac{b^2}{a^2}}=\sqrt{1+\tan^2\theta}=\frac{1}{\cos\theta}
\]
\[
\frac{b}{a}=\sqrt{e^2-1}\Rightarrow y=\pm\sqrt{e^2-1}\,x
\]
旁注：$130^\circ\to 50^\circ$

板书三区：超级韦达通法
设定点 $P(x_0,y_0)$，连线斜率为 $k=\dfrac{y-y_0}{x-x_0}$。
\[
\text{曲线方程} + \text{动直线方程} \Rightarrow \text{关于 }k\text{ 的二次方程}
\]
\[
\text{两根 }k_1,k_2\Rightarrow k_1+k_2,\ k_1k_2\text{ 用韦达}
\]
\[
PA\perp PB\Rightarrow k_1k_2=-1
\]
\[
\overrightarrow{TP}\cdot\overrightarrow{TQ}\le 0\Rightarrow \text{化为两根和、积的不等式}
\]

板书四区：非对称处理提醒
\[
\text{非对称} \neq \text{乱算}
\]
\[
\text{先判角色是否对等\text{，}再决定代换顺序}
\]
\[
\text{常用：半代换、配凑、保留一边、整体消参}
\]

板书五区：本课结论树
\[
\text{渐近线} \Rightarrow \frac{b}{a} \Rightarrow e
\]
\[
\text{定点/定值} \Rightarrow k_1,k_2 \Rightarrow \text{韦达}
\]
\[
\text{非对称} \Rightarrow \text{结构改造} \Rightarrow \text{再消元}
\]

\subsection*{易错点}

\begin{itemize}

\item 把钝角倾斜角直接代入 $e=\dfrac{1}{\cos\theta}$，导致离心率出现负值或错误值。应先转化为正斜率渐近线的锐角倾斜角。

\item 横轴型与纵轴型双曲线混淆，把渐近线斜率误写为同一种形式。

\item 见到动直线与圆锥曲线相交，就机械先求两个交点坐标，忽视题目真正需要的是 $k_1+k_2$ 或 $k_1k_2$。

\item 使用“超级韦达”时忘记检查前提：动直线是否经过定点、两条连线斜率是否都存在。

\item 把非对称问题硬做成对称问题，导致式子越来越长，最后无法消参。

\item 在中点、倍长、交点构造类问题中，没有意识到某一个交点会被特殊对待，从而错失最自然的代换顺序。

\item 点积条件只会展开坐标乘法，不会主动联系根与系数关系，错过整体化简机会。

\end{itemize}

\subsection*{教师备注}

一、时间实施建议：第1课时完成“渐近线倾角与离心率”及快练，第2课时完成“超级韦达”“非对称处理”与综合归纳。

二、课堂执行提醒：
\[
\text{讲结构\text{，}不代替学生算完所有细节\text{。}}
\]
尤其是 \texttt{2017-全国I卷-理科::20}、\texttt{2022-全国乙卷-理科::20}、\texttt{2023-新高考II卷::21} 三题，教师应抓“为什么这样设、为什么这样消元”，而不是把板书变成长篇运算。

三、提问层次设计：
先问“看见了什么结构”，再问“为什么选这个工具”，最后问“规范书写怎么落地”。避免一上来只问答案。

四、正式高考书写规范提醒：
关于渐近线题，只写
\[
\frac{b}{a}=\tan\theta\quad\text{或}\quad e=\sqrt{1+\frac{b^2}{a^2}}
\]
等合法公式；关于定点题，只写联立、设斜率、韦达、消参过程；关于非对称题，只写必要代换与化简，不写“射影视角”“结构预判”等课堂术语。

五、分层要求：
基础层学生要求会完成 \texttt{2019-全国I卷-文科::10}、\texttt{2020-全国III卷-文科::14}、\texttt{2021-新高考II卷::13} 的互化与 \texttt{2018-全国III卷-理科::16} 的斜率积翻译；提升层学生要求能独立复述 \texttt{2017-全国I卷-理科::20} 的“超级韦达”主线，并能对 \texttt{2022-全国乙卷-理科::20}、\texttt{2023-新高考II卷::21} 说明非对称来源。

六、课后任务安排建议：
第1组：完成三道基础题 \texttt{2019-全国I卷-文科::10}、\texttt{2020-全国III卷-文科::14}、\texttt{2021-新高考II卷::13}，每题写“入口量是什么”。
第2组：完成 \texttt{2017-全国I卷-理科::20}、\texttt{2018-全国III卷-理科::16}，每题写出“斜率二次方程的两根分别代表什么”。
第3组：完成 \texttt{2024-天津卷::18}、\texttt{2022-全国乙卷-理科::20}、\texttt{2023-新高考II卷::21}，每题写出“本题为什么具有非对称或准非对称特征”。

七、评价标准：
学生若能准确说出
\[
\frac{b}{a},\quad k_1+k_2,\quad k_1k_2,\quad \text{角色是否对等}
\]
这四个关键词在对应题型中的作用，本课目标基本达成。

\clearpage

\section{第10课：高考综合题的结构识别与规范表达}

\noindent 建议课时：2课时。

\subsection*{教学目标}

\begin{itemize}

\item 能从题设中的“焦点、准线、定比、斜率积、切线、弦中点、过定点、面积最值”等信息，快速识别对应的结构类型，并完成从“高观点预判”到“高中方法求解”的转化。

\item 能将综合题的思路归入以下可执行通道：圆锥曲线第一定义、第二定义、第三定义、隐圆识别、点差法、韦达定理、切线点式、向量数量积与解析几何联立。

\item 能按照高考规范完成解答书写：设元清晰、等价变形合法、关键结论有依据、结论回扣题意、范围与特殊情形有复核。

\item 能对解答进行二次复核，重点检查：定义域、斜率是否存在、线与曲线位置关系、所证定点是否与参数无关、最值是否取到。

\end{itemize}

\subsection*{重点与难点}

\begin{itemize}

\item 结构识别先行：见到 $|MF_1|+|MF_2|$、$\bigl||MF_1|-|MF_2|\bigr|$、$|MF|=d(M,l)$，优先联想到圆锥曲线定义；见到 $k_{PA}k_{PB}$ 为定值，优先联想到第三定义或韦达处理；见到二元二次式中 $x^2,y^2$ 系数相同，优先检查是否可化为圆。

\item 正式求解必须落回高中方法：坐标法、定义法、韦达定理、点差法、向量法、平面几何；“高观点”只用于统一视角、结构预判和概念溯源。

\item 规范表达的核心是“先判型，再选法，再书写依据，再检验范围”。

\item 综合题复核的核心是“结果正确”与“表达可得分”并重。

\end{itemize}

\subsection*{高观点}

从统一视角看，许多高考综合题都在处理“点、线、二次曲线”之间的不变量：如弦过定点、切点弦过定点、对偶式的点线对应、由定比或交比痕迹引出的圆或圆锥曲线。这里可借助射影几何作结构预判：例如“切线交点的轨迹”“过定点的弦”“两束连线的斜率和或积不变”常提示存在更高层面的点线对应关系。但课堂中所有正式解答必须回到高中课标允许的方法，不直接使用射影结论作证明依据。

\subsection*{高中落地}

将高观点落地为高中方法时，可执行的翻译规则是：若预判为“隐圆”，就设点坐标、列距离关系、平方化简并配方；若预判为“第一定义/第二定义”，就把题设翻译为点到定点、点到定直线的距离关系；若预判为“第三定义或手电筒模型”，就联立直线与曲线，利用韦达定理求斜率和、斜率积或证明过定点；若题目给出弦中点，就优先用点差法；若涉及切线，就优先调用点式切线方程或切点弦方程；若涉及面积、垂直、最值，则再接向量、距离公式或基本不等式完成收尾。

\subsection*{教学流程}

\paragraph{8分钟：导入与目标明确}

教师出示本节主线：$\text{结构识别}\rightarrow\text{方法落地}\rightarrow\text{规范表达}\rightarrow\text{复核}$。强调：同一道题，先判断属于“定义型、隐圆型、斜率积型、切线型、点差型、焦点面积型”中的哪一类，再选择坐标法或几何法。

\noindent\textbf{教师追问：}综合题一上来最容易失分的环节是什么？\\

\textbf{预期回答：}没有先识别结构，直接设元硬算，导致式子很乱。\par

\noindent\textbf{教师追问：}为什么有些题会“算得动但写不稳”？\\

\textbf{预期回答：}因为方法选得不经济，或者关键依据没有写清，虽然想到结论但得分点不完整。\par

\paragraph{12分钟：结构图谱建构}

围绕参考技巧 $43\sim57$ 建立识别表：
$\bullet$ 若出现 $|MF_1|+|MF_2|=\text{常数}$、$\bigl||MF_1|-|MF_2|\bigr|=\text{常数}$、$|MF|=d(M,l)$，走定义法；
$\bullet$ 若出现 $k_{PA}k_{PB}=\text{常数}$，优先考虑第三定义或设点化简；
$\bullet$ 若二次式化为 $x^2+y^2+Dx+Ey+F=0$，识别为隐圆；
$\bullet$ 若已知弦中点，优先点差法；
$\bullet$ 若给出切点，优先点式切线；若给出曲线外一点引两切线，优先切点弦。

\noindent\textbf{教师追问：}见到“斜率之积为定值”，你第一反应应该是什么？\\

\textbf{预期回答：}先判断是否能归入第三定义，或者用设点与韦达定理处理。\par

\noindent\textbf{教师追问：}见到“化简后 $x^2,y^2$ 系数相等”，下一步最关键的动作是什么？\\

\textbf{预期回答：}立刻配方，提取圆心和半径，转入圆的问题。\par

\paragraph{20分钟：例题1 结构识别示范}

选用题目：`2017-全国II卷-理科::20`。本题第（1）问的核心不是“椭圆上动点带参数”，而是通过向量关系把点 $P$ 表示出来，再化为二元二次式，识别隐圆。第（2）问则转入向量数量积与直线垂直关系，最终回到解析几何证明定点在线上。
教师示范书写主骨架：设 $M(x,y)$，由 $M\in C$ 得 $\dfrac{x^2}{2}+y^2=1$；过 $M$ 作 $x$ 轴垂线，垂足 $N(x,0)$；由 $\overrightarrow{NP}=\sqrt{2}\overrightarrow{NM}$，可将 $P$ 坐标表示为 $\bigl(x,\sqrt{2}y\bigr)$；代回得 $x^2+y_P^2=2$，从而轨迹为圆。
第（2）问重点演示“结论不抢跑”：先设 $Q(-3,t)$，再由 $\overrightarrow{OP}\cdot\overrightarrow{PQ}=1$ 建立方程，推出 $t$ 与 $P$ 的关系，再写 $OQ$ 的斜率，最后写过 $P$ 且垂直于 $OQ$ 的直线方程，并验证左焦点 $F$ 满足之。

\noindent\textbf{教师追问：}第（1）问如果一开始就把 $M$ 参数化为三角函数，有没有必要？\\

\textbf{预期回答：}没有必要，直接坐标表示更简洁。\par

\noindent\textbf{教师追问：}为什么这一问应优先归入‘隐圆’而不是继续在椭圆背景里纠缠？\\

\textbf{预期回答：}因为变换后方程很快变成 $x^2+y^2=\text{常数}$，圆的结构已经显现。\par

\noindent\textbf{教师追问：}第（2）问证明一条直线过定点，规范书写通常落在哪一步？\\

\textbf{预期回答：}写出所求直线方程后，把定点坐标代入验证恒成立。\par

\paragraph{5分钟：小结1}

归纳：当题目背景来自椭圆、抛物线或双曲线时，不要被“原曲线身份”束缚；真正的轨迹可能在中途转化为圆。书写上必须出现“由已知可得点 $P$ 坐标……代入化简得……故轨迹方程为……”的完整链条。

\noindent\textbf{教师追问：}这一题的高观点预判是什么？高中落地方法是什么？\\

\textbf{预期回答：}预判是存在二次曲线之间的简单仿射变换痕迹；落地方法是设点坐标、向量表示、代数化简、配方识圆。\par

\paragraph{18分钟：例题2 斜率积到曲线判定}

选用题目：`2019-全国II卷-理科::21`。第（1）问用第三定义视角做结构预判：两定点 $A(-2,0),B(2,0)$ 与动点 $M$ 满足 $k_{AM}k_{BM}=-\dfrac12$，负常数，预判为以 $AB$ 为左右顶点相关的椭圆。正式解答回到高中方法：设 $M(x,y)$，则
\[
\frac{y}{x+2}\cdot\frac{y}{x-2}=-\frac12.
\]
整理得
\[
\frac{x^2}{4}+\frac{y^2}{2}=1.
\]
再说明去掉 $x=\pm2$ 对应的顶点情形。
第（2）问强调“联立+韦达+面积表达”的规范链条，不追求课堂内全部算完，而着重示范如何组织证明：过原点直线设为 $y=kx$，与椭圆联立得交点 $P,Q$；根据 $E$ 为 $P$ 到 $x$ 轴垂足构造 $QE$，再与曲线交于 $G$；证明直角时先找斜率积或向量数量积；求面积最值时再写成单变量函数。

\noindent\textbf{教师追问：}本题为什么可用第三定义作预判，但不能把第三定义本身当作完整证明？\\

\textbf{预期回答：}因为高考解答要把关系代数化并推出标准方程，预判只帮助选法。\par

\noindent\textbf{教师追问：}由斜率积为负常数，你能先预判曲线类型吗？\\

\textbf{预期回答：}可以，通常预判为椭圆。\par

\noindent\textbf{教师追问：}写出轨迹后，为什么还要补一句‘是椭圆’？\\

\textbf{预期回答：}因为题目要求不仅是方程，还要说明曲线名称，结论要完整。\par

\paragraph{22分钟：第二课时导入与规范表达训练}

回顾上节两类典型：隐圆型与斜率积型。进入本课重点：同学能做出思路，但失分多在表达。教师给出综合题书写模板：
\[
\text{判型}\rightarrow\text{设元}\rightarrow\text{列关系}\rightarrow\text{化简}\rightarrow\text{结论}\rightarrow\text{检验}.
\]
同时给出复核清单：是否说明参数范围；是否说明斜率存在；是否说明唯一交点未必是切线；是否说明定点与参数无关。

\noindent\textbf{教师追问：}高考阅卷中，‘想到但没写’最常见丢分点有哪些？\\

\textbf{预期回答：}关键关系式来历没写、曲线类型没说明、特殊情形没讨论、最后结论没回扣题意。\par

\noindent\textbf{教师追问：}为什么综合题要有‘检验与说明’这一环？\\

\textbf{预期回答：}因为有些式子变形会引入增根，或者遗漏斜率不存在、切线退化等特殊情形。\par

\paragraph{18分钟：例题3 切线与定点结构}

选用题目：`2019-全国III卷-理科::21`。本题适合训练“高观点预判为点线对应，高中落地为切线点式与消元”。对抛物线 $C:y=\dfrac{x^2}{2}$，若切点为 $(t,\dfrac{t^2}{2})$，则切线方程可由导数背景外的高中方式写成
\[
y=tx-\frac{t^2}{2}.
\]
设动点 $D(u,-\dfrac12)$，过 $D$ 作两条切线，对应切点参数为 $t_1,t_2$。由 $D$ 在切线上得
\[
-\frac12=t_i u-\frac{t_i^2}{2}
\quad (i=1,2),
\]
即
\[
t^2-2ut-1=0.
\]
故 $t_1+t_2=2u,\ t_1t_2=-1$。再由两切点弦 $AB$ 方程组织证明其过定点。
教学重点不在算快，而在示范：切线问题优先参数切点或点式切线，定点问题优先把直线写出来再代定点验证。

\noindent\textbf{教师追问：}为什么这题不能直接凭‘两条切线’就开始设两条一般直线？\\

\textbf{预期回答：}因为那样未知量太多，利用切点参数能把两条切线统一起来。\par

\noindent\textbf{教师追问：}由 $t_1+t_2$、$t_1t_2$ 已知后，证明过定点最稳的动作是什么？\\

\textbf{预期回答：}先写出弦 $AB$ 的方程，再把候选定点代入检验恒成立。\par

\paragraph{15分钟：例题4 焦点三角形与参数范围}

选用题目：`2019-全国II卷-文科::20`。本题训练“焦点三角形面积问题”的识别与规范表达。第（1）问由 $\triangle POF_2$ 为等边三角形，可把边、角关系转化为椭圆离心率关系；第（2）问由 $PF_1\perp PF_2$ 可知 $\angle F_1PF_2=90^\circ$，面积为
\[
S_{\triangle F_1PF_2}=\frac12|PF_1|\,|PF_2|.
\]
再结合椭圆中焦点三角形的常用结论或由定义、余弦定理推导，得到 $b$ 并进一步求 $a$ 的范围。
本环节强调：涉及范围题时，最终答案不能只给一个代数式，还要说明取值依据，例如由椭圆存在条件、点在曲线上、或三角函数取值范围得到区间。

\noindent\textbf{教师追问：}见到 $PF_1\perp PF_2$，你应立刻想到哪一个标准角？\\

\textbf{预期回答：}想到 $\angle F_1PF_2=90^\circ$。\par

\noindent\textbf{教师追问：}范围题为什么不能只算出一个不等式就结束？\\

\textbf{预期回答：}因为还要说明不等式来源、端点是否能取到、是否与曲线存在性条件一致。\par

\paragraph{10分钟：当堂练写与互评}

学生分组完成两项短写：
其一，围绕 `2023-北京卷::6`，用圆锥曲线第二定义写出简洁求解链；
其二，围绕 `2022-全国甲卷-理科::10`，先写“结构预判句”，再写“正式求解入口句”。
教师巡视，重点检查是否出现“先判型后落地”的表达。

\noindent\textbf{教师追问：}对于 `2023-北京卷::6`，最短的得分链是什么？\\

\textbf{预期回答：}由抛物线定义，把到焦点距离转化为到准线距离，再结合到直线 $x=-3$ 的距离求得。\par

\noindent\textbf{教师追问：}对于 `2022-全国甲卷-理科::10`，结构预判句应如何写得既准又不过界？\\

\textbf{预期回答：}可写‘由左顶点与关于 $y$ 轴对称的两点连线斜率积为定值，预判可用圆锥曲线第三定义；正式解答设点坐标并化简求离心率。’\par

\paragraph{7分钟：课堂总结}

教师带领学生完成本节总结：
\[
\text{先识别结构\text{，}再调用高中方法\text{；}先写依据\text{，}再写结论\text{；}先得结果\text{，}再做复核\text{。}}
\]
最后回扣本节主题“高考综合题的结构识别与规范表达”。

\noindent\textbf{教师追问：}今天最重要的两个动作是什么？\\

\textbf{预期回答：}一是判型，二是按规范书写并复核。\par

\noindent\textbf{教师追问：}以后遇到综合题，第一分钟应该做什么？\\

\textbf{预期回答：}先从题干中找结构信号，再决定用定义法、坐标法、韦达、点差还是切线公式。\par

\subsection*{课堂真题}

\paragraph{2017 年高考数学 全国II卷-理科第20题}

设 $O$ 为坐标原点，动点 $M$ 在椭圆 $C:\frac{x^2}{2}+y^2=1$ 上，过 $M$ 作 $x$ 轴的垂线，垂足为 $N$，点 $P$ 满足 $\overrightarrow{NP}=\sqrt{2}\overrightarrow{NM}$．
（1）求点 $P$ 的轨迹方程；
（2）设点 $Q$ 在直线 $x=-3$ 上，且 $\overrightarrow{OP}\cdot\overrightarrow{PQ}=1$．证明：过点 $P$ 且垂直于 $OQ$ 的直线 $l$ 过 $C$ 的左焦点 $F$．

\textbf{答案：}（1）$x^2+y^2=2$；（2）略\par

\textbf{解答：}（1）设 $M(x_0,y_0)$，由题意可得 $N(x_0,0)$，设 $P(x,y)$，由点 $P$ 满足 $\overrightarrow{NP}=\sqrt{2}\overrightarrow{NM}$．可得 $(x-x_0,y)=\sqrt{2}(0,y_0)$，可得 $x-x_0=0$，$y=\sqrt{2}y_0$，即有 $x_0=x$，$y_0=\frac{y}{\sqrt{2}}$，代入椭圆方程 $\frac{x^2}{2}+y^2=1$，可得 $\frac{x^2}{2}+\frac{y^2}{2}=1$，即有点 $P$ 的轨迹方程为圆 $x^2+y^2=2$；
（2）证明：设 $Q(-3,m)$，$P(\sqrt{2}\cos\alpha,\sin\alpha)$（$0\le\alpha<2\pi$），$\overrightarrow{OP}\cdot\overrightarrow{PQ}=1$，可得 $(\sqrt{2}\cos\alpha,\sin\alpha)\cdot(-3-\sqrt{2}\cos\alpha,m-\sin\alpha)=1$，即为 $-3\sqrt{2}\cos\alpha-2\cos^2\alpha+\sqrt{2}m\sin\alpha-2\sin^2\alpha=1$，当 $\alpha=0$ 时，上式不成立，则 $0<\alpha<2\pi$，解得 $m=\frac{3(1+\sqrt{2}\cos\alpha)}{\sqrt{2}\sin\alpha}$，即有 $Q(-3,\frac{3(1+\sqrt{2}\cos\alpha)}{\sqrt{2}\sin\alpha})$，椭圆 $\frac{x^2}{2}+y^2=1$ 的左焦点 $F(-1,0)$，由 $\overrightarrow{PF}\cdot\overrightarrow{OQ}=(-1-\sqrt{2}\cos\alpha,-\sin\alpha)\cdot(-3,\frac{3(1+\sqrt{2}\cos\alpha)}{\sqrt{2}\sin\alpha})=3+3\sqrt{2}\cos\alpha-3(1+\sqrt{2}\cos\alpha)=0$．可得过点 $P$ 且垂直于 $OQ$ 的直线 $l$ 过 $C$ 的左焦点 $F$．

\paragraph{2019 年高考数学 全国II卷-理科第21题}

已知点 $A(-2,0)$，$B(2,0)$，动点 $M(x,y)$ 满足直线 $AM$ 与 $BM$ 的斜率之积为 $-\frac{1}{2}$．记 $M$ 的轨迹为曲线 $C$．
（1）求 $C$ 的方程，并说明 $C$ 是什么曲线；
（2）过坐标原点的直线交 $C$ 于 $P$，$Q$ 两点，点 $P$ 在第一象限，$PE\perp x$ 轴，垂足为 $E$，连结 $QE$ 并延长交 $C$ 于点 $G$．
     （i）证明：$\triangle PQG$ 是直角三角形；
     （ii）求 $\triangle PQG$ 面积的最大值．

\textbf{答案：}（1）$\frac{x^2}{4}+\frac{y^2}{2}=1\ (|x|\neq 2)$，$C$ 为中心在坐标原点，焦点在 $x$ 轴上的椭圆，不含左右顶点；（2）（i）证明见解析；（ii）$\frac{16}{9}$\par

\textbf{解答：}（1）直线 $AM$ 的斜率为 $\frac{y}{x+2}\ (x\neq -2)$，直线 $BM$ 的斜率为 $\frac{y}{x-2}\ (x\neq 2)$，由题意可知 $\frac{y}{x+2}\cdot\frac{y}{x-2}=-\frac{1}{2}$，化简得 $\frac{x^2}{4}+\frac{y^2}{2}=1\ (x\neq\pm 2)$，所以 $C$ 为中心在坐标原点，焦点在 $x$ 轴上的椭圆，不含左右顶点．
（2）（i）设直线 $PQ$ 的斜率为 $k$，则其方程为 $y=kx\ (k>0)$．由 $\begin{cases} y=kx \\ \frac{x^2}{4}+\frac{y^2}{2}=1 \end{cases}$ 得 $x=\pm\frac{2}{\sqrt{1+2k^2}}$．记 $u=\frac{2}{\sqrt{1+2k^2}}$，则 $P(u,uk)$，$Q(-u,-uk)$，$E(u,0)$．于是直线 $QG$ 的斜率为 $\frac{k}{2}$，方程为 $y=\frac{k}{2}(x-u)$．由 $\begin{cases} y=\frac{k}{2}(x-u) \\ \frac{x^2}{4}+\frac{y^2}{2}=1 \end{cases}$ 得 $(2+k^2)x^2-2uk^2x+k^2u^2-8=0$．(1) 设 $G(x_G,y_G)$，则 $-u$ 和 $x_G$ 是方程(1)的解，故 $x_G=\frac{u(3k^2+2)}{2+k^2}$，$y_G=\frac{uk^3}{2+k^2}$．从而直线 $PG$ 的斜率为 $\frac{\frac{uk^3}{2+k^2}-uk}{\frac{u(3k^2+2)}{2+k^2}-u}=-\frac{1}{k}$．所以 $PQ\perp PG$，即 $\triangle PQG$ 是直角三角形．
（ii）由（i）得 $|PQ|=2u\sqrt{1+k^2}$，$|PG|=\frac{2uk\sqrt{1+k^2}}{2+k^2}$，所以 $\triangle PQG$ 的面积 $S=\frac{1}{2}|PQ||PG|=\frac{8k(1+k^2)}{(1+2k^2)(2+k^2)}=\frac{8(\frac{1}{k}+k)}{1+2(\frac{1}{k}+k)^2}$．设 $t=k+\frac{1}{k}$，则由 $k>0$ 得 $t\geq 2$，当且仅当 $k=1$ 时取等号．因为 $S=\frac{8t}{1+2t^2}$ 在 $[2,+\infty)$ 单调递减，所以当 $t=2$，即 $k=1$ 时，$S$ 取得最大值 $\frac{16}{9}$．因此，$\triangle PQG$ 面积的最大值为 $\frac{16}{9}$．

\paragraph{2019 年高考数学 全国III卷-理科第21题}

已知曲线 $C: y=\dfrac{x^2}{2}$，$D$ 为直线 $y=-\dfrac{1}{2}$ 上的动点，过 $D$ 作 $C$ 的两条切线，切点分别为 $A$，$B$。
（1）证明：直线 $AB$ 过定点；
（2）若以 $E\left(0,\dfrac{5}{2}\right)$ 为圆心的圆与直线 $AB$ 相切，且切点为线段 $AB$ 的中点，求四边形 $ADBE$ 的面积。

\textbf{答案：}（1）定点为 $\left(0,\dfrac{1}{2}\right)$；（2）四边形 $ADBE$ 的面积为 $3$ 或 $4\sqrt{2}$。\par

\textbf{解答：}（1）设 $D\left(t,-\dfrac{1}{2}\right)$，$A(x_1,y_1)$，则 $x_1^2=2y_1$。由于 $y'=x$，所以切线 $DA$ 的斜率为 $x_1$，故 $\dfrac{y_1+\frac{1}{2}}{x_1-t}=x_1$，整理得 $2tx_1-2y_1+1=0$。设 $B(x_2,y_2)$，同理可得 $2tx_2-2y_2+1=0$。故直线 $AB$ 的方程为 $2tx-2y+1=0$，所以直线 $AB$ 过定点 $\left(0,\dfrac{1}{2}\right)$。
（2）由（1）得直线 $AB$ 的方程为 $y=tx+\dfrac{1}{2}$。与 $y=\dfrac{x^2}{2}$ 联立得 $x^2-2tx-1=0$，于是 $x_1+x_2=2t$，$x_1x_2=-1$，$y_1+y_2=t(x_1+x_2)+1=2t^2+1$，$|AB|=\sqrt{1+t^2}|x_1-x_2|=\sqrt{1+t^2}\times\sqrt{(x_1+x_2)^2-4x_1x_2}=2(t^2+1)$。设 $d_1,d_2$ 分别为点 $D$，$E$ 到直线 $AB$ 的距离，则 $d_1=\sqrt{t^2+1}$，$d_2=\dfrac{2}{\sqrt{t^2+1}}$。因此四边形 $ADBE$ 的面积 $S=\dfrac{1}{2}|AB|(d_1+d_2)=(t^2+3)\sqrt{t^2+1}$。设 $M$ 为线段 $AB$ 的中点，则 $M\left(t,t^2+\dfrac{1}{2}\right)$。由于 $\overrightarrow{EM}\perp\overrightarrow{AB}$，而 $\overrightarrow{EM}=(t,t^2-2)$，$\overrightarrow{AB}$ 与向量 $(1,t)$ 平行，所以 $t+(t^2-2)t=0$，解得 $t=0$ 或 $t=\pm1$。当 $t=0$ 时，$S=3$；当 $t=\pm1$ 时，$S=4\sqrt{2}$。因此，四边形 $ADBE$ 的面积为 $3$ 或 $4\sqrt{2}$。

\paragraph{2019 年高考数学 全国II卷-文科第20题}

已知 $F_1,F_2$ 是椭圆 $C:\dfrac{x^2}{a^2}+\dfrac{y^2}{b^2}=1(a>b>0)$ 的两个焦点，$P$ 为 $C$ 上一点，$O$ 为坐标原点．
（1）若 $\triangle POF_2$ 为等边三角形，求 $C$ 的离心率；
（2）如果存在点 $P$，使得 $PF_1\perp PF_2$，且 $\triangle F_1PF_2$ 的面积等于 $16$，求 $b$ 的值和 $a$ 的取值范围．

\textbf{答案：}（1）$e=\sqrt{3}-1$；（2）$b=4$，$a$ 的取值范围为 $[4\sqrt{2},+\infty)$\par

\textbf{解答：}（1）连结 $PF_1$，由 $\triangle POF_2$ 为等边三角形可知：在 $\triangle F_1PF_2$ 中，$\angle F_1PF_2=90^\circ$，$PF_2=c$，$PF_1=\sqrt{3}c$，于是 $2a=PF_1+PF_2=c+\sqrt{3}c$，故椭圆 $C$ 的离心率为 $e=\dfrac{c}{a}=\dfrac{1}{1+\sqrt{3}}=\sqrt{3}-1$．
（2）由题意可知，满足条件的点 $P(x,y)$ 存在，当且仅当 $\dfrac{1}{2}|y|\cdot2c=16$，$\dfrac{y}{x+c}\cdot\dfrac{y}{x-c}=-1$，$\dfrac{x^2}{a^2}+\dfrac{y^2}{b^2}=1$，即 $c|y|=16$ (1)，$x^2+y^2=c^2$ (2)，$\dfrac{x^2}{a^2}+\dfrac{y^2}{b^2}=1$ (3)，由(2)(3)以及 $a^2=b^2+c^2$ 得 $y^2=\dfrac{b^4}{c^2}$，又由(1)知 $y^2=\dfrac{16^2}{c^2}$，故 $b=4$；由(2)(3)得 $x^2=\dfrac{a^2}{c^2}(c^2-b^2)$，所以 $c^2\ge b^2$，从而 $a^2=b^2+c^2\ge 2b^2=32$，故 $a\ge 4\sqrt{2}$；当 $b=4$，$a\ge 4\sqrt{2}$ 时，存在满足条件的点 $P$．所以 $b=4$，$a$ 的取值范围为 $[4\sqrt{2},+\infty)$．

\paragraph{2023 年高考数学 北京卷第6题}

已知抛物线 $C: y^2 = 8x$ 的焦点为 $F$ ，点 $M$ 在 $C$ 上．若 $M$ 到直线 $x = -3$ 的距离为 $5$，则 $|MF| =$ （ ）

\textbf{答案：}D\par

\textbf{解答：}因为抛物线 $C: y^2 = 8x$ 的焦点 $F(2,0)$ ，准线方程为 $x = -2$ ，点 $M$ 在 $C$ 上，所以 $M$ 到准线 $x = -2$ 的距离为 $|MF|$，又 $M$ 到直线 $x = -3$ 的距离为 $5$，所以 $|MF| + 1 = 5$，故 $|MF| = 4$.

\paragraph{2022 年高考数学 全国甲卷-理科第10题}

椭圆 $C:\frac{x^2}{a^2}+\frac{y^2}{b^2}=1(a>b>0)$ 的左顶点为A，点P，Q均在C上，且关于y轴对称. 若直线AP, AQ的斜率之积为 $\frac{1}{4}$，则C的离心率为（ ）

\textbf{答案：}A\par

\textbf{解答：}设$P(x_1,y_1)$，则$Q(-x_1,y_1)$，则由$k_{AP}\cdot k_{AQ}=\frac{1}{4}$得：$k_{AP}\cdot k_{AQ}=\frac{y_1}{x_1+a}\cdot\frac{y_1}{-x_1+a}=\frac{y_1^2}{-x_1^2+a^2}=\frac{1}{4}$. 由$\frac{x_1^2}{a^2}+\frac{y_1^2}{b^2}=1$得$y_1^2=\frac{b^2(a^2-x_1^2)}{a^2}$，代入得$\frac{b^2(a^2-x_1^2)}{a^2(-x_1^2+a^2)}=\frac{1}{4}$，即$\frac{b^2}{a^2}=\frac{1}{4}$，所以椭圆C的离心率$e=\frac{c}{a}=\sqrt{1-\frac{b^2}{a^2}}=\sqrt{1-\frac{1}{4}}=\frac{\sqrt{3}}{2}$. 故选A.

\subsection*{课后作业}

\begin{enumerate}

\item 2020 年高考数学 全国I卷-理科第20题

\item 2022 年高考数学 全国乙卷-理科第20题

\item 2021 年高考数学 全国乙卷-理科第21题

\item 2024 年高考数学 天津卷第18题

\item 2023 年高考数学 新高考II卷第21题

\item 2020 年高考数学 全国II卷-理科第19题

\end{enumerate}

\noindent\textbf{作业答案索引：}

\begin{itemize}

\item 2020 年高考数学 全国I卷-理科第20题：(1)$\frac{x^2}{9}+y^2=1$；(2)直线$CD$过定点$\left(\frac{3}{2},0\right)$

\item 2022 年高考数学 全国乙卷-理科第20题：(1) $\frac{y^2}{4}+\frac{x^2}{3}=1$
(2) $(0,-2)$

\item 2021 年高考数学 全国乙卷-理科第21题：（1）$p=2$；（2）$20\sqrt{5}$

\item 2024 年高考数学 天津卷第18题：(1) $\frac{x^2}{12}+\frac{y^2}{9}=1$; (2) 存在 $T(0,t)$ 且 $-3\leq t\leq \frac{3}{2}$, 使得 $\overrightarrow{TP}\cdot\overrightarrow{TQ}\leq 0$ 恒成立.

\item 2023 年高考数学 新高考II卷第21题：（1）$\dfrac{x^{2}}{4}-\dfrac{y^{2}}{16}=1$；（2）证明见解析.

\item 2020 年高考数学 全国II卷-理科第19题：（1）$\dfrac12$；（2）$C_1:\dfrac{x^2}{36}+\dfrac{y^2}{27}=1$，$C_2:y^2=12x$

\end{itemize}

\subsection*{板书设计}

一、主题：$\text{结构识别}\rightarrow\text{方法落地}\rightarrow\text{规范表达}\rightarrow\text{复核}$

二、识别表
\[
\begin{array}{c|c}
\text{题干信号} & \text{优先方法}\\
\hline
|MF_1|+|MF_2|=\text{常数} & \text{第一定义}\\
\bigl||MF_1|-|MF_2|\bigr|=\text{常数} & \text{第一定义}\\
|MF|=d(M,l) & \text{第二定义}\\
k_{PA}k_{PB}=\text{常数} & \text{第三定义/韦达}\\
\text{弦中点已知} & \text{点差法}\\
\text{切点已知} & \text{点式切线}\\
x^2,y^2\text{系数相等} & \text{配方识圆}
\end{array}
\]

三、例题1：`2017-全国II卷-理科::20`
\[
M(x,y),\ N(x,0),\ \overrightarrow{NP}=\sqrt{2}\overrightarrow{NM}
\]
\[
P(x,\sqrt{2}y),\ \frac{x^2}{2}+y^2=1\Rightarrow x^2+y_P^2=2
\]
结论：$\text{轨迹是圆 }x^2+y^2=2$

四、例题2：`2019-全国II卷-理科::21`
\[
\frac{y}{x+2}\cdot\frac{y}{x-2}=-\frac12
\Rightarrow \frac{x^2}{4}+\frac{y^2}{2}=1
\]
结论：$\text{椭圆}$

五、例题3：`2019-全国III卷-理科::21`
抛物线切线参数式：
\[
y=tx-\frac{t^2}{2}
\]
若 $D(u,-\frac12)$ 在切线上，则
\[
t^2-2ut-1=0,
\quad t_1+t_2=2u,\ t_1t_2=-1
\]

六、例题4：`2019-全国II卷-文科::20`
\[
PF_1\perp PF_2\Rightarrow \angle F_1PF_2=90^\circ
\]
\[
S_{\triangle F_1PF_2}=\frac12|PF_1|\,|PF_2|
\]

七、规范模板
\[
\text{设元}\to\text{列式}\to\text{化简}\to\text{结论}\to\text{检验}
\]

八、复核清单
1. 是否说明范围；
2. 是否排除斜率不存在；
3. 是否区分唯一交点与切线；
4. 是否回扣“求方程/证定点/求最值”要求。

\subsection*{易错点}

\begin{itemize}

\item 把“高观点预判”直接当成正式证明。例如预判是第三定义，但解答中没有设点化简或没有写出标准方程。

\item 见到椭圆、双曲线背景就一直在原曲线上纠缠，没有识别中途已转化为圆，如 `2017-全国II卷-理科::20` 第（1）问。

\item 把“有唯一公共点”误认为“相切”。如切线问题中必须核查是否真为切线，不能只凭交点个数判断。

\item 写出轨迹方程后不说明曲线名称，或说明曲线名称但不检查退化与范围。

\item 使用斜率时忽视分母为 $0$ 的情形，导致论证不完整。

\item 面积最值题只会列式，不会把面积进一步化为单变量函数，或不会说明最值是否取得。

\item 证明‘过定点’时只凭观察报出结论，没有把定点代入所求直线方程验证恒成立。

\item 点差法使用时只记结论，不会从两点代入方程后作差推得斜率与中点关系。

\end{itemize}

\subsection*{教师备注}

1. 本课不是按曲线种类分题，而是按“结构信号”组织复习，这样更贴近高考综合题的真实解题过程。
2. 射影视角只用于提醒教师把零散技巧统一到“点线对应、不变量、对偶痕迹”上，不进入学生正式作答层面。
3. 课堂实施时，例题1与例题2必须完整展开；例题3与例题4重在示范书写骨架，不必全算到底，但关键关系式必须板演。
4. 当堂互评时，建议按四项给分：$\text{判型是否准确}$、$\text{入口是否经济}$、$\text{推导是否有依据}$、$\text{结论是否完整}$。
5. 课后作业安排建议：
第一组 `2020-全国I卷-理科::20`、`2022-全国乙卷-理科::20` 训练“定点结构与直线方程组织”；
第二组 `2021-全国乙卷-理科::21`、`2020-全国II卷-理科::19` 训练“切线与焦点的综合转化”；
第三组 `2024-天津卷::18`、`2023-新高考II卷::21` 训练“参数范围、定直线与复核意识”。
6. 课后讲评时，应特别追问学生三个问题：
\[
\text{这题先识别出了什么结构？}
\quad
\text{为什么选这种高中方法？}
\quad
\text{你的答案哪里体现了规范表达？}
\]
7. 若班级基础较弱，可把 `2019-全国III卷-理科::21` 改为只讲第（1）问，把切线参数法练熟；若班级较强，可补充让学生对比 `2017-全国I卷-理科::20` 与 `2020-全国I卷-理科::20`，总结“斜率和为定值 $\Rightarrow$ 弦过定点”与“外点引弦 $\Rightarrow$ 交点连线过定点”的共同组织方式。

\clearpage

\end{document}
